\documentclass[11pt]{article}
\usepackage[ngerman]{babel}
\usepackage[a4paper,margin=0.9in]{geometry}
\usepackage[T1]{fontenc}
\usepackage[utf8]{inputenc}
\usepackage{lmodern}
\usepackage{amsmath,amssymb,mathtools,mathrsfs}
\usepackage{microtype}
\usepackage{fancyhdr}
\usepackage{array,longtable,enumitem,multicol}
\usepackage{tikz}
\setlength{\parindent}{1.5em}
\setlength{\parskip}{0.18em}
\emergencystretch=6em
\pagestyle{fancy}
\fancyhf{}
\cfoot{\thepage}
\newcommand{\Jac}[2]{\left(\frac{#1}{#2}\right)}
\newcommand{\disc}{\operatorname{disc}}
\newcommand{\Norm}{\operatorname{N}}
\providecommand{\dd}{\,d}
\providecommand{\tht}{\vartheta}
\providecommand{\vark}{\varkappa}
\providecommand{\sn}{\operatorname{sn}}
\providecommand{\cn}{\operatorname{cn}}
\providecommand{\dn}{\operatorname{dn}}
\providecommand{\am}{\operatorname{am}}
\providecommand{\ctg}{\operatorname{cotg}}
\providecommand{\cotg}{\operatorname{cotg}}
\providecommand{\eps}{\varepsilon}
\providecommand{\ii}{\mathrm{i}}
\providecommand{\tg}{\operatorname{tg}}
\providecommand{\jac}[2]{\left(\frac{#1}{#2}\right)}
\newcommand{\Om}{\Omega}
\newcommand{\Omm}{\Omega_m}
\newcommand{\Hm}{H_m}
\newcommand{\Lm}{\Omega_m}
\newcommand{\Normone}{\mathrm{N}_1}
\newcommand{\frp}{\mathfrak p}
\newcommand{\frP}{\mathfrak P}
\newcommand{\frG}{\mathfrak G}
\newcommand{\frA}{\mathfrak A}
\newcommand{\frB}{\mathfrak B}
\newcommand{\Gg}{\mathfrak G}
\newcommand{\Aa}{\mathfrak A}
\newcommand{\Bb}{\mathfrak B}
\newcommand{\pp}{\mathfrak p}
\newcommand{\G}{\mathfrak G}
\newcommand{\B}{\mathfrak B}
\newcommand{\A}{\mathfrak A}
\newcommand{\frakA}{\mathfrak a}
\newcommand{\Na}{N_a}
\newcommand{\abs}[1]{\left|#1\right|}
\newcommand{\idxitem}[1]{\par\hangindent=1.4em\hangafter=1 #1}
\providecommand{\Lzero}{\mathfrak L_0}
\providecommand{\Kgrp}{\mathfrak K}
\lhead{Weber Bd. I bis §33}
\rhead{Deutsche Quelle kumulativ}
\begin{document}
\begin{center}
{\large Erster Band.}\\[0.5em]
{\Large Einleitung.}\\[0.4em]
\end{center}

Wir setzen bei unseren Betrachtungen die natürlichen Zahlen $1,2,3,\ldots$ und die Regeln, nach denen mit diesen Zahlen gerechnet wird, als bekannt und gegeben voraus. Die fundamentalen Rechenarten, die sogenannten vier Species, sind die Addition, die Multiplication, zu der als Wiederholung das Potenziren gehört, die Subtraction und die Division. Die beiden ersten heissen die directen Rechenoperationen; sie sind dadurch ausgezeichnet, dass sie im Reiche der natürlichen Zahlen unbegrenzt ausgeführt werden können. Die erste der indirecten oder inversen Operationen, die Subtraction, lässt sich nur dann ausführen, wenn der Minuend grösser ist als der Subtrahend.

Die Aufgabe der Division kann man auf zwei Arten auffassen. Bei der ersten elementaren Auffassung wird gefragt, wie oft der Divisor im Dividenden enthalten ist. Eine Zahl ist nicht in einer kleineren enthalten. Ist aber der Dividend gleich oder grösser als der Divisor, so giebt die Beantwortung der Frage einen Quotienten und in den meisten Fällen einen Rest, der kleiner als der Divisor ist. Wenn kein Rest bleibt, so sagt man, die Division geht auf, oder der Dividend ist durch den Divisor theilbar, oder der Divisor ist ein Factor oder Theiler der Zahl, die den Dividenden bildet.

Diese Aufgabe, die sich schon auf den ersten Stufen der Rechenkunst einstellt, führt zu einer tiefer liegenden Unterscheidung der Zahlen, die das Fundament aller Zahlentheorie ist.

Da ein Factor nie grösser sein kann als die Zahl, deren Factor sie ist, so hat jede Zahl nur eine endliche Anzahl von Factoren. Jede Zahl ist durch $1$ und durch sich selbst theilbar, und eine Zahl, die sonst keinen Theiler hat, heisst eine Primzahl. Die Zahl $1$ selbst pflegt man aus Zweckmässigkeitsgründen nicht als Primzahl zu bezeichnen. Sind zwei Zahlen durch eine dritte theilbar, so ist auch die Summe und die Differenz der beiden ersten durch die dritte theilbar; und ist eine Zahl durch eine zweite, diese durch eine dritte theilbar, so ist auch die erste durch die dritte theilbar.

Zwei Zahlen haben immer den gemeinsamen Theiler $1$. Wenn sie keinen anderen gemeinsamen Theiler haben, wie z. B. die Zahlenpaare $5$ und $7$ oder $21$ und $38$, so heissen die beiden Zahlen relative Primzahlen oder auch theilerfremde Zahlen.

Unter den gemeinsamen Theilern von irgend zwei gegebenen Zahlen wird einer der grösste sein und es ist eine sehr wichtige Aufgabe, diesen grössten gemeinschaftlichen Theiler zu finden. Dazu führt ein Verfahren, das unter dem Namen Algorithmus des grössten gemeinschaftlichen Theilers bekannt ist und sich schon bei Euklid findet.\footnote{Elemente, Buch VII, 11, Bd. II der Heiberg'schen Ausgabe.}

Sind $a,a_1$ die beiden Zahlen, so nehme man die grössere von ihnen, die $a$ sei, als Dividenden, die kleinere $a_1$ als Divisor, und bestimme den Quotienten $q_1$; und wenn die Division nicht aufgeht, den Rest $a_2$, also
\[
  a=q_1a_1+a_2,
\]
so dass $a_2<a_1$ ist. Jeder gemeinschaftliche Theiler von $a$ und $a_1$ ist dann auch gemeinschaftlicher Theiler von $a_1$ und $a_2$ und umgekehrt. Verfährt man mit $a_1,a_2$ ebenso wie mit $a$ und $a_1$ und setzt, wenn die Division nicht aufgeht,
\[
  a_1=q_2a_2+a_3,
\]
so ist wieder $a_3<a_2$, und jeder gemeinschaftliche Theiler von $a_1$ und $a_2$ ist auch gemeinschaftlicher Theiler von $a_2$ und $a_3$ und umgekehrt. Fährt man auf diese Weise fort zu dividiren, so muss, da die Zahlen $a,a_1,a_2,a_3,\ldots$ immer abnehmen, nothwendiger Weise die Division nach einer endlichen Anzahl von Schritten aufgehen, und es muss also zuletzt ein Paar von Gleichungen
\[
  a_{\nu-2}=q_{\nu-1}a_{\nu-1}+a_\nu,
  \qquad
  a_{\nu-1}=q_\nu a_\nu,
\]
auftreten. Dann ist $a_\nu$ ein gemeinschaftlicher Theiler aller vorausgehenden $a$, also auch von $a$ und $a_1$, und jeder gemeinschaftliche Theiler von $a$ und $a_1$ ist Theiler von $a_\nu$. Also ist $a_\nu$ der grösste gemeinschaftliche Theiler von $a$ und $a_1$, und wir sind zugleich zu dem Satze gelangt, dass jeder gemeinschaftliche Theiler zweier Zahlen in ihrem grössten gemeinschaftlichen Theiler aufgehen muss.

Sind $a$ und $a_1$ relative Primzahlen, so ist der letzte Divisor $a_\nu=1$. Multiplicirt man unter dieser Voraussetzung die vorstehenden Gleichungen mit irgend einer Zahl $b$, so folgt, dass $b$ der grösste gemeinschaftliche Theiler von $ab$ und $a_1b$ ist, und dass also jeder gemeinschaftliche Theiler von $ab$ und $a_1$ oder von $a$ und $a_1b$ Theiler von $b$ sein muss. Ist also $a_1$ relativ prim zu $a$ und zu $b$, so ist es auch relativ prim zu $ab$, und aus der speciellen Annahme, dass $a_1$ eine Primzahl sei, ergiebt sich, dass ein Product nur dann durch eine Primzahl theilbar sein kann, wenn wenigstens einer der Factoren durch sie theilbar ist. Ist also das Product $ab$ durch $a_1$ theilbar und ist $a_1$ relativ prim zu $a$, so muss $b$ durch $a_1$ theilbar sein.

Sind $a,b$ irgend zwei Zahlen mit dem grössten gemeinschaftlichen Theiler $d$ und ist $a=da'$, $b=db'$, so sind $a'$ und $b'$ relativ prim zu einander. Jede Zahl $m$, die zugleich ein Vielfaches von $a$ und von $b$ ist, hat die Form $m=am'=da'm'$ und darin muss $a'm'$ ein Vielfaches von $b'$ sein, also muss auch $m'$ ein Vielfaches von $b'$ sein, d. h. jede Zahl, die zugleich ein Vielfaches von $a$ und von $b$ ist, ist durch $a'b'd$ theilbar. Diese Zahl $a'b'd$, die selbst ein gemeinschaftliches Vielfaches von $a$ und $b$ ist, heisst daher das kleinste gemeinschaftliche Vielfache von $a$ und $b$.

Wenn eine Zahl durch zwei relative Primzahlen theilbar ist, so ist sie auch durch ihr Product theilbar, und wenn also eine Zahl $m$ durch mehrere Zahlen theilbar ist, von denen je zwei zu einander relativ prim sind, so ist sie auch durch das Product aller dieser Zahlen theilbar.

Hierauf gründet sich der Beweis des wichtigen Satzes, dass eine Zahl $m$ immer und nur auf eine einzige Weise als ein Product von Primzahlen dargestellt werden kann. Denn da der Theiler nicht grösser sein kann als der Dividend, so kann $m$ sicher nur durch eine endliche Anzahl von Primzahlen theilbar sein. Ist $a$ eine von diesen Primzahlen und $a^\alpha$ die höchste Potenz von $a$, die in $m$ aufgeht, so ist auch $\alpha$ eine bestimmte Zahl. Es seien nun ebenso $b^\beta,c^\gamma,\ldots$ die höchsten Potenzen der übrigen in $m$ aufgehenden Primzahlen $b,c,\ldots$, durch die sich $m$ theilen lässt, dann muss, da je zwei der Zahlen $a^\alpha,b^\beta,c^\gamma,\ldots$ relativ prim sind, $m$ auch durch das Product $a^\alpha b^\beta c^\gamma\cdots$ theilbar sein, und es muss $m$ diesem Producte gleich sein, da sonst noch eine andere Primzahl oder eine höhere Potenz einer der Primzahlen $a,b,c,\ldots$ in $m$ aufgehen müsste.

Wir geben nun einen Ueberblick über die in der Mathematik nothwendigen und allmählich eingeführten Erweiterungen des Zahlenbegriffes.

Wir verstehen unter einer Mannigfaltigkeit oder Menge, oder dem abkürzenden Zeichen $M$, ein System von Objecten oder Elementen irgend welcher Art, das so in sich abgegrenzt und vollendet ist, dass von jedem beliebigen Object vollkommen bestimmt ist, ob es zu dem System gehört oder nicht, gleichviel, ob wir im Stande sind, in jedem besonderen Falle die Entscheidung wirklich zu treffen oder nicht.

Eine Menge heisst geordnet, wenn von irgend zwei unterschiedenen ihrer Elemente immer ein in sich vollkommen bestimmtes als das grössere gilt, und zwar so, dass aus $a>b$, $b>c$ stets $a>c$ folgt. Ist $a>b$ und $b>c$ oder $a<b$ und $b<c$, so sagen wir, dass $b$ zwischen $a$ und $c$ liegt.

Die natürlichen Zahlen bilden eine geordnete Menge; zwischen zwei auf einander folgenden ihrer Elemente liegt kein weiteres Element. Eine solche Mannigfaltigkeit heisst eine discrete. Eine geordnete Menge von der Eigenschaft, dass zwischen je zwei Elementen immer noch andere Elemente gefunden werden, heisst dicht. Eine dichte Menge kann man bilden, wenn man die natürlichen Zahlen in Paaren zusammenfasst, und diese Paare als Elemente einer Menge auffasst. Diese Paare sollen Brüche genannt und mit $m:n$ oder $m/n$ bezeichnet werden, und zwei solche Brüche $m:n$ und $m':n'$ werden einander gleich gesetzt, wenn $mn'=nm'$ ist. Fasst man alle unter einander gleichen Brüche zu einem Element zusammen, so erhält man eine Mannigfaltigkeit, die geordnet ist, wenn man noch festsetzt, dass $m:n$ grösser als $m':n'$ ist, wenn $mn'>nm'$ ist. Dass diese Mannigfaltigkeit dicht ist, sieht man so ein: sind $\mu=m:n$, $\mu'=m':n'$ zwei Brüche und $\mu>\mu'$, so kann man, wenn $h$ eine willkürliche Zahl ist,
\[
 {hmn'\over hnn'},
 \qquad
 {hm'n\over hnn'}
\]
setzen, und darin ist $hmn'>hm'n$. Man kann $h$ immer so annehmen, dass zwischen $hmn'$ und $hm'n$ noch Zahlen liegen, und wenn $p$ eine solche Zahl ist, so liegt $p:hnn'$ zwischen $\mu$ und $\mu'$.

Die Punkte einer geraden Linie kann man auch als eine geordnete Menge auffassen, wenn man unter grösser und kleiner irgend eine Ortsbezeichnung, z. B. weiter rechts und weiter links oder höher und tiefer versteht.

Eine Eintheilung einer geordneten Menge $M$ in zwei Theile $A,B$ der Art, dass jedes Element $a$ von $A$ kleiner ist als jedes Element $b$ von $B$, wird ein Schnitt in $M$ genannt und wird passend durch $(A,B)$ bezeichnet. Ein solcher Schnitt entsteht, wenn man irgend ein Element $\mu$ von $M$ herausgreift, alle kleineren Elemente zu $A$, alle grösseren zu $B$ und $\mu$ selbst nach Belieben zu $A$ oder zu $B$ rechnet. Es entstehen, genau gesagt, je nachdem man das eine oder das andere thut, zwei Schnitte, die wir aber immer als gleich betrachten wollen. Wenn in einem Schnitt $(A,B)$ entweder $A$ ein grösstes oder $B$ ein kleinstes Element $\mu$ enthält, so sagen wir, dass $\mu$ den Schnitt $(A,B)$ erzeugt. Es kann aber auch der Fall vorkommen, dass weder $A$ ein grösstes noch $B$ ein kleinstes Element besitzt.

Wenn jeder Schnitt in einer dichten Menge durch ein bestimmtes Element $\mu$ erzeugt wird, so heisst die Menge stetig.

Stetigkeit sowohl als Dichtigkeit sind Eigenschaften, die der Natur der Sache nach unserer Sinneswahrnehmung unzugänglich sind; sie lassen sich daher auch an Dingen der Aussenwelt, an Raumgrössen, Zeiträumen, Massen niemals mit Strenge nachweisen, wie sehr sie uns auch im Wesen unserer Anschauung zu liegen scheinen. Es lassen sich aber sehr wohl reine Begriffssysteme construiren, denen die Dichtigkeit ohne die Stetigkeit oder auch Dichtigkeit und Stetigkeit zukommen.\footnote{Dies ist von Dedekind nachgewiesen, dem wir überhaupt die oben gegebene Definition der Stetigkeit verdanken. Vgl. die Schriften von Dedekind, ``Stetigkeit und irrationale Zahlen'', Braunschweig 1872, 1892. ``Was sind und was sollen die Zahlen?'', Braunschweig 1888, 1893. Andere Mannigfaltigkeiten, denen die Stetigkeit zukommt, sind von Weierstrass und G. Cantor gebildet.}


Diese sehr abstracte Betrachtungsweise giebt uns die Sicherheit, dass die Annahme einer stetigen Menge keinen Widerspruch enthält, dass solche Mengen wenigstens im Reiche der Gedanken existiren. Die Geometrie wie die Analysis, die immer gern an die geometrische Anschauung anknüpft, hat lange stillschweigend die Existenz stetiger Mengen, z. B. bei den Punkten einer geraden Linie oder irgend eines anderen zusammenhängenden Linienzuges, als eine Art von Axiom angenommen. Auch der Unterschied zwischen dichter und stetiger Menge, der der Unterscheidung commensurabler und incommensurabler Strecken zu Grunde liegt, ist den Alten nicht entgangen.\footnote{Euklid, Elemente, Buch X.}

Auch wir wollen in der Folge nicht auf das Hülfsmittel der geometrischen Anschauung verzichten, und z. B. die Punkte einer geraden Linie unbedenklich als eine stetige Menge betrachten.

Eine geordnete Menge $M$ heisst messbar unter folgenden Voraussetzungen. Addition und Vervielfältigung sind in $M$ allgemein ausführbar, ebenso Subtraction eines kleineren von einem grösseren Element, d. h. aus irgend zwei Elementen $a,b$ (die auch identisch sein können) kann nach einer bestimmten Vorschrift ein neues Element $a+b$ von $M$ abgeleitet werden, so dass $a+b$ grösser als $a$ und als $b$ ist, und dass die bekannten, in den Formeln
\[
 a+b=b+a,
 \qquad (a+b)+c=a+(b+c)
\]
ausgedrückten Regeln der Addition gelten. Und zu zwei Elementen $a,c$, von denen das zweite grösser ist, kann ein drittes Element $b$ gefunden werden, so dass $a+b=c$ ist, was auch durch das Zeichen der Subtraction $b=c-a$ ausgedrückt wird. Zwei ungleiche Elemente von $M$ haben also immer eine bestimmte Differenz. Aus diesen Voraussetzungen folgt, dass eine Summe grösser wird, wenn einer der Summanden sich vergrössert. Die wiederholte, etwa $m$-malige Addition desselben Elementes $a$ heisst Vervielfältigung und ihr Ergebniss wird mit $ma$ bezeichnet.

Es kommt endlich noch eine Voraussetzung hinzu, nämlich die, dass bei jedem gegebenen $a$ ein hinlänglich hohes Vielfaches $ma$ grösser ist als ein beliebig gegebenes anderes Element $b$. Unter den Elementen einer messbaren Menge giebt es also kein grösstes.

In einer dichten messbaren Menge giebt es auch kein kleinstes Element. Denn wäre $a$ das kleinste und $b$ ein beliebiges Element, so könnten zwischen $b$ und $b+a$ keine Elemente liegen; denn wenn $b<c<b+a$ wäre, so würde aus der Definition der Messbarkeit folgen, dass $a'=c-b$ kleiner als $a$ wäre. Es folgt auch umgekehrt, dass eine messbare Menge, in der kein kleinstes Element vorkommt, dicht ist. Denn sind $a$ und $a+b$ irgend zwei Elemente, so braucht man ja zu $a$ nur ein Element zu addiren, was kleiner als $b$ ist, um ein Element zwischen $a$ und $a+b$ zu erhalten.

Ist eine Menge stetig, so lassen sich die Voraussetzungen für die Messbarkeit noch vereinfachen, weil dann die Subtraction eine Folge der Addition ist. Sind nämlich $a$ und $c$ zwei Elemente einer stetigen geordneten Menge $M$, in der die Addition besteht, und ist $c>a$, so erhält man einen Schnitt $(A,B)$ in $M$, wenn man alle Elemente $x$, für die $a+x<c$ ist, nach $A$, und für die $a+x>c$ ist, nach $B$ verweist. Dieser Schnitt wird durch ein Element $b$ erzeugt, für das $a+b=c$ sein muss.

Die natürlichen Zahlen bilden nach unserer Definition eine messbare Menge, in der ein kleinstes Element, nämlich $1$, vorkommt. Die Mannigfaltigkeit der rationalen Brüche wird ebenfalls messbar, wenn man Addition und Subtraction nach den bekannten Regeln der Bruchrechnung erklärt. Besonders wichtig und gleichsam typisch für die messbaren Mengen ist die Mannigfaltigkeit der geradlinigen Strecken oder Längen einer Linie, die einfach durch Aneinanderlegen addirt werden. Auch Stoffmengen, durch die Wage verglichen, und Zeiträume, mit der Uhr gemessen, liefern Beispiele messbarer Mengen. Die Art des Messens liegt nicht in der Natur der Mannigfaltigkeit selbst, sondern wird durch den denkenden Beobachter hineingelegt; so würde es z. B. ebenso gut zulässig sein, unter der Summe $a+b$ zweier Strecken $a$ und $b$ die Hypotenuse eines rechtwinkligen Dreiecks mit den Katheten $a,b$ zu verstehen, statt, wie es gewöhnlich angenommen wird, die aus $a$ und $b$ durch Aneinanderlegen zusammengesetzte Strecke.

Um die stetige Menge $B$ der Schnitte in der Mannigfaltigkeit der rationalen Brüche $R$ zu einer messbaren zu machen, beachte man zunächst Folgendes. Ist $\alpha=(A,B)$ ein Element in $B$ und $\mu$ ein beliebig gegebener rationaler Bruch, so kann man immer ein Element $a'$ in $A$ so bestimmen, dass $a'+\mu=b'$ in $B$ enthalten ist. Denn wählt man zwei beliebige Elemente $a,b$, so kann man die natürliche Zahl $m$ bestimmen, dass $m\mu>b-a$ ist, so dass $a+m\mu$ in $B$ enthalten ist. Ist dann $h$ die kleinste ganze Zahl, für die $a+h\mu$ in $B$ enthalten ist, so ist $a+(h-1)\mu=a'$ in $A$ und also $a'+\mu=b'$ in $B$.

Wir verstehen nun unter der Summe $(A,B)+(A',B')=(A'',B'')$ oder $\alpha+\alpha'=\alpha''$ den Schnitt in $R$, den man erhält, wenn man einen rationalen Bruch $a''$ nur dann nach $A''$ verweist, wenn ein $a$ in $A$ und ein $a'$ in $A'$ existirt von der Beschaffenheit, dass $a''\le a+a'$ ist. In der That ist $(A'',B'')$ ein Schnitt. Denn ist $a''$ in $A''$ enthalten, so gilt dasselbe von jedem kleineren Bruch, und es giebt Brüche, die in $A''$, und Brüche, die nicht in $A''$ enthalten sind, nämlich die Brüche von der Form $a+a'$ und $b+b'$. Es ist ferner $\alpha''$ grösser als $\alpha$ und $\alpha'$. Denn $A''$ enthält zunächst alle $A$, und wenn $a'$ ein beliebiger Bruch in $A'$ ist, so kann man $a$ in $A$ so wählen, dass das Element $a+a'$ von $A'$ in $B$ enthalten ist; also ist $A''$ umfassender als $A$. Die durch rationale Brüche $\mu,\mu'$ erzeugten Schnitte ergeben durch Addition den durch $\mu+\mu'$ erzeugten Schnitt.

Wir gehen nun über zu der Definition der Verhältnisse, die von Alters her als Grundlage der Zahlenlehre betrachtet werden, und folgen dabei zunächst Euklid.\footnote{Elemente, Buch V.}

Wenn man die Elemente einer messbaren Menge $M$ zu Paaren verbindet, und diese Paare an sich als Elemente betrachtet, so entsteht eine neue Mannigfaltigkeit; wir bezeichnen ein solches Paar mit $a:b$, oder auch mit $a/b$, unterscheiden aber, wenn $a$ und $b$ verschiedene Elemente sind, $a:b$ von $b:a$, und nennen $a$ den Zähler und $b$ den Nenner von $a:b$. Diese Paare wollen wir Verhältnisse nennen und wollen diese neue Menge nun ordnen und messbar machen.

Nehmen wir zunächst an, dass zwei ganze Zahlen $m,n$ existiren, so dass $na=mb$ wird, wie es z. B. immer der Fall ist, wenn $M$ das System der natürlichen Zahlen ist, oder wenn $a,b$ zwei commensurable Strecken sind; dann ist, wenn $p,q$ zwei andere ganze Zahlen sind, $qa=pb$ dann und nur dann, wenn $mq=np$ ist. Das Zahlenpaar $p,q$ ist durch diese Forderung vollständig bestimmt, wenn noch die Bedingung hinzukommt, dass $p,q$ relative Primzahlen sein sollen. Dann kann, wenn $h$ eine beliebige ganze Zahl ist, $m=hp$, $n=hq$ sein. In diesem Falle nennen wir das Verhältniss $a:b$ ein rationales und setzen es gleich dem rationalen Bruch $m:n$ oder $p:q$. Diese rationalen Brüche können hiernach als Verhältnisse ganzer Zahlen aufgefasst werden.

Alle unter einander gleichen rationalen Verhältnisse bilden eine rationale Zahl, und die rationalen Zahlen bilden, wie die rationalen Brüche, eine geordnete, dichte und messbare Mannigfaltigkeit. In die Mannigfaltigkeit der rationalen Zahlen ordnen sich die natürlichen Zahlen selbst mit ein, wenn man unter einer natürlichen Zahl $m$ das Verhältniss $m:1$ versteht.

Wir kehren jetzt zu irgend einer messbaren Mannigfaltigkeit $M$ zurück und nehmen aus ihr irgend zwei Elemente $a$ und $b$. Wählt man, was immer möglich ist, zwei natürliche Zahlen $m,n$, so dass $na>mb$, so heisst das Verhältniss $a:b$ grösser als das rationale Verhältniss $m:n$ oder
\[
 {a\over b}>{m\over n},
\]
und wenn $m:n=p:q$, so ist auch $a:b>p:q$. Ebenso folgt, wenn $n'a<m'b$ ist,
\[
 {a\over b}<{m'
\over n'}.
\]
Ist $a:b>m:n$, so kann man eine und folglich auch beliebig viele rationale Zahlen $m_1:n_1$ finden, so dass
\[
 {a\over b}>{m_1\over n_1}>{m\over n}
\]
ist. Um dies zu zeigen, wähle man eine beliebige ganze Zahl $k$ und bestimme die ganze Zahl $h$ so, dass $h(na-mb)>kb$ wird, was immer möglich ist; dann ist
\[
 {a\over b}>{hm+k\over hn}>{m\over n}.
\]
Und ebenso folgt, wenn $n'a<m'b$, $h'(m'b-n'a)>k'a$ ist,
\[
 {a\over b}<{h'm'
\over h'n'+k'}<{m'
\over n'}.
\]

Wenn nun $a:b$ und $\alpha:\beta$ irgend zwei Verhältnisse sind, die wir der Kürze wegen auch mit $\epsilon,\epsilon'$ bezeichnen wollen, deren Elemente derselben oder auch verschiedenen Mannigfaltigkeiten angehören, so sind zwei Fälle möglich: 1) Es liegt kein rationales Verhältniss $\mu$ zwischen $\epsilon$ und $\epsilon'$, oder 2) es liegt ein rationales Verhältniss zwischen $\epsilon$ und $\epsilon'$.

Im Falle 1) heissen die beiden Verhältnisse $\epsilon$ und $\epsilon'$ einander gleich, und man sieht, dass zwei Verhältnisse, die einem dritten gleich sind, auch unter einander gleich sind. Im Falle 2) heissen die Verhältnisse $\epsilon,\epsilon'$ ungleich. Es kann also entweder $\epsilon<\mu<\epsilon'$ oder $\epsilon>\mu'>\epsilon'$ sein, und diese beiden Fälle schliessen sich aus. Die Grössenbeziehung bleibt auch bestehen, wenn $\epsilon$ oder $\epsilon'$ durch ein ihm gleiches Element ersetzt wird. Im ersten Falle heisst $\epsilon$ kleiner als $\epsilon'$, im zweiten grösser als $\epsilon'$. Zwischen zwei ungleichen Verhältnissen kann man eine beliebige Anzahl rationaler Verhältnisse einschieben.

Wenn wir nun alle unter einander gleichen Verhältnisse zusammenfassen, so erhalten wir einen Gattungsbegriff, den wir als Zahl im allgemeinen Sinne des Wortes bezeichnen. Die Zahl ist also ein Name oder Zeichen für eine gewisse Mannigfaltigkeit, deren Elemente eben die mit einem unter ihnen gleichen Verhältnisse sind.\footnote{Auf den Gattungsbegriff lassen sich auch die natürlichen Zahlen in einfacher und folgerichtiger Weise zurückführen.} Unter diesem Zahlbegriff sind die rationalen Verhältnisse und folglich auch die natürlichen Zahlen als die Verhältnisse $m:1$ mit enthalten und bilden die rationalen Zahlen. Zahlen, die nicht aus rationalen Verhältnissen entspringen, heissen irrationale Zahlen. Nach dem, was bis jetzt ausgeführt ist, bilden die Zahlen eine geordnete Menge, und man kann ihre Ordnung feststellen, wenn für jede Zahl irgend eines der darunter enthaltenen Verhältnisse als Repräsentant gewählt wird.

Von zwei Verhältnissen mit demselben Nenner und ungleichen Zählern ist das das grössere, dessen Zähler grösser ist, und von zwei Verhältnissen mit demselben Zähler und ungleichen Nennern ist das das kleinere, dessen Nenner grösser ist. Sind nämlich $a,a',b$ beliebige Elemente einer messbaren Menge und $a'>a$, so wähle man zunächst eine ganze Zahl $n$ so, dass $na>b$ und $n(a'-a)>b$. Hierauf nehme man die kleinste ganze Zahl $m$, die der Bedingung $mb>na$ genügt; dann ist $na<mb$, aber $mb<na'$. Dann ist
\[
 {a\over b}<{m\over n}<{a'
\over b},
\]
also $a:b<a':b$. Ganz ebenso kann man beweisen, dass, wenn $b'>b$ ist, $a:b>a:b'$ wird. Man drückt diesen Satz auch so aus, dass ein Verhältniss zugleich mit dem Zähler wächst und mit wachsendem Nenner abnimmt.

Hieraus ergiebt sich auch leicht der folgende Satz: Sind $a,b,c,d$ Elemente derselben messbaren Menge und ist $a:b=c:d$, so ist auch $a:c=b:d$. Denn angenommen, es wäre $a:c<b:d$, so müsste es zwei ganze Zahlen $m,n$ geben, so dass
\[
 na<mc,
 \qquad nb>md.
\]
Dann aber wäre nach dem eben bewiesenen Satze $na:nb<mc:md$, also auch $a:b<c:d$, entgegen der Voraussetzung.

Hieran schliesst sich nun folgender Hauptsatz. Wenn von den vier Grössen $a,b,c,d$ irgend drei aus einer stetigen messbaren Menge beliebig gegeben sind, so lässt sich das vierte in derselben Menge so bestimmen, dass $a:b=c:d$ ist.

Der Satz ist eine unmittelbare Folge der vorausgesetzten Stetigkeit. Denn wenn man z. B. ein Element $x$ von $M$ in $A$ oder in $B$ aufnimmt, je nachdem $x:b$ kleiner oder grösser als $c:d$ ist, so erhält man einen Schnitt, der durch ein Element $a$ erzeugt wird, das der Bedingung $a:b=c:d$ genügt. Es gilt aber dieser Satz auch in gewissen nicht stetigen Mengen, z. B. für die rationalen Brüche.

Hieraus folgt, dass man als Repräsentanten zweier Zahlen immer zwei Verhältnisse wählen kann, deren Elemente derselben Mannigfaltigkeit angehören, und die denselben beliebig zu wählenden Nenner haben. Die Addition wird dann so erklärt, dass
\[
 {a\over c}+{b\over c}={a+b\over c}
\]
ist. Diese Regel umfasst als speciellen Fall die Addition der rationalen Brüche. Es bleibt zu zeigen, dass, wenn $a:c=a':c'$ und $b:c=b':c'$ ist, auch $(a+b):c=(a'+b'):c'$ sein muss. Der Beweis geschieht dadurch, dass aus $a:c=a':c'$ und $(a+b):c>(a'+b'):c'$ auch $b:c>b':c'$ folgt. Ist nämlich
\[
 {a+b\over c}>{m\over n}>{a'+b'
\over c'},
\]
so ist $n(a+b)>mc$, und daher
\[
 {b\over c}>{mc-na\over nc}.
\]
Andererseits folgt aus $a:c=a':c'$ leicht
\[
 {mc-na\over nc}={mc'-na'\over nc'},
\]
und also $b:c>b':c'$, w. z. b. w.

Hiermit ist also nachgewiesen, dass auch die Zahlen, wie wir sie definirt haben, eine messbare Menge bilden. Sind $a$ und $c$ einer stetigen Mannigfaltigkeit entnommen, so bilden auch bei feststehendem $c$ die Verhältnisse $a:c$ eine stetige Menge, und es folgt also, da es überhaupt stetige Mengen giebt, dass auch die Zahlen eine stetige Menge bilden.

Sind $\alpha,\beta,\gamma,\delta$ jetzt Zahlen, so kann man aus der Proportion $\alpha:\beta=\gamma:\delta$ eine beliebige der vier Zahlen durch die drei anderen gegebenen bestimmen. Setzt man $\delta=1$ und sucht $\alpha$, so erhält man die Multiplication $\alpha=\beta\gamma$, und die Vertauschbarkeit der Factoren ist eine Folge des Satzes, dass $\alpha:\gamma=\beta:\delta$ aus $\alpha:\beta=\gamma:\delta$ folgt. Sucht man $\gamma$, so erhält man die Division; und aus der oben gegebenen Definition der Addition folgt die Grundformel $\alpha(\beta+
\gamma)=\alpha\beta+\alpha\gamma$. Die vier Grundrechnungsarten sind also in dem Gebiete der Zahlen ausführbar mit der einzigen Beschränkung, dass bei der Subtraction der Subtrahend kleiner sein muss als der Minuend.

Die Construction eines Schnittes in der Reihe der Zahlen liefert stets den Beweis für die Existenz einer Zahl, die bestimmten Anforderungen genügt. So erhält man einen Schnitt $(A,B)$, wenn man alle und nur die Zahlen, deren Quadrat kleiner als eine bestimmte Zahl $\alpha$ ist, in $A$ aufnimmt; diesem Schnitt entspricht eine bestimmte Zahl, deren Quadrat gleich $\alpha$ ist und die mit $\sqrt\alpha$ bezeichnet wird, und dadurch wird die Existenz der Quadratwurzeln nachgewiesen.

Auf die Schnitte lassen sich auch die von G. Cantor zur Definition der Irrationalzahlen eingeführten Zahlenreihen zurückführen.\footnote{Cantor, Ueber die Ausdehnung eines Satzes aus der Theorie der trigonometrischen Reihen. Mathematische Annalen, Bd. 5 (1872); vergl. auch Heine, Elemente der Functionenlehre. Journal für Mathematik, Bd. 74 (1872).}

Nach Cantor ist unter einer Zahlenreihe irgend ein unbegrenztes, in bestimmter Weise geordnetes System von Zahlen zu verstehen
\[
 S=x_1,x_2,x_3,x_4,\ldots
\]
von der Beschaffenheit, dass es eine bestimmte Zahl $g$ giebt, unter die keine der Zahlen $S$ heruntersinkt, und dass, wenn $\delta$ eine beliebig gewählte Zahl ist, und $x_n,x_m$ verschiedene Elemente aus $S$ sind, $x_n-x_m$ oder $x_m-x_n$ immer kleiner bleibt als $\delta$, wenn $m$ und $n$ eine hinlänglich grosse Zahl überschritten haben.

Es giebt immer Zahlen, die von den Zahlen einer solchen Reihe nicht überschritten werden; denn hat man $\delta$ gewählt und $n$ passend bestimmt, so überschreitet $x_m$, welchen Werth auch $m$ haben mag, niemals die grösste der Zahlen $x_1,x_2,\ldots,x_n,x_n+
\delta$. Man erhält nun einen Schnitt $(A,B)$, wenn man die Zahlen nach $B$ wirft, die, wenn $n$ einen hinlänglich hohen Werth hat, von keinem $x_n$ mehr überschritten werden, und alle anderen Zahlen (die also von unendlich vielen $x_n$ überschritten werden) nach $A$. Wird dieser Schnitt durch die Zahl $\alpha$ erzeugt, so giebt es, wie klein auch $\varepsilon$ sei, immer unendlich viele Zahlen $x_n$ zwischen $\alpha-\varepsilon$ und $\alpha$, und man kann sagen, dass diese durch $S$ vollkommen bestimmte Zahl $\alpha$ durch die Zahlenreihe $S$ erzeugt wird. Nach Cantor ist die Zahlenreihe $S$ geradezu die Definition der Zahl $\alpha$. Selbstverständlich kann eine und dieselbe Zahl $\alpha$ durch sehr verschiedene Zahlenreihen erzeugt werden. Diese Zahlenreihen sind aber alle als unter einander gleich zu betrachten. Man kann unter Anderem die Zahlen von $S$ alle als rationale Brüche annehmen. Es lässt sich auch umgekehrt leicht nachweisen, dass man zu jeder gegebenen Zahl $\alpha$ immer Zahlenreihen $S$ angeben kann, durch die $\alpha$ erzeugt wird, so dass also die Gesammtheit der Zahlenreihen $S$ gleichfalls eine stetige Menge bildet.

Bei der Erklärung der Grundrechnungsarten hat sich bei der Subtraction eine unbequeme Beschränkung ergeben, von der wir uns frei machen, indem wir die negativen Zahlen einführen.

Es möge $x$ jedes Element des bisher definirten Zahlensystems bedeuten, das wir jetzt als das System der positiven Zahlen bezeichnen wollen. Wir nehmen dies Zahlensystem ein zweites Mal und bezeichnen zum Unterschied in diesem zweiten System, das als das System der negativen Zahlen bezeichnet werden soll, jedes Element mit $-x$. Das zweite System ordnen wir nun dem ersten gerade entgegengesetzt, so dass überall, wo in dem System $x$ ``grösser'' steht, in dem System $-x$ ``kleiner'' gesetzt wird, und umgekehrt. Addition und Subtraction werden in $-x$ ebenso erklärt wie in $x$, so dass $(-x)+(-y)=-(x+y)$, $(-x)-(-y)=-(x-y)$ sein soll.

Wir wollen aber diese beiden Zahlensysteme in der Weise zusammenordnen, dass jedes $-x$ kleiner sein soll als jedes $x$. Wir erhalten so eine geordnete Menge, in der kein grösstes und kein kleinstes Element vorhanden ist. Dieses System ist auch im Allgemeinen stetig, nur der einzige Schnitt $(-x,x)$ wird durch kein Element erzeugt, und hier, wo beide Systeme zusammenstossen, ist also noch eine Verletzung der Stetigkeit vorhanden. Um die Stetigkeit herzustellen, müssen wir also dem Schnitt $(-x,x)$ entsprechend noch eine Zahl Null oder $0$ hinzufügen, die eben durch diesen Schnitt definirt ist. Dann haben wir eine geordnete stetige, beiderseits unbegrenzte Menge, die vollständige Reihe der reellen Zahlen.

In dem so erweiterten Zahlenbereiche erklären wir nun die Addition allgemein, indem wir definitionsweise setzen:
\[
 x+(-x)=0,
 \qquad x+0=x,
\]
\[
 x+(-y)=x-y \quad (x>y),
 \qquad
 x+(-y)=-(y-x) \quad (y>x).
\]
Bei dieser Erklärung der Addition gelten, wenn $z_1,z_2,z_3$ irgend drei Zahlen des ganzen Zahlenbereiches sind, die Gesetze
\[
 z_1+z_2=z_2+z_1,
 \qquad (z_1+z_2)+z_3=z_1+(z_2+z_3),
\]
die man das commutative und das associative Gesetz nennt. Man bildet die Summe aus einer beliebigen Anzahl von Summanden, indem man nach Belieben der Reihe nach je zwei Summanden zu einer Summe vereinigt. Die Subtraction braucht nicht mehr besonders berücksichtigt zu werden, wenn man $z_1-z_2$ durch $z_1+(-z_2)$ erklärt und $-(-z)=z$ setzt.

Man stellt die Zahlenreihe anschaulich durch Punkte dar, indem man von einem festen mit $0$ bezeichneten Punkte einer geraden Linie die positiven Zahlen als Strecken nach der einen, etwa der rechten, die negativen Zahlen nach der anderen, linken Seite aufträgt. Das Bild der Summe zweier Strecken $z_1+z_2$ erhält man, wenn man von dem Punkte $z_1$ aus die Strecke von der Länge $+z_2$ nach der rechten oder nach der linken Seite abträgt, je nachdem $z_2$ positiv oder negativ ist.

Die Multiplication und Division wird in den erweiterten Zahlenbereich durch die Gleichungen
\[
 x(-y)=(-x)y=-xy,
 \qquad (-x)(-y)=xy,
 \qquad 0x=0
\]
erklärt. Die Division als die Umkehrung der Multiplication ist immer möglich, ausser wenn der Divisor Null ist.

Eine fernere Erweiterung des Zahlbegriffes besteht in der Einführung der complexen Grössen. Wir combiniren je zwei Zahlen der gesammten Zahlenreihe zu Paaren $(x,y)$, und betrachten zwei solche Paare $(x,y)$ und $(a,b)$ nur dann als gleich, wenn $x=a$, $y=b$ ist. Diese Zahlenpaare bilden eine Mannigfaltigkeit, deren Elemente zwar nicht geordnet werden, mit denen aber die Rechenoperationen der Addition, Subtraction, Multiplication und Division vorgenommen werden sollen, nach folgenden Regeln:
\[
 (x,y)+(a,b)=(x+a,y+b),
\]
\[
 (x,y)(a,b)=(xa-yb,xb+ya).
\]
Wir setzen ausserdem fest, dass $(x,0)=x$ sei, was diesen Gleichungen nicht widerspricht. Es ist $(x,y)$ nur dann $=0$, wenn $x$ und $y$ beide gleich Null sind. Ferner bezeichnen wir zur Abkürzung $(0,1)$ mit $i$. Dann ergeben obige Gleichungen die Folgerungen:
\[
 (x,0)(0,1)=(0,x)=ix,
 \qquad
 (x,0)+(0,y)=(x,y)=x+yi,
\]
\[
 x+yi+a+bi=(x+a)+(y+b)i,
\]
und die Umkehrung der Addition
\[
 x+yi-(a+bi)=(x-a)+i(y-b),
\]
ferner die Multiplication
\[
 (x+yi)(a+bi)=xa-yb+i(xb+ya),
 \qquad i^2=-1,
\]
\[
 (x+yi)(x-yi)=x^2+y^2,
\]
\[
 {x+yi\over a+bi}={ax+by+i(ay-bx)\over a^2+b^2},
\]
wodurch die Division erklärt ist, ausser wenn $a+bi=0$ ist.

Zahlen von der Form $ix$ heissen rein imaginäre Zahlen und $i$ die imaginäre Einheit. Die Zahlen $a+bi$ heissen imaginär oder complex. Das System der reellen und der rein imaginären Zahlen sind darunter als Specialfälle enthalten. Es sind also in dem Gebiete der complexen Zahlen $x+yi$ die Grundrechnungsarten unbegrenzt auszuführen (mit Ausnahme der Division durch Null), und die Rechnung mit den reellen Zahlen ist ein Specialfall davon.

Man stellt die complexen Zahlen $z=x+yi$ nach Gauss geometrisch durch die Punkte einer Ebene dar, indem man ein rechtwinkliges Coordinatensystem zu Grunde legt und den Punkt mit den Coordinaten $x,y$ als Bild des Zahlwerthes $z$ betrachtet. Die Punkte der $x$-Axe stellen in der oben besprochenen Weise die reellen Zahlen $x$ dar. Die Punkte der $y$-Axe sind die Bilder der rein imaginären Zahlen $yi$. Der Coordinatenanfangspunkt ist das Bild der Zahl $0$. Der Radiusvector vom Nullpunkte nach dem Punkte $z$ hat den Zahlwerth $\rho=\sqrt{x^2+y^2}$, der der absolute Werth oder der Betrag oder, nach älterer Ausdrucksweise, der Modulus der complexen Zahl $z$ genannt wird.

Die einzige Zahl $0$ hat den absoluten Werth $0$. Jede positive Zahl kommt bei unendlich vielen complexen Zahlen als absoluter Werth vor und die Bildpunkte aller Zahlen mit demselben absoluten Werthe liegen auf einem Kreise, dessen Mittelpunkt im Coordinatenanfangspunkte liegt.

Zwei imaginäre Zahlen, die sich nur durch das Vorzeichen von $i$ unterscheiden, also $x+yi$ und $x-yi$, heissen conjugirt imaginär. Ihr Product ist das Quadrat des absoluten Werthes von jeder von ihnen. Wenn von zwei conjugirten imaginären Zahlen die eine gleich Null ist, so ist auch die andere gleich Null, und man kann also in jeder richtigen Zahlengleichung $i$ durch $-i$ ersetzen, ohne dass die Richtigkeit gestört wird.

Wir wollen noch den oft angewandten Satz anführen, dass der absolute Werth einer Summe zweier von Null verschiedener complexer Zahlen niemals grösser ist, als die Summe der absoluten Werthe der Summanden, und nur dann gleich, wenn das Verhältniss (der Quotient) beider Summanden reell und positiv ist. Sei nämlich
\[
 z=x+yi,
 \qquad c=a+bi,
 \qquad Z=z+c=(x+a)+(y+b)i,
\]
\[
 \rho^2=x^2+y^2,
 \qquad r^2=a^2+b^2,
 \qquad R^2=(x+a)^2+(y+b)^2.
\]
Dann ist
\[
 (\rho+r-R)(\rho+r+R)=(\rho+r)^2-R^2=2(\rho r-ax-by),
\]
was sicher positiv ist, wenn $ax+by<\rho r$ ist. Wenn aber $ax+by>\rho r$ ist, so folgt aus
\[
 \rho^2r^2-(ax+by)^2=(ay-bx)^2,
\]
dass $\rho r\ge ax+by$ und nur dann gleich, wenn $ay-bx=0$. Daraus also ergiebt sich, dass $\rho+r\ge R$ ist und nur in dem besonderen Falle $\rho+r=R$, wenn $ay-bx=0$, $ax+by>0$, woraus das Gesagte folgt.

Bei der geometrischen Darstellung ist dieser Satz ein Ausdruck dafür, dass in einem Dreieck eine Seite kleiner ist als die Summe der beiden anderen. Der Satz hat noch die andere Folge, dass der absolute Werth einer Summe nicht kleiner sein kann als die Differenz der absoluten Werthe der Summanden. Denn wenden wir den vorigen Satz auf die Summe $c=Z-z$ an, so folgt $r\le R+\rho$ oder
\[
 R\ge r-\rho.
\]
Die Gleichheit findet hier nur dann statt, wenn der Quotient $z:c$ reell und negativ ist.

Es ist noch ein Wort über das wichtigste Hülfsmittel der Algebra, die Buchstabenrechnung, zu sagen. Die Anwendung dieses Hülfsmittels ist so allgemein, dass man bisweilen das Wort Buchstabenrechnung geradezu synonym mit Algebra gebraucht. Die Regeln, wie mit solchen Buchstabenausdrücken gerechnet wird, setzen wir als bekannt voraus. Gleichungen zwischen Buchstabenausdrücken können von zweierlei Art sein: entweder es sind sogenannte Identitäten, d. h. die zwei einander gleich gesetzten Ausdrücke können durch Anwendung der Rechenregeln so umgeformt werden, dass beide Ausdrücke genau übereinstimmen. Man erhält dann aus solchen Buchstabengleichungen richtige Zahlengleichungen, wenn die Buchstaben durch irgend welche, sei es reelle, sei es complexe Zahlen ersetzt werden, vorausgesetzt, dass dabei nicht die Forderung der Division durch Null auftritt. Die Buchstaben in solchen Gleichungen werden oft auch als Variable bezeichnet, weil man sich vorstellen kann, ohne je zu einem Widerspruch zu gelangen, dass für die Buchstaben nach und nach andere und andere Zahlwerthe gesetzt werden.

Eine andere Art von Gleichungen zwischen Buchstabenausdrücken haben nicht diesen Charakter der Identität. Sie enthalten vielmehr eine Forderung, die an die Zahlen gestellt werden, die man ohne einen Fehler zu begehen für die Buchstaben einsetzen darf. Die Algebra hat die Aufgabe, Zahlwerthe zu ermitteln, die einer solchen Forderung genügen, die Gleichung zu lösen. In diesen Gleichungen werden die Buchstaben auch als ``Unbekannte'' bezeichnet. Es kommen sehr häufig in ein und derselben Gleichung Buchstaben von zwei Arten vor, solche, für die beliebige Zahlwerthe gesetzt werden sollen, und andere, deren Zahlwerth erst ermittelt werden soll.

\clearpage
\clearpage

\section*{Erster Abschnitt. Rationale Functionen.}

\section*{\S~1. Ganze Functionen.}

Nächster Gegenstand der Betrachtung sind ganze rationale oder auch kurz ganze Functionen einer Veränderlichen. Wir verstehen darunter Ausdrücke von folgender Form:
\begin{equation}
 f(x)=a_0x^n+a_1x^{n-1}+a_2x^{n-2}+\cdots+a_{n-1}x+a_n,
\tag{1}
\end{equation}
worin $n$ ein ganzzahliger Exponent, der Grad der Function $f(x)$ ist. Der Grad ist eine natürliche Zahl. Bisweilen ist aber auch nützlich, von ganzen rationalen Functionen $0$ten Grades zu sprechen, worunter eine Constante verstanden wird. $x$ heisst die Veränderliche, $a_0,a_1,a_2\ldots a_{n-1},a_n$ die Coefficienten. Sowohl $x$ als $a_0,a_1,\ldots,a_n$ sind Symbole für unbestimmte Grössen, mit denen nach den Regeln der Buchstabenrechnung verfahren wird, für die auch unter Umständen bestimmte Zahlwerthe gesetzt werden können (vgl. die Einleitung).

Wenn die Function $f(x)$ in der Weise wie in (1) geschrieben ist, so nennen wir sie nach absteigenden Potenzen von $x$ geordnet. Die Summanden können in jeder beliebigen anderen Reihenfolge angeordnet, also z. B. auch nach aufsteigenden Potenzen von $x$ geordnet sein:
\[
 f(x)=a_n+a_{n-1}x+\cdots+a_1x^{n-1}+a_0x^n.
\]
Durch Addition (Subtraction) und Multiplication ganzer rationaler Functionen entstehen wieder ganze rationale Functionen. Die Vorschriften der Buchstabenrechnung geben unmittelbar die Bildungsgesetze dieser neuen Functionen.

Bei der Addition ist der Coefficient irgend einer Potenz $x^\nu$ in der Summe gleich der Summe der Coefficienten von $x^\nu$ in den einzelnen Summanden. Der Grad der entstandenen Function ist gleich dem höchsten der Grade der Summanden und kann sich nur in dem besonderen Falle erniedrigen, wenn der höchste Grad in mehreren Summanden vorkommt und die Summe der Coefficienten der höchsten Potenzen gleich Null ist.

Bei der Multiplication zweier oder mehrerer ganzer rationaler Functionen entsteht eine Function, deren Grad gleich der Summe der Grade der Factoren ist.

Um für das Product das Bildungsgesetz der Coefficienten zu übersehen, setzen wir
\begin{equation}
\begin{aligned}
 A(x)&=a_0x^m+a_1x^{m-1}+a_2x^{m-2}+\cdots,\\
 B(x)&=b_0x^n+b_1x^{n-1}+b_2x^{n-2}+\cdots,
\end{aligned}
\tag{2}
\end{equation}
\begin{equation}
 A(x)B(x)=C(x)=c_0x^{m+n}+c_1x^{m+n-1}+c_2x^{m+n-2}+\cdots
\tag{3}
\end{equation}
und erhalten
\[
 c_0=a_0b_0,
\]
\[
 c_1=a_0b_1+a_1b_0,
\]
\[
 c_2=a_0b_2+a_1b_1+a_2b_0,
\]
\[
 \cdots,
\]
oder in allgemeinen Zeichen, wenn $\nu$ eine der Zahlen $0,1,2,\ldots,m+n$ bedeutet:
\begin{equation}
 c_\nu=a_0b_\nu+a_1b_{\nu-1}+a_2b_{\nu-2}+\cdots+a_{\nu-1}b_1+a_\nu b_0,
\tag{4}
\end{equation}
worin alle $a$, deren Index grösser als $m$, und alle $b$, deren Index grösser als $n$ ist, gleich Null zu setzen sind. Denn $c_\nu$ ist der Coefficient von $x^{m+n-\nu}$, und es entsteht also $c_\nu x^{m+n-\nu}$ durch Multiplication aller Glieder von der Form $a_\mu x^{m-\mu}$ mit allen Gliedern der Form $b_{\nu-\mu}x^{n-\nu+\mu}$ und darauf folgende Addition. Sind $a_0$ und $b_0=1$, so ist auch $c_0=1$.

Hat in allen Factoren eines solchen Products die höchste Potenz von $x$ den Coefficienten $1$, so ist auch im Product der Coefficient der höchsten Potenz $=1$.

Aus diesen Vorschriften für die Rechnung mit ganzen Functionen folgt, dass die Regeln des Rechnens, die sich in Formeln wie $ab=ba$, $(ab)c=a(bc)$, $(a+b)c=ac+bc$ und ähnlichen aussprechen, auch wenn $a,b,c$ ganze Functionen von $x$ sind, gelten, und zwar in dem Sinne, dass ganze Functionen nur dann als gleich gelten, wenn gleich hohe Potenzen gleiche Coefficienten haben.

Unter ganzen Functionen mehrerer Veränderlichen verstehen wir Summen von Producten von ganzen, positiven Potenzen der Veränderlichen $x,y,z\ldots$ mit irgend welchen Coefficienten, die als constant gelten. Man kann sich diese Functionen entstanden denken aus den Functionen einer Veränderlichen $x$, wenn man darin die Coefficienten $a_0,a_1\ldots a_n$ selbst wieder als ganze rationale Functionen von anderen Veränderlichen $y,z\ldots$ auffasst. So entstehen Functionen von $m$ Veränderlichen aus Functionen von $m-1$ Veränderlichen. Alle Glieder einer solchen Function sind von der Form $x^r y^s z^t\cdots$, multiplicirt mit einem Coefficienten, und wenn die Function gehörig zusammengefasst ist, so kommt jede Combination der Exponenten $r,s,t\ldots$ nur einmal vor. Zwei so geordnete ganze Functionen gelten nur dann als einander gleich, wenn sie dieselben Producte $x^r y^s z^t\cdots$ mit denselben Coefficienten enthalten.

\section*{\S~2. Ein Satz von Gauss.}

Wir wollen sogleich eine Anwendung der Multiplicationsregel zweier ganzer rationaler Functionen machen zum Beweise eines Satzes von Gauss, der uns später noch nützlich sein wird, hier aber zur Einführung in die Rechnungsweise und als Beispiel dienen soll.\footnote{Gauss, \emph{Disquisitiones arithmeticae}, Art. 42.}

Wir betrachten hier den Fall, dass die Coefficienten in den Functionen $A(x),B(x)$ ganze Zahlen sind, so dass nach (4) auch die Coefficienten von $C(x)=A(x)B(x)$ ganze Zahlen sind.

Wenn die sämmtlichen ganzzahligen Coefficienten $a_0,a_1,\ldots,a_m$ einer Function $A(x)$ keinen gemeinschaftlichen Theiler haben, so heisst die Function $A(x)$ eine ursprüngliche oder primitive Function und der Satz, den wir beweisen wollen, lautet:

Wenn $A(x)$ und $B(x)$ ursprüngliche Functionen sind, so ist auch ihr Product $C(x)$ eine ursprüngliche Function.

Der Beweis ergiebt sich fast unmittelbar aus dem Anblick der Formel (4).

Wenn nämlich die sämmtlichen Coefficienten $c_0,c_1,c_2\ldots c_{m+n}$, wie sie in (4) angegeben sind, einen gemeinschaftlichen Theiler haben, der grösser als $1$ ist, so muss es auch wenigstens eine Primzahl geben, die in allen diesen Coefficienten aufgeht. Es sei $p$ eine solche Primzahl; diese kann nach der Voraussetzung, dass $A(x),B(x)$ ursprünglich seien, weder in allen $a_0,a_1\ldots a_m$, noch in allen $b_0,b_1\ldots b_n$ aufgehen.

Es möge nun $p$ aufgehen in
\[
 a_0,a_1\ldots a_{r-1},\quad \hbox{aber nicht in }a_r,
\]
in
\[
 b_0,b_1\ldots b_{s-1},\quad \hbox{aber nicht in }b_s,
\]
dann ist nach Voraussetzung $r\le m$ und kann auch gleich Null sein, wenn $p$ schon in $a_0$ nicht aufgeht. Ebenso ist $s\le n$.

Bilden wir nun nach (4) $c_{r+s}$ und ordnen es in folgender Weise
\begin{equation}
\begin{aligned}
 c_{r+s}={}&a_rb_s+a_{r-1}b_{s+1}+a_{r-2}b_{s+2}+\cdots\\
&+a_{r+1}b_{s-1}+a_{r+2}b_{s-2}+\cdots,
\end{aligned}
\tag{1}
\end{equation}
so sieht man unmittelbar, dass $c_{r+s}$ nicht durch $p$ theilbar sein kann, wie doch angenommen war; denn das erste Glied $a_rb_s$ ist durch $p$ nicht theilbar, während alle anderen Glieder, da sie mit einem der Coefficienten $a_0,a_1\ldots a_{r-1},b_0,b_1\ldots b_{s-1}$ multiplicirt sind, durch $p$ theilbar sind. Die Annahme also, dass $C(x)$ nicht ursprünglich sei, während es $A(x)$ und $B(x)$ sind, führt zu einem Widerspruch.

Dieser Satz lässt sich übertragen auf Functionen von mehreren Veränderlichen. Wir nennen eine ganze rationale Function von $m$ Veränderlichen mit ganzzahligen Coefficienten ursprünglich oder primitiv, wenn die Coefficienten keinen gemeinsamen Theiler haben, und wir sprechen den Satz aus, dass das Product von zwei ursprünglichen Functionen wieder eine ursprüngliche Function ist.

Um seine Wahrheit einzusehen, brauchen wir nur in der oben durchgeführten Betrachtung die Coefficienten $a_0,a_1,a_2\ldots$, $b_0,b_1,b_2\ldots$ nicht als ganze Zahlen, sondern als ganze rationale Functionen von $m-1$ Veränderlichen $y,z\ldots$ anzunehmen und eine solche Function durch eine Primzahl $p$ theilbar zu nennen, wenn alle ihre Coefficienten durch $p$ theilbar sind. Setzen wir dann voraus, der zu beweisende Satz sei bereits für Functionen von $m-1$ Variablen bewiesen, dann ergiebt die Formel (1) seine Richtigkeit für Functionen von $m$ Variablen, also seine allgemeine Gültigkeit durch den Schluss der vollständigen Induction oder den Schluss von $m-1$ auf $m$.

Eine imprimitive Function ist eine solche, deren ganzzahlige Coefficienten alle einen gemeinsamen Theiler haben. Diesen Theiler nennen wir den Theiler der Function. Dann können wir den bewiesenen Satz auch so aussprechen:

Der Theiler eines Productes zweier ganzer Functionen ist gleich dem Product der Theiler beider Functionen.

Denn sind $PA$ und $QB$ zwei ganze Functionen mit den Theilern $P$ und $Q$, so sind $A$ und $B$ ursprüngliche Functionen. Also ist auch $AB=C$ eine ursprüngliche Function und $PQ$ ist der Theiler der Function $PA\cdot QB=PQC$.

Wir können dem Satze, insofern er sich auf Functionen einer Veränderlichen bezieht, ohne seinen Inhalt wesentlich zu ändern, folgende Fassung geben, in der er besonders nützlich ist. Sind
\[
 \varphi(x)=x^m+\alpha_1x^{m-1}+\alpha_2x^{m-2}+\cdots+\alpha_m,
\]
\[
 \psi(x)=x^n+\beta_1x^{n-1}+\beta_2x^{n-2}+\cdots+\beta_n
\]
zwei ganze rationale Functionen, in denen die höchsten Potenzen von $x$ den Coefficienten $1$ haben, während die übrigen Coefficienten rationale Zahlen sind, so können in dem Product
\[
 \varphi(x)\psi(x)=x^{m+n}+\gamma_1x^{m+n-1}+\gamma_2x^{m+n-2}+\cdots+\gamma_{m+n}
\]
die Coefficienten $\gamma$ nicht alle ganze Zahlen sein, wenn die Coefficienten $\alpha,\beta$ in $\varphi(x)$ und $\psi(x)$ nicht alle ganze Zahlen sind.

Denn bezeichnen wir den kleinsten Hauptnenner der Coefficienten $\alpha$ von $\varphi$ mit $a_0$, der Coefficienten $\beta$ von $\psi$ mit $b_0$, so sind $a_0\varphi(x)=A(x)$, $b_0\psi(x)=B(x)$ primitive Functionen von $x$. Ihr Product $a_0b_0\varphi(x)\psi(x)$ hätte, wenn die Coefficienten $\gamma$ ganze Zahlen wären, den Theiler $a_0b_0$ und wäre also, wenn $a_0$ und $b_0$ nicht beide gleich $1$ wären, nicht primitiv; dies aber wäre ein Widerspruch mit dem oben bewiesenen Satze.

\section*{\S~3. Division.}

Es seien, wie bisher,
\begin{equation}
\begin{aligned}
 A&=A(x)=a_0x^m+a_1x^{m-1}+\cdots,\\
 B&=B(x)=b_0x^n+b_1x^{n-1}+\cdots
\end{aligned}
\tag{1}
\end{equation}
zwei ganze rationale Functionen von $x$; es soll aber jetzt vorausgesetzt werden, dass $m\ge n$ sei und dass $a_0$ und $b_0$ von Null verschieden sind. Dann ist die Differenz
\begin{equation}
 A-{a_0\over b_0}x^{m-n}B
\tag{2}
\end{equation}
auch eine ganze rationale Function von $x$, deren Grad aber kleiner ist als $m$, da die höchste Potenz in beiden Gliedern der Differenz denselben Coefficienten hat und also herausfällt. Wir setzen diese Differenz:
\begin{equation}
 A'=A'(x)=a'_0x^{m'}+a'_1x^{m'-1}+\cdots,
 \qquad m'<m.
\tag{3}
\end{equation}
Ist nun $m'$ noch $\ge n$, so können wir in (2) $A'$ an Stelle von $A$ setzen und dieselbe Schlussweise wiederholen.

So ergiebt sich eine Kette von Gleichungen:
\begin{equation}
\begin{aligned}
 A-{a_0\over b_0}x^{m-n}B&=A',\\
 A'-{a'_0\over b_0}x^{m'-n}B&=A'',\\
 A''-{a''_0\over b_0}x^{m''-n}B&=A''',\\
 &\cdots,
\end{aligned}
\tag{4}
\end{equation}
und diese Kette lässt sich so lange fortsetzen, bis der Grad der entstandenen Function kleiner als $n$ geworden ist. Da nun in der Reihe der Functionen $A,A',A''\ldots$ der Grad jeder folgenden mindestens um eine Einheit erniedrigt ist, so besteht die Kette der Gleichungen (4) höchstens aus $m-n+1$ Gliedern; sie kann aber auch weniger Glieder enthalten, wenn sich gleichzeitig mehrere Potenzen herausheben. Addiren wir die sämmtlichen Gleichungen (4), bezeichnen die letzte der Functionen $A,A'\ldots$ mit $C$ und setzen zur Abkürzung
\begin{equation}
 Q={a_0\over b_0}x^{m-n}+{a'_0\over b_0}x^{m'-n}+\cdots,
\tag{5}
\end{equation}
so dass auch $Q$ eine ganze rationale Function von $x$ ist, so folgt
\begin{equation}
 A=QB+C.
\tag{6}
\end{equation}
Die hier geschilderte Operation, durch die aus $A,B$ die Functionen $Q,C$ gefunden werden, heisst Division. $A$ ist der Dividendus, $B$ der Divisor, $C$ der Rest und $Q$ der Quotient. Der Grad des Restes ist immer niedriger als der Grad des Divisors.

Die Coefficienten der Functionen $Q$ und $C$ sind aus den Coefficienten $a$ und $b$ durch Addition, Subtraction, Multiplication und Theilung zusammengesetzt. Im Nenner kommen aber nur Potenzen von $b_0$ vor, und wenn also $b_0=1$ ist, so sind die Coefficienten von $Q$ und $C$ ganze Functionen der $a$ und $b$. Die höchste Potenz von $b_0$, die im Nenner auftreten kann, ist die $(m-n+1)$te, da in der Kette (4) in jeder folgenden Gleichung im Nenner einmal der Factor $b_0$ hinzukommt. Es kann aber in besonderen Fällen die höchste Potenz von $b_0$ in allen Nennern eine niedrigere sein.

Nehmen wir z. B. für $A$ eine Function dritten Grades
\begin{equation}
 f(x)=a_0x^3+a_1x^2+a_2x+a_3
\tag{7}
\end{equation}
und für $B$ die sogenannte erste Derivirte von $f(x)$, die vom zweiten Grade ist,
\begin{equation}
 f'(x)=3a_0x^2+2a_1x+a_2,
\tag{8}
\end{equation}
so erhält man:
\begin{equation}
 Q={1\over 3}x+{a_1\over 9a_0},
\tag{9}
\end{equation}
\begin{equation}
 C={6a_0a_2-2a_1^2\over 9a_0}x+{9a_0a_3-a_1a_2\over 9a_0}.
\tag{10}
\end{equation}

Wie man die in den Gleichungen (4) vorgeschriebene Rechnung zweckmässig anordnet, darf hier aus den Elementen als bekannt vorausgesetzt werden. Wir machen auf die Analogie aufmerksam, die zwischen dieser Rechnung und der Division im dekadischen Zahlsystem besteht. Eine Function $f(x)$ stellt eine dekadisch geschriebene Zahl dar, wenn die Coefficienten $a_0,a_1\ldots$ ganze Zahlen zwischen Null (einschliesslich) und $10$ (ausschliesslich) sind, und $x=10$ gesetzt wird. Lässt man in den Coefficienten auch Zahlen zu, die grösser als $10$ sind, so kann man eine Zahl auf verschiedene Arten durch $f(x)$ darstellen. Man wendet dies bei dem Divisionsverfahren an, um auf die einfachste Weise in den Resultaten gebrochene und negative Coefficienten zu vermeiden.

\section*{\S~4. Theilung durch eine lineare Function.}

Wir wollen die allgemeine Vorschrift für die Division noch auf einen besonderen Fall anwenden. Es sei der Dividend
\begin{equation}
 f(x)=a_0x^n+a_1x^{n-1}+a_2x^{n-2}+\cdots+a_n
\tag{1}
\end{equation}
beliebig, dagegen der Divisor vom ersten Grade oder, wie man auch sagt, linear. Wir wollen auch den Coefficienten der ersten Potenz von $x$ im Divisor $=1$ voraussetzen und also den Divisor in der einfachen Form $(x-\alpha)$ annehmen, worin $\alpha$ beliebig bleibt. Der Quotient $Q$ ist in diesem Falle vom $(n-1)$ten Grade und der Rest $C$ vom $0$ten Grade, d. h. von $x$ unabhängig. Setzen wir also
\begin{equation}
 f(x)=(x-\alpha)Q+C,
\tag{2}
\end{equation}
so enthält $C$ das $x$ nicht mehr, und wenn wir
\begin{equation}
 Q=q_0x^{n-1}+q_1x^{n-2}+\cdots+q_{n-2}x+q_{n-1}
\tag{3}
\end{equation}
setzen, so folgt aus (2):
\begin{equation}
\begin{aligned}
 f(x)={}&q_0x^n+q_1x^{n-1}+\cdots+q_{n-2}x^2+q_{n-1}x+C\\
&-\alpha q_0x^{n-1}-\cdots-\alpha q_{n-3}x^2-\alpha q_{n-2}x-\alpha q_{n-1},
\end{aligned}
\tag{4}
\end{equation}
und aus der Vergleichung mit (1):
\begin{equation}
\begin{aligned}
q_0&=a_0,\\
q_1-\alpha q_0&=a_1,\\
q_2-\alpha q_1&=a_2,\\
&\cdots,\\
q_{n-1}-\alpha q_{n-2}&=a_{n-1},\\
C-\alpha q_{n-1}&=a_n.
\end{aligned}
\tag{5}
\end{equation}
Daraus erhält man
\begin{equation}
\begin{aligned}
q_0&=a_0,\\
q_1&=a_0\alpha+a_1,\\
q_2&=a_0\alpha^2+a_1\alpha+a_2,\\
&\cdots,\\
q_{n-1}&=a_0\alpha^{n-1}+a_1\alpha^{n-2}+\cdots+a_{n-1},\\
C&=a_0\alpha^n+a_1\alpha^{n-1}+\cdots+a_{n-1}\alpha+a_n=f(\alpha).
\end{aligned}
\tag{6}
\end{equation}
$C$ entsteht aus $f(x)$, wenn man $x=\alpha$ setzt, und kann also auch mit $f(\alpha)$ bezeichnet werden.

Demnach haben wir auch die Formel
\begin{equation}
 {f(x)-f(\alpha)\over x-\alpha}=Q(x),
\tag{7}
\end{equation}
worin $Q(x)$ eine ganze Function vom Grade $n-1$ ist.

Wenn wir in den Ausdrücken (6) an Stelle der unbestimmten Grösse $\alpha$ das Zeichen $x$ setzen, so entsteht daraus eine Reihe von ganzen rationalen Functionen von $x$, die wir, wenn wir der Einfachheit halber $a_0=1$ setzen, so schreiben:
\begin{equation}
\begin{aligned}
f_0&=1,\\
f_1&=x+a_1,\\
f_2&=x^2+a_1x+a_2,\\
&\cdots,\\
f_{n-1}&=x^{n-1}+a_1x^{n-2}+a_2x^{n-3}+\cdots+a_{n-1}.
\end{aligned}
\tag{8}
\end{equation}
Diese Functionen $f_0,f_1\ldots f_{n-1}$ werden uns später noch gute Dienste leisten. Für jetzt fügen wir noch folgende Bemerkungen bei.

Man kann nach (8) die Potenzen $1,x,x^2\ldots x^{n-1}$ von $x$ linear ausdrücken durch die Functionen $f_0,f_1\ldots f_{n-1}$ und zwar so, dass in den Coefficienten nur ganze rationale Verbindungen der $a$ vorkommen, z. B.
\[
 1=f_0,
\]
\[
 x=f_1-a_1f_0,
\]
\[
 x^2=f_2-a_1f_1+(a_1^2-a_2)f_0,
\]
\[
 \cdots,
\]
und daraus folgt, dass man jede ganze rationale Function von $x$, deren Grad nicht grösser als $n-1$ ist, gleichfalls linear durch $f_0,f_1,\ldots,f_{n-1}$ ausdrücken kann in der Form
\begin{equation}
 y_0f_0+y_1f_1+\cdots+y_{n-1}f_{n-1},
\tag{9}
\end{equation}
worin die Coefficienten $y_0,y_1\ldots y_{n-1}$ von $x$ unabhängig sind.

Ist also $F(x)$ eine beliebige ganze rationale Function von $x$, so kann man nach \S~3, indem man $f(x)$ als Divisor betrachtet,
\begin{equation}
 F(x)=Qf(x)+y_0f_0+y_1f_1+\cdots+y_{n-1}f_{n-1}
\tag{10}
\end{equation}
setzen, worin auch $Q$ eine ganze rationale Function von $x$ ist.

Zur recurrenten Berechnung der Functionen $f_r(x)$ ergiebt sich aus (8) die Relation:
\begin{equation}
 f_r(x)-xf_{r-1}(x)=a_r.
\tag{11}
\end{equation}

\clearpage
\clearpage

\section*{\S~5. Gebrochene Functionen; Theilbarkeit.}

Wenn $F(x)$ und $f(x)$ zwei ganze rationale Functionen von $x$ sind, so heisst der Bruch
\begin{equation}
 {F(x)\over f(x)}
\tag{1}
\end{equation}
eine \emph{gebrochene rationale} oder auch kurz \emph{gebrochene} oder \emph{rationale Function} von $x$.

Ist der Grad des Zählers niedriger als der Grad des Nenners, so heisst die Function \emph{echt gebrochen}, im entgegengesetzten Falle \emph{unecht gebrochen}.

Nach \S~3 lassen sich die ganzen rationalen Functionen $Q$ und $\varphi(x)$ so bestimmen, dass
\begin{equation}
 F(x)=Qf(x)+\varphi(x)
\tag{2}
\end{equation}
und der Grad von $\varphi(x)$ kleiner als der Grad von $f(x)$ ist; demnach ist
\begin{equation}
 {F(x)\over f(x)}=Q+{\varphi(x)\over f(x)}
\tag{3}
\end{equation}
und daraus der Satz:

Jede gebrochene Function kann in die Summe aus einer ganzen und einer echt gebrochenen Function zerlegt werden.

Ist $n$ der Grad von $f(x)$, so ist der Grad von $\varphi(x)$ höchstens $n-1$, er kann aber auch niedriger sein; insbesondere kann auch der Fall eintreten, dass $\varphi(x)$ identisch verschwindet.

In diesem Falle heisst die Function $F(x)$ durch $f(x)$ theilbar. Die Function
\[
 {F(x)\over f(x)}
\]
ist in diesem Falle nur scheinbar gebrochen, in Wirklichkeit der ganzen Function $Q$ gleich.

Die Formel
\[
 {x^m-1\over x-1}=x^{m-1}+x^{m-2}+x^{m-3}+\cdots+1
\]
giebt hierfür ein einfaches Beispiel.

Für die Theilbarkeit von Functionen gelten dieselben Gesetze, wie für die Theilbarkeit der Zahlen, insbesondere die folgenden:

1. Wenn die Function $F(x)$ durch die Function $f(x)$, $f(x)$ durch eine dritte Function $\varphi(x)$ theilbar ist, so ist auch $F(x)$ durch $\varphi(x)$ theilbar.

Denn ist
\[
 F=Qf,
 \qquad
 f=q\varphi,
\]
worin $Q$ und $q$ ganze rationale Functionen sind, so ist
\[
 F=Qq\varphi,
\]
und da $Qq$ eine ganze rationale Function ist, $F$ durch $\varphi$ theilbar.

2. Ist $F(x)$ durch $f(x)$ theilbar und $Q$ eine beliebige ganze rationale Function, so ist auch $QF(x)$ durch $f(x)$ theilbar.

3. Ist $F(x)$ und $f(x)$ durch $\varphi(x)$ theilbar, so ist auch $F(x)\pm f(x)$ durch $\varphi(x)$ theilbar, oder allgemeiner:

4. Sind $F_1,F_2\ldots$ durch $f(x)$ theilbar und $Q_1,Q_2\ldots$ beliebige ganze rationale Functionen, so ist auch
\[
 Q_1F_1+Q_2F_2+\cdots
\]
durch $f(x)$ theilbar.

Der letzte Satz umfasst die beiden vorhergehenden und wird einfach so bewiesen.

Sind $F_1,F_2\ldots$ durch $f$ theilbar, so kann man die ganzen rationalen Functionen $\Phi_1,\Phi_2\ldots$ so bestimmen, dass
\[
 F_1=\Phi_1f,\qquad F_2=\Phi_2f,\ldots
\]
und folglich
\[
 Q_1F_1+Q_2F_2+\cdots=(Q_1\Phi_1+Q_2\Phi_2+\cdots)f.
\]
Da nun $Q_1\Phi_1+Q_2\Phi_2+\cdots$ eine ganze rationale Function ist, so ist der Satz bewiesen.

5. Jede Function ist durch sich selbst theilbar.

In Bezug auf die Theilbarkeit oder Untheilbarkeit von Functionen wird nichts geändert, wenn die Functionen mit beliebigen, von $x$ unabhängigen Factoren multiplicirt werden.

Eine von $x$ unabhängige, von Null verschiedene Grösse kann als Function nullten Grades aufgefasst werden. Nennen wir eine solche Grösse eine Constante, so können wir sagen:

6. Jede Function ist durch jede Constante theilbar.

Wenn eine Function durch eine andere theilbar ist, so ist der Grad des Quotienten gleich der Differenz des Grades des Dividenden und des Grades des Divisors. Ersterer kann also nicht kleiner sein als letzterer. Sind die Grade gleich, so ist der Quotient eine Constante, und daraus folgt:

7. Wenn von zwei ganzen rationalen Functionen gleichen Grades die eine durch die andere theilbar ist, so unterscheiden sie sich nur durch einen constanten Factor von einander, und es ist auch die zweite durch die erste theilbar.

8. Nach \S~4 ist die nothwendige und hinreichende Bedingung dafür, dass die Function $f(x)$ durch die lineare Function $x-\alpha$ theilbar ist, die dass $f(\alpha)=0$ sei.

\section*{\S~6. Grösster gemeinschaftlicher Theiler.}

Es ist eine Aufgabe von fundamentaler Bedeutung, zu entscheiden, ob zwei ganze rationale Functionen ausser den Constanten noch einen anderen gemeinschaftlichen Theiler haben. Man findet die Lösung durch den Algorithmus des grössten gemeinschaftlichen Theilers ganz in derselben Weise, wie die entsprechende Frage in Bezug auf die Theilbarkeit ganzer Zahlen beantwortet wird. (S. die Einleitung.)

Es seien
\begin{equation}
 f(x)=A,\qquad \varphi(x)=A'
\tag{1}
\end{equation}
zwei gegebene Functionen; der Grad von $\varphi(x)$ möge niedriger oder wenigstens nicht höher als der von $f(x)$ sein.

Wir dividiren $A$ durch $A'$ und bezeichnen den Rest, dessen Grad niedriger ist als der Grad von $A'$, mit $A''$, also:
\[
 A=Q'A'+A'';
\]
nun dividiren wir $A'$ durch $A''$ und bezeichnen den Rest mit $A'''$, und fahren so fort in der Bildung der Functionenreihe
\begin{equation}
 A,A',A'',A''',\ldots,
\tag{2}
\end{equation}
deren Grade $n,n',n'',n'''\ldots$ immer abnehmen und folglich nach einer endlichen Anzahl von Divisionen auf Null heruntergehen. Der letzte Rest vom Grade Null, der also eine Constante ist, sei $A^{(\nu)}$. Dann haben wir die Kette der Gleichungen:
\begin{equation}
\begin{aligned}
A&=Q'A'+A'',\\
A'&=Q''A''+A''',\\
&\cdots\\
A^{(\nu-3)}&=Q^{(\nu-2)}A^{(\nu-2)}+A^{(\nu-1)},\\
A^{(\nu-2)}&=Q^{(\nu-1)}A^{(\nu-1)}+A^{(\nu)},
\end{aligned}
\tag{3}
\end{equation}
worin $A^{(\nu)}$ eine Constante ist.

Wenn nun $A,A'$ irgend einen gemeinsamen Theiler haben, so ist dieser nach den Sätzen des vorigen Paragraphen, wie der Anblick der Gleichungen (3) lehrt, auch Theiler von $A'',A''',A''''$ u. s. f. bis $A^{(\nu)}$. Ist $A^{(\nu)}$ eine von Null verschiedene Constante, so kann also kein von $x$ abhängiger Theiler von $f(x)$ und $\varphi(x)$ existiren. Solche Functionen heissen \emph{theilerfremd} oder \emph{relativ prim}. Man sagt auch, indem man nur die von $x$ abhängigen Theiler berücksichtigt, die Functionen haben keinen gemeinsamen Theiler.

Die Bedingung, dass die beiden Functionen einen gemeinsamen Theiler haben, ist also:
\begin{equation}
 A^{(\nu)}=0.
\tag{4}
\end{equation}
In diesem Falle ist $A^{(\nu-1)}$ Theiler von $A^{(\nu-2)}$, wie die letzte Gleichung (3) zeigt, und nach der vorletzten dieser Gleichungen Theiler von $A^{(\nu-3)}$ u. s. f., also auch Theiler von $A$ und $A'$, d. h. von $f$ und $\varphi$. Und da umgekehrt jeder gemeinsame Theiler von $A,A'$ auch Theiler von $A^{(\nu-1)}$ ist, so heisst $A^{(\nu-1)}$ der grösste gemeinsame Theiler von $f$ und $\varphi$. Der Algorithmus (3) zeigt, dass man $A^{(\nu)}$ oder $A^{(\nu-1)}$ aus den Coefficienten von $f$ und $\varphi$ durch die rationalen Rechenoperationen ableiten kann, und zwar so, dass immer nur durch die Coefficienten der höchsten Potenzen von $x$ in den Functionen $A,A',A'',\ldots$, die von Null verschieden sind, dividirt wird.

Wir wollen diese Betrachtungen auf zwei Beispiele anwenden.

Es seien zunächst:
\begin{equation}
\begin{aligned}
 A&=a_0x^2+a_1x+a_2,\\
 B&=b_0x^2+b_1x+b_2
\end{aligned}
\tag{5}
\end{equation}
zwei Functionen zweiten Grades, also $a_0,b_0$ von Null verschieden.

Der erste Schritt ist die Bildung der Gleichung
\begin{equation}
 A={a_0\over b_0}B+C,
\tag{6}
\end{equation}
worin
\begin{equation}
 C=c_0x+c_1
\tag{7}
\end{equation}
und
\begin{equation}
 c_0={a_1b_0-b_1a_0\over b_0},\qquad
 c_1={a_2b_0-b_2a_0\over b_0}.
\tag{8}
\end{equation}
Ist $c_0=0$, also $C$ constant, so haben $A,B$ nur dann einen gemeinsamen Factor, wenn auch $c_1=0$ ist, und dann ist $A$ durch $B$ theilbar, d. h. $A$ und $B$ unterscheiden sich nur durch einen constanten Factor. Ist aber $c_0$ von Null verschieden, so setzen wir die Rechnung fort, indem wir
\[
 B=QC+D
\]
setzen, worin (nach \S~4, wenn dort $\alpha=-c_1:c_0$ gesetzt wird)
\begin{equation}
 D={b_0c_1^2-b_1c_0c_1+b_2c_0^2\over c_0^2}.
\tag{9}
\end{equation}
Ist $D$ von Null verschieden, so sind $A,B$ ohne gemeinschaftlichen Theiler. Ist aber $D=0$, so ist $C$ der grösste gemeinschaftliche Factor von $A$ und $B$.

Setzt man die Werthe $c_0,c_1$ aus (8) in (9) ein, so erhält man mit Weglassung des Nenners $b_0c_0^2$ die Bedingung für die Existenz eines gemeinsamen Theilers $C$ von $A$ und $B$ in der Form
\begin{equation}
\begin{aligned}
&a_0^2b_2^2+a_2^2b_0^2-2a_0a_2b_0b_2-a_1a_2b_0b_1-a_0a_1b_1b_2\\
&\qquad{}+a_0a_2b_1^2+a_1^2b_0b_2=0
\end{aligned}
\tag{10}
\end{equation}
oder
\begin{equation}
 (a_0b_2-b_0a_2)^2+(a_0b_1-a_1b_0)(a_2b_1-a_1b_2)=0.
\tag{11}
\end{equation}
Diese Bedingung ist auch erfüllt, wenn $c_0$ und $c_1$ gleich Null sind, und ist also die nothwendige und hinreichende Bedingung dafür, dass $A,B$ einen gemeinsamen Theiler haben.

Die linke Seite von (10) oder (11) heisst die \emph{Resultante} der Functionen $A$ und $B$.

Als zweites Beispiel nehmen wir das schon im \S~3 gewählte. Setzen wir
\begin{equation}
\begin{aligned}
 f(x)&=a_0x^3+a_1x^2+a_2x+a_3=A,\\
 f'(x)&=3a_0x^2+2a_1x+a_2=B,
\end{aligned}
\tag{12}
\end{equation}
und setzen $a_0$ von Null verschieden voraus, so haben wir nach \S~3:
\begin{equation}
 A=QB+C,
\tag{13}
\end{equation}
\begin{equation}
 C=c_0x+c_1,
\tag{14}
\end{equation}
worin
\begin{equation}
 c_0={6a_0a_2-2a_1^2\over 9a_0},\qquad
 c_1={9a_0a_3-a_1a_2\over 9a_0}.
\tag{15}
\end{equation}
Wenn $c_0$ gleich Null ist, so ist hiermit der Algorithmus schon geschlossen; wenn $c_1$ von Null verschieden ist, dann haben $A$ und $B$ keinen gemeinsamen Theiler; ist aber $c_1$ auch gleich Null, so ist $B$ selbst der grösste gemeinsame Theiler von $A$ und $B$, d. h. $A$ ist durch $B$ theilbar. Die Bedingungen hierfür sind also:
\begin{equation}
 3a_0a_2-a_1^2=0,
 \qquad
 9a_0a_3-a_1a_2=0.
\tag{16}
\end{equation}
Ist $c_0$ nicht gleich Null, so gehen wir einen Schritt weiter und setzen:
\begin{equation}
 B=PC+D,
\tag{17}
\end{equation}
worin $D$ constant wird und den Ausdruck erhält:
\begin{equation}
 D={a_2c_0^2-2a_1c_0c_1+3a_0c_1^2\over c_0^2}.
\tag{18}
\end{equation}
Ist dieser Ausdruck von Null verschieden, so sind $A$ und $B$ theilerfremd, ist er gleich Null, so haben $A$ und $B$ den grössten gemeinschaftlichen Theiler $C$. Setzen wir für $c_0,c_1$ die Werthe (15) ein, lassen den Nenner weg, heben noch den Null verschwindenden Factor $9a_0$ heraus und kehren das Vorzeichen um, so erhält diese Bedingung nach einfacher Rechnung die Gestalt:
\begin{equation}
 a_1^2a_2^2+18a_0a_1a_2a_3-4a_0a_2^3-4a_1^3a_3-27a_0^2a_3^2=0.
\tag{19}
\end{equation}
Sie ist, wie man leicht durch Rechnung oder auch aus (18) sieht, auch dann erfüllt, wenn die Bedingungen (16) bestehen, und ist also die nothwendige und hinreichende Bedingung dafür, dass $f(x)$ und $f'(x)$ einen gemeinsamen Theiler haben. Die linke Seite von (19), die eine ganze rationale und homogene Function der Coefficienten von $f(x)$ ist, heisst die \emph{Discriminante} der Function $f(x)$.

Wir leiten noch einen Satz aus dem Algorithmus (3) her, der oft angewandt wird.

Aus der ersten dieser Gleichungen folgt:
\begin{equation}
 A''=A-Q'A',
\tag{20}
\end{equation}
und wenn man diesen Werth von $A''$ in die folgende Gleichung einsetzt:
\[
 A'''=(1+Q'Q'')A'-Q''A,
\]
also, wenn mit $p,p'$ ganze rationale Functionen bezeichnet werden:
\begin{equation}
 A'''=pA+p'A'.
\tag{21}
\end{equation}
Setzt man die Ausdrücke (20), (21) in die dritte Gleichung (3) ein, so ergiebt sich für $A''''$ wieder ein Ausdruck von der Form (21) und so kann man fortfahren, und erhält schliesslich:
\begin{equation}
 A^{(\nu)}=PA+P'A',
\tag{22}
\end{equation}
worin $P,P'$ ganze rationale Functionen sind, deren Coefficienten durch rationale Rechenoperationen aus den Coefficienten von $A$ und $A'$ zusammengesetzt sind.

In der Formel (22) ist $A^{(\nu)}$ eine Constante. Besonders wichtig ist dieser Satz in dem Falle, wo $A^{(\nu)}$ von Null verschieden, also $A,A'$ relativ prim sind. Setzen wir in diesem Falle
\[
 P=A^{(\nu)}F(x),\qquad P'=A^{(\nu)}\Phi(x),
\]
so können wir nach Weglassung des Factors $A^{(\nu)}$ dem Satz folgenden Ausdruck geben:

I. Sind $f(x)$ und $\psi(x)$ zwei ganze Functionen ohne gemeinsamen Theiler, so kann man zwei andere ganze Functionen $F(x)$ und $\Phi(x)$ bestimmen, die der Gleichung
\begin{equation}
 F(x)f(x)+\Phi(x)\psi(x)=1
\tag{23}
\end{equation}
identisch genügen.

Der Satz lässt sich noch verallgemeinern. Multipliciren wir die Gleichung (23) mit einer beliebigen ganzen Function $\chi(x)$, so folgt:
\begin{equation}
 F(x)\chi(x)f(x)+\Phi(x)\chi(x)\psi(x)=\chi(x),
\tag{24}
\end{equation}
und nach \S~3 können wir
\begin{equation}
 \Phi(x)\chi(x)=Q(x)f(x)+\varphi(x)
\tag{25}
\end{equation}
setzen, so dass $Q(x),\varphi(x)$ ganze rationale Functionen von $x$ sind und der Grad von $\varphi(x)$ kleiner ist als der von $f(x)$. Setzen wir dies in (24) ein und setzen an Stelle von $F(x)\chi(x)+Q(x)\psi(x)$ wieder $F(x)$, so erhalten wir:
\[
 F(x)f(x)+\varphi(x)\psi(x)=\chi(x).
\]
Wir können daher den vorigen Satz so verallgemeinern:

II. Sind $f(x),\psi(x),\chi(x)$ gegebene ganze rationale Functionen, und $f(x)$ und $\psi(x)$ ohne gemeinsamen Theiler, so lassen sich die ganzen rationalen Functionen $F(x),\varphi(x)$ und zwar $\varphi(x)$ von niedrigerem Grade als $f(x)$ so bestimmen, dass die Gleichung
\begin{equation}
 F(x)f(x)+\varphi(x)\psi(x)=\chi(x)
\tag{26}
\end{equation}
identisch befriedigt ist.

\section*{\S~7. Producte linearer Factoren.}

Nach \S~1 erhalten wir durch Multiplication ganzer Functionen ebensolche Functionen von höherem Grade, und zwar bestimmt sich der Grad des Productes durch die Summe der Grade der einzelnen Factoren.

Wenn wir also $n$ lineare Factoren mit einander multipliciren, so entsteht eine Function $n$ten Grades, deren Bildungsweise wir etwas genauer untersuchen müssen.

Wir wollen die linearen Factoren in der einfachen Form $x-\alpha_1$, $x-\alpha_2$, $\ldots$, $x-\alpha_n$ annehmen, und setzen, da der Coefficient der höchsten Potenz (nach \S~1) $=1$ ist,
\begin{equation}
\begin{aligned}
 f(x)&=(x-\alpha_1)(x-\alpha_2)\cdots(x-\alpha_n)\\
 &=x^n+a_1x^{n-1}+a_2x^{n-2}+\cdots+a_n.
\end{aligned}
\tag{1}
\end{equation}
Eine leichte Ueberlegung lässt folgendes Bildungsgesetz der Coefficienten $a_1,a_2\ldots a_n$ erkennen:

Es ist $-a_1$ gleich der Summe der $\alpha$, $a_2$ gleich der Summe der Producte von je zweien der $\alpha$, $-a_3$ die Summe der Producte von je dreien der $\alpha$, allgemein $(-1)^\nu a_\nu$ die Summe der Producte von je $\nu$ der Grössen $\alpha$, oder in Formeln ausgedrückt:
\begin{equation}
\begin{aligned}
-a_1&=\sum \alpha_1,\\
+a_2&=\sum \alpha_1\alpha_2,\\
&\cdots\\
(-1)^\nu a_\nu&=\sum \alpha_1\alpha_2\cdots\alpha_\nu,\\
&\cdots\\
(-1)^n a_n&=\alpha_1\alpha_2\cdots\alpha_n.
\end{aligned}
\tag{2}
\end{equation}
Um aber noch deutlicher die Richtigkeit dieses Bildungsgesetzes einzusehen, bedient man sich der vollständigen Induction.

Man bestätigt die Richtigkeit zunächst in den ersten Fällen durch wirkliche Ausführung der Multiplication
\[
 (x-\alpha_1)(x-\alpha_2)=x^2-(\alpha_1+\alpha_2)x+\alpha_1\alpha_2,
\]
\[
\begin{aligned}
 (x-\alpha_1)(x-\alpha_2)(x-\alpha_3)&=x^3-(\alpha_1+\alpha_2+\alpha_3)x^2\\
&\quad +(\alpha_1\alpha_2+\alpha_1\alpha_3+\alpha_2\alpha_3)x-\alpha_1\alpha_2\alpha_3.
\end{aligned}
\]
Nehmen wir die Richtigkeit unseres Bildungsgesetzes bei $n-1$ Factoren als bewiesen an und setzen
\[
 (x-\alpha_1)(x-\alpha_2)\cdots(x-\alpha_{n-1})=x^{n-1}+a'_1x^{n-2}+\cdots+a'_{n-1},
\]
so findet man durch Multiplication mit $x-\alpha_n$:
\begin{equation}
\begin{aligned}
 a_1&=a'_1-\alpha_n,\\
 a_2&=a'_2-\alpha_n a'_1,\\
 a_3&=a'_3-\alpha_n a'_2,\\
&\cdots\\
 a_\nu&=a'_\nu-\alpha_n a'_{\nu-1},\\
&\cdots\\
 a_n&=-\alpha_n a'_{n-1},
\end{aligned}
\tag{3}
\end{equation}
und daraus ist die Richtigkeit der Formeln (2) unmittelbar ersichtlich.

Es ist von Wichtigkeit, die Anzahl der Glieder zu bestimmen, die in jeder der Summen (2) vorkommen. Wir bezeichnen die Anzahl der Terme, die in der Summe $(-1)^\nu a_\nu$ vorkommen, mit $B_\nu^{(n)}$; es ist die Anzahl der Combinationen ohne Wiederholung von $n$ Elementen zur $\nu$ten Classe (d. h. von je $\nu$ verschiedenen der $n$ Elemente). Um sie zu bestimmen, denke man sich zunächst die $B_{\nu-1}^{(n)}$ Combinationen zur $(\nu-1)$ten Classe gebildet. Aus jeder dieser Combinationen kann man durch Hinzufügung je eines der fehlenden $n-\nu+1$ Elemente $n-\nu+1$ Combinationen zur $\nu$ten Classe ableiten. Auf diese Art aber wird jede Combination $\nu$mal, nämlich durch Hinzufügung jedes ihrer Elemente gebildet, so dass man die Relation
\begin{equation}
 B_\nu^{(n)}=B_{\nu-1}^{(n)}{n-\nu+1\over \nu}
\tag{4}
\end{equation}
erhält, während $B_1^{(n)}$ offenbar den Werth $n$ hat. Wenn man also die aus (4) folgenden Gleichungen
\begin{equation}
\begin{aligned}
 B_1^{(n)}&=n,\\
 B_2^{(n)}&=B_1^{(n)}{n-1\over 2},\\
&\cdots\\
 B_\nu^{(n)}&=B_{\nu-1}^{(n)}{n-\nu+1\over \nu}
\end{aligned}
\tag{5}
\end{equation}
multiplicirt, so folgt:
\begin{equation}
 B_\nu^{(n)}={n(n-1)(n-2)\cdots(n-\nu+1)\over 1\cdot 2\cdot 3\cdots \nu}.
\tag{6}
\end{equation}
Wir werden in der Folge oft, wenn $m$ eine beliebige positive ganze Zahl ist, das Zeichen
\begin{equation}
 \Pi(m)=1\cdot 2\cdot 3\cdots m,
 \qquad
 \Pi(0)=1
\tag{7}
\end{equation}
benutzen, so dass
\begin{equation}
 \Pi(m)=m\Pi(m-1).
\tag{8}
\end{equation}
Mit Hülfe dieses Zeichens lässt sich der Ausdruck für $B_\nu^{(n)}$ übersichtlicher so darstellen:
\begin{equation}
 B_\nu^{(n)}=B_{n-\nu}^{(n)}={\Pi(n)\over \Pi(\nu)\Pi(n-\nu)},
\tag{9}
\end{equation}
der, wenn $B_0^{(n)}=1$ gesetzt wird, auch noch für $\nu=0$ und $\nu=n$ gilt, und die Unveränderlichkeit von $B_\nu^{(n)}$ bei der Vertauschung von $\nu$ mit $n-\nu$ erkennen lässt.

Die Ableitung der Formel (4), die wir soeben nach den Vorschriften der Combinationslehre gegeben haben, ist zwar vollkommen richtig und einleuchtend, erfordert aber zur genauen Begründung einige Ueberlegung, die sich nicht gut in kurze Worte fassen lässt. Wir wollen daher nachträglich noch durch das Mittel der vollständigen Induction die Richtigkeit beweisen.

Die Formeln (3) geben nämlich die folgenden Recursionsformeln
\begin{equation}
\begin{aligned}
B_1^{(n)}&=B_1^{(n-1)}+1,\\
B_2^{(n)}&=B_2^{(n-1)}+B_1^{(n-1)},\\
&\cdots\\
B_\nu^{(n)}&=B_\nu^{(n-1)}+B_{\nu-1}^{(n-1)},\\
&\cdots\\
B_n^{(n)}&=B_{n-1}^{(n-1)}.
\end{aligned}
\tag{10}
\end{equation}
Nun lässt sich die Formel (6) für die ersten Werthe von $n$ sehr leicht durch Abzählen bestätigen. Nimmt man sie für $n-1$ als richtig an, so ergiebt (10):
\[
 B_\nu^{(n)}={\Pi(n-1)\over \Pi(\nu)\Pi(n-\nu-1)}+{\Pi(n-1)\over \Pi(\nu-1)\Pi(n-\nu)},
\]
woraus nach (8)
\[
 B_\nu^{(n)}={\Pi(n)\over \Pi(\nu)\Pi(n-\nu)},
\]
und hierdurch ist die Allgemeingültigkeit der Formel (9) bewiesen.

Aus der Bedeutung von $B_\nu^{(n)}$ ergiebt sich, und wird auch aus den Formeln (10) durch vollständige Induction bewiesen, dass es die $B_\nu^{(n)}$ positive ganze Zahlen sind.

\section*{\S~8. Der binomische Lehrsatz.}

Wenn wir in der Formel (1) des vorigen Paragraphen die bisher willkürlichen Grössen $\alpha_1,\alpha_2\ldots \alpha_n$ einander gleich setzen, so erhalten wir den binomischen Lehrsatz, der in nichts Anderem besteht, als in der Ordnung der $n$ten Potenz eines Binomiums $x+y$ nach Potenzen von $x$. Wenn wir nämlich
\[
 \alpha_1=\alpha_2=\cdots=\alpha_n=-y
\]
setzen, so ergiebt die Formel (2):
\[
 a_\nu=(-1)^\nu \sum \alpha_1\alpha_2\cdots\alpha_\nu=y^\nu B_\nu^{(n)},
\]
worin nach der Definition des vorigen Paragraphen $B_\nu^{(n)}$ die Anzahl der Terme der Summe bedeutet, die alle einander gleich und gleich $(-1)^\nu y^\nu$ werden.

Demnach ergiebt sich:
\begin{equation}
 (x+y)^n=x^n+B_1^{(n)}x^{n-1}y+B_2^{(n)}x^{n-2}y^2+\cdots+B_n^{(n)}y^n,
\tag{1}
\end{equation}
oder entwickelt geschrieben:
\[
\begin{aligned}
(x+y)^n
&=x^n+nx^{n-1}y+{n(n-1)\over 1\cdot2}x^{n-2}y^2+\cdots+y^n\\
&=\Pi(n)\sum {x^{n-\nu}y^\nu\over \Pi(n-\nu)\Pi(\nu)}
 =\Pi(n)\sum {x^\alpha y^\beta\over \Pi(\alpha)\Pi(\beta)},
\end{aligned}
\]
die letzte Summe über alle nicht negativen ganzzahligen Werthe $\alpha,\beta$ erstreckt, die der Bedingung $\alpha+\beta=n$ genügen.

Hiernach heissen die Coefficienten $B_\nu^{(n)}$ die Binomialcoefficienten.

Wir setzen der Uebersicht wegen eine kleine Tabelle der ersten Werthe der Binomialcoefficienten hierher:
\[
\begin{array}{c|rrrrrrrr}
 n=1&1&1\\
 n=2&1&2&1\\
 n=3&1&3&3&1\\
 n=4&1&4&6&4&1\\
 n=5&1&5&10&10&5&1\\
 n=6&1&6&15&20&15&6&1\\
 n=7&1&7&21&35&35&21&7&1
\end{array}
\]
Wir wollen unter den verschiedenen Eigenschaften der Binomialcoefficienten zwei ableiten, von denen wir nachher eine interessante Anwendung machen werden.

Aus (1) ergeben sich, wenn man $x,y$ durch $1,x$ ersetzt, und $n=0,1,2,3\ldots$ nimmt, die Formeln
\begin{equation}
\begin{aligned}
1&=B_0^{(0)},\\
1+x&=B_0^{(1)}+B_1^{(1)}x,\\
(1+x)^2&=B_0^{(2)}+B_1^{(2)}x+B_2^{(2)}x^2,\\
&\cdots\\
(1+x)^n&=B_0^{(n)}+B_1^{(n)}x+B_2^{(n)}x^2+\cdots+B_n^{(n)}x^n.
\end{aligned}
\tag{2}
\end{equation}
Wir machen nun von der bekannten Summenformel der geometrischen Reihe Gebrauch:
\[
 1+(1+x)+(1+x)^2+\cdots+(1+x)^n={(1+x)^{n+1}-1\over x},
\]
worin, wenn man $(1+x)^{n+1}$ wieder nach der Binomialformel entwickelt, indem man in der letzten Formel (2) $n$ in $n+1$ verwandelt und $B_0^{(n+1)}=1$ setzt:
\begin{equation}
{(1+x)^{n+1}-1\over x}=B_1^{(n+1)}+B_2^{(n+1)}x+\cdots+B_{n+1}^{(n+1)}x^n.
\tag{3}
\end{equation}
Vergleicht man dies mit der Summe der rechten Seiten von (2) und setzt den Vorschriften des \S~1 gemäss die Coefficienten entsprechender Potenzen von $x$ einander gleich, so folgt:
\begin{equation}
\begin{aligned}
 B_0^{(0)}+B_0^{(1)}+B_0^{(2)}+\cdots+B_0^{(n)}&=B_1^{(n+1)},\\
 B_1^{(1)}+B_1^{(2)}+\cdots+B_1^{(n)}&=B_2^{(n+1)},\\
&\cdots\\
 B_n^{(n)}&=B_{n+1}^{(n+1)},
\end{aligned}
\tag{4}
\end{equation}
oder allgemein
\begin{equation}
 B_\nu^{(\nu)}+B_\nu^{(\nu+1)}+\cdots+B_\nu^{(n)}=B_{\nu+1}^{(n+1)}.
\tag{5}
\end{equation}
Wenn man aber die Gleichungen (2) der Reihe nach mit
\[
B_0^{(n)},\quad -B_1^{(n)},\quad +B_2^{(n)},\ldots,\quad \pm B_n^{(n)}
\]
multiplicirt (wo das obere Zeichen bei geradem, das untere bei ungeradem $n$ gilt), so ergiebt die Summe der linken Seiten:
\[
 B_0^{(n)}-B_1^{(n)}(1+x)+B_2^{(n)}(1+x)^2-\cdots\pm B_n^{(n)}(1+x)^n
 =[1-(1+x)]^n=(-x)^n
\]
und die Gleichsetzung der Coefficienten gleich hoher Potenzen auf der rechten und linken Seite liefert das Formelsystem
\begin{equation}
\begin{aligned}
0&=B_0^{(0)}B_0^{(n)}-B_0^{(1)}B_1^{(n)}+B_0^{(2)}B_2^{(n)}-\cdots\pm B_0^{(n)}B_n^{(n)},\\
0&=-B_1^{(1)}B_1^{(n)}+B_1^{(2)}B_2^{(n)}-\cdots\pm B_1^{(n)}B_n^{(n)},\\
&\cdots\\
\pm 1&=\pm B_n^{(n)}B_n^{(n)},
\end{aligned}
\tag{6}
\end{equation}
und dies Formelsystem lässt sich in umgekehrter Reihenfolge mit Benutzung der Relation $B_\nu^{(n)}=B_{n-\nu}^{(n)}$ auch so darstellen:
\begin{equation}
\begin{aligned}
1&=B_0^{(n)}B_0^{(n)},\\
0&=B_1^{(n)}B_0^{(n)}-B_0^{(n-1)}B_1^{(n)},\\
0&=B_2^{(n)}B_0^{(n)}-B_1^{(n-1)}B_1^{(n)}+B_0^{(n-2)}B_2^{(n)},\\
&\cdots\\
0&=B_n^{(n)}B_0^{(n)}-B_{n-1}^{(n-1)}B_1^{(n)}+B_{n-2}^{(n-2)}B_2^{(n)}-\cdots\pm B_0^{(0)}B_n^{(n)},
\end{aligned}
\tag{7}
\end{equation}
oder auch in zusammenfassender Bezeichnung:
\begin{equation}
 \sum_{\nu=0}^{\mu}(-1)^\nu B_{\mu-\nu}^{(n-\nu)}B_\nu^{(n)}
 =\begin{cases}
 1,&\mu=0,\\
 0,&\mu=1,2,\ldots,n.
 \end{cases}
\tag{8}
\end{equation}


\clearpage
\section*{Band I. Erster Abschnitt. Rationale Functionen.}

\section*{\S~9. Interpolation.}

Wir wollen von den zuletzt gewonnenen Formeln eine Anwendung machen auf eine Aufgabe aus der Theorie der Interpolation.

Es soll eine ganze rationale Function $n$ten Grades bestimmt werden, die für $n+1$ gegebene Werthe der Veränderlichen $x$ vorgeschriebene Werthe hat.

Es handelt sich also um die Bestimmung der $n+1$ Coefficienten $a_0,a_1,\ldots,a_n$ in
\begin{equation}
 f(x)=a_0x^n+a_1x^{n-1}+\cdots+a_{n-1}x+a_n
\tag{1}
\end{equation}
aus den $n+1$ Gleichungen
\begin{equation}
 f(\alpha_0)=A_0,\quad f(\alpha_1)=A_1,\quad \ldots,\quad f(\alpha_n)=A_n,
\tag{2}
\end{equation}
wenn $\alpha_0,\alpha_1,\ldots,\alpha_n$ die gegebenen Werthe von $x$, und $A_0,A_1,\ldots,A_n$ die entsprechenden Werthe von $f(x)$ sind.

Die Gleichungen (2) sind ein System von Gleichungen ersten Grades für die Unbekannten $a_0,a_1,\ldots,a_n$, um die sich im Allgemeinen durch Determinanten auflösen lassen, wie wir im nächsten Abschnitt sehen werden. Wir kommen später auf diese Aufgabe zurück, behandeln sie aber jetzt unter der besonderen Voraussetzung, dass die gegebenen Werthe von $x$ die $n+1$ ersten ganzen Zahlen $0,1,2,\ldots,n$ seien.

Die Binomialcoefficienten $B_\nu^{(x)}$, wie sie durch \S~7, (6) definirt sind, behalten ihren guten Sinn, auch wenn $x$ keine ganze Zahl ist, wie bisher vorausgesetzt war, sondern eine beliebige veränderliche Grösse. Es ist dann
\begin{equation}
 B_\nu^{(x)}={x(x-1)\cdots(x-\nu+1)\over 1\cdot2\cdot3\cdots\nu}
\tag{3}
\end{equation}
eine ganze rationale Function $\nu$ten Grades von $x$.
\footnote{Die Bedeutung dieser verallgemeinerten Binomialcoefficienten für die Entwickelung der Potenzen des Binoms gehört nicht hierher, sondern in die Analysis.}
Der Coefficient von $x^\nu$ ist von Null verschieden und daher kann auch $x^\nu$ ausgedrückt werden durch $B_\nu^{(x)}$ und durch niedrigere Potenzen von $x$. Es lässt sich also auch jede ganze Function $n$ten Grades $f(x)$ von $x$ linear ausdrücken durch $B_0^{(x)},B_1^{(x)},\ldots,B_n^{(x)}$ in der Form
\begin{equation}
 f(x)=M_0B_0^{(x)}+M_1B_1^{(x)}+\cdots+M_nB_n^{(x)},
\tag{4}
\end{equation}
worin $M_0,M_1,\ldots,M_n$ Constanten sind, und die Function $f(x)$ ist bestimmt, wenn diese Constanten bestimmt sind.

Es sei nun nach unserer Voraussetzung $f(0),f(1),f(2),\ldots,f(n)$ gegeben; da nach (3) $B_\nu^{(x)}$ immer verschwindet, wenn $x$ einen der Werthe $0,1,2,\ldots,\nu-1$ hat, so ergeben sich aus (4) die folgenden linearen Gleichungen für die Unbekannten $M$:
\begin{equation}
\begin{aligned}
 f(0)&=M_0B_0^{(0)},\\
 f(1)&=M_0B_0^{(1)}+M_1B_1^{(1)},\\
 f(2)&=M_0B_0^{(2)}+M_1B_1^{(2)}+M_2B_2^{(2)},\\
 &\quad\cdots\\
 f(n)&=M_0B_0^{(n)}+M_1B_1^{(n)}+\cdots+M_nB_n^{(n)}.
\end{aligned}
\tag{5}
\end{equation}
Diese Gleichungen sind nun in Bezug auf $M_0,M_1,\ldots,M_n$ aufzulösen, was sehr leicht mit Hülfe der Schlussgleichungen (8) des vorigen Paragraphen geschieht. Die erste Gleichung (5) ergiebt nämlich direct
\begin{equation}
 M_0=f(0).
\tag{6}
\end{equation}
Multiplicirt man die erste Gleichung (5) mit $B_0^{(1)}$, die zweite mit $-B_1^{(1)}$ und addirt, so erhält man nach dem erwähnten Formelsystem (auf $n=1$ angewandt)
\[
 -M_1=B_0^{(1)}f(0)-B_1^{(1)}f(1),
\]
und so allgemein, indem man die $\nu$ ersten Gleichungen (5) der Reihe nach mit
\[
 B_0^{(\nu)},\quad -B_1^{(\nu)},\quad +B_2^{(\nu)},\quad \ldots,\quad \pm B_\nu^{(\nu)}
\]
multiplicirt und addirt,
\begin{equation}
 \pm M_\nu=B_0^{(\nu)}f(0)-B_1^{(\nu)}f(1)+B_2^{(\nu)}f(2)-\cdots\pm B_\nu^{(\nu)}f(\nu),
\tag{7}
\end{equation}
wodurch nach (4) die Function $f(x)$ bestimmt und die Aufgabe gelöst ist. Es ist klar, dass, so lange wir über die Werthe $f(0),f(1),\ldots,f(n)$ keine besondere Voraussetzung machen, in dieser Form jede beliebige ganze rationale Function von $x$ dargestellt werden kann.

\section*{\S~10. Lösung des Interpolationsproblems durch die Differenzen.}

Die Definition der ganzen rationalen Function
\begin{equation}
 B_\nu^{(x)}={x(x-1)\cdots(x-\nu+1)\over 1\cdot2\cdot3\cdots\nu}
\tag{1}
\end{equation}
giebt die früher schon für den speciellen Fall eines ganzen positiven $x$ bewiesene Relation
\begin{equation}
 B_\nu^{(x+1)}=B_\nu^{(x)}+B_{\nu-1}^{(x)},\qquad B_0^{(x+1)}=B_0^{(x)}=1.
\tag{2}
\end{equation}
Daraus erhalten wir für die Coefficienten $M_0,M_1,\ldots,M_n$ eine Bestimmungsweise, die für die praktische Rechnung viel bequemer ist, als die Anwendung der Formeln des vorigen Paragraphen.

Es ist nämlich nach den Formeln (4) und (6), \S~9
\begin{equation}
 f(x)=f(0)+M_1B_1^{(x)}+M_2B_2^{(x)}+\cdots+M_nB_n^{(x)},
\tag{3}
\end{equation}
und wenn wir darin $x$ durch $x+1$ ersetzen und die Differenz
\begin{equation}
 \Delta_x=f(x+1)-f(x)
\tag{4}
\end{equation}
bilden, mit Rücksicht auf (2),
\begin{equation}
 \Delta_x=M_1+M_2B_1^{(x)}+M_3B_2^{(x)}+\cdots+M_nB_{n-1}^{(x)},
\tag{5}
\end{equation}
woraus sich ergiebt (da (5) eine Gleichung von derselben Art wie (3) ist, nur dass $n-1$ an Stelle von $n$ getreten ist):
\[
 M_1=\Delta_0=f(1)-f(0).
\]
Setzen wir
\begin{equation}
 \Delta_x'=\Delta_{x+1}-\Delta_x,
 \qquad \Delta_x''=\Delta_{x+1}'-\Delta_x',\quad\ldots,
\tag{6}
\end{equation}
so wird also hiernach
\[
 M_0=f(0),\quad M_1=\Delta_0,
 \quad M_2=\Delta_0',\quad\ldots,\quad M_n=\Delta_0^{(n-1)},
\]
und die Formel (3) ergiebt
\begin{equation}
 f(x)=f(0)+\Delta_0B_1^{(x)}+\Delta_0'B_2^{(x)}+\Delta_0''B_3^{(x)}+\cdots+\Delta_0^{(n-1)}B_n^{(x)}.
\tag{7}
\end{equation}
und ebenso kann man die Functionen $\Delta_x,\Delta_x',\Delta_x'',\ldots$ ausdrücken:
\begin{equation}
\begin{aligned}
 \Delta_x&=\Delta_0+\Delta_0'B_1^{(x)}+\Delta_0''B_2^{(x)}+\cdots+\Delta_0^{(n-1)}B_{n-1}^{(x)},\\
 \Delta_x'&=\Delta_0'+\Delta_0''B_1^{(x)}+\cdots+\Delta_0^{(n-1)}B_{n-2}^{(x)}.
\end{aligned}
\tag{8}
\end{equation}
Die $\Delta_x,\Delta_x',\Delta_x'',\ldots,\Delta_x^{(n-1)}$ sind ganze rationale Functionen der Grade $n-1,n-2,\ldots,0$, die letzte von ihnen also constant. Um sie alle darzustellen, braucht man nur die Werthe
\[
 f(0),\quad \Delta_0,\quad \Delta_0',\quad \Delta_0'',\quad\ldots,\quad \Delta_0^{(n-1)},
\]
die man am leichtesten berechnet, wenn man eine Tabelle anlegt, die für den Fall $n=3$ z. B. folgende Form haben würde:
\[
\begin{array}{cccc}
f(0) &&&\\
f(1)&\Delta_0&&\\
f(2)&\Delta_1&\Delta_0'&\\
f(3)&\Delta_2&\Delta_1'&\Delta_0''
\end{array}
\]

\section*{\S~11. Arithmetische Reihen höherer Ordnung.}

Die vorstehenden Betrachtungen lassen sich noch in einer anderen Form aussprechen, die uns später von Nutzen sein wird. Ist
\begin{equation}
 u_0,u_1,u_2,\ldots
\tag{1}
\end{equation}
eine Reihe von Grössen, so nennen wir
\begin{equation}
 \Delta_0=u_1-u_0,
 \Delta_1=u_2-u_1,
 \Delta_2=u_3-u_2,\ldots
\tag{2}
\end{equation}
die Reihe ihrer Differenzen, mit
\begin{equation}
 \Delta_0'=\Delta_1-\Delta_0,
 \Delta_1'=\Delta_2-\Delta_1,\ldots
\tag{3}
\end{equation}
die Reihe ihrer zweiten Differenzen u. s. f.

Es ist klar, dass die Reihe (1) vollständig bestimmt ist, wenn ihr erstes Glied und die Reihe ihrer ersten Differenzen gegeben ist; denn es ist
\[
 u_1=u_0+\Delta_0,
 \quad u_2=u_0+\Delta_0+\Delta_1,
 \quad\ldots,
\]
\[
 u_m=u_0+\Delta_0+\Delta_1+\cdots+\Delta_{m-1}.
\]
Ebenso ist die Reihe (1) völlig bestimmt, wenn die beiden ersten Glieder und die Reihe ihrer zweiten Differenzen gegeben ist u. s. f. Die Reihe (1) wird eine arithmetische Reihe $n$ter Ordnung genannt, wenn die Reihe ihrer $n$ten Differenzen constant ist, also die Reihe der $(n+1)$ten Differenzen aus lauter Nullen besteht.

Man erhält eine arithmetische Reihe $n$ter Ordnung, wenn man in einer ganzen Function $n$ten Grades $f(x)$ für $x$ die Zahlen $0,1,2,3,\ldots$ einsetzt. Denn setzt man
\[
 \Delta_x=f(x+1)-f(x),\quad
 \Delta_x^{(n-1)}=\Delta_{x+1}^{(n-2)}-\Delta_x^{(n-2)},
\]
so ist $\Delta_x$ vom $(n-1)$ten, $\Delta_x'$ vom $(n-2)$ten Grade in Bezug auf $x$ und also $\Delta_x^{(n-1)}$ constant.

Ist nun
\[
 u_0,u_1,u_2,\ldots
\]
eine arithmetische Reihe $n$ter Ordnung, so ist die ganze Reihe vollständig bestimmt, wenn die $n$ Werthe $u_0,u_1,u_2,\ldots,u_{n-1}$ und ausserdem die constante $n$te Differenz gegeben sind. Diese letztere ist aber bestimmt, wenn auch noch das $(n+1)$te Glied $u_n$ gegeben ist. Wir können also den Satz aussprechen:

Eine arithmetische Reihe $n$ter Ordnung ist vollständig bestimmt, wenn ihre $n+1$ ersten Glieder gegeben sind.

Da nun eine Function $f(x)$ vom $n$ten Grade gleichfalls durch die willkürlich gegebenen Werthe
\[
 f(0),f(1),f(2),\ldots,f(n)
\]
völlig bestimmt ist, so folgt, dass aus den ganzen rationalen Functionen $n$ten Grades $f(x)$ alle arithmetischen Reihen $n$ter Ordnung erzeugt werden, wenn man darin für $x$ die Reihe der natürlichen Zahlen setzt.

Der Ausdruck des allgemeinen Gliedes ist dann durch die Formel (7) des vorigen Paragraphen gegeben.

Die Summen der $m+1$ ersten Glieder einer arithmetischen Reihe $n$ter Ordnung
\[
 s_m=u_0+u_1+\cdots+u_m
\]
bilden eine arithmetische Reihe $(n+1)$ter Ordnung, da ihre ersten Differenzen
\[
 s_{m+1}-s_m=u_{m+1}
\]
eine arithmetische Reihe $n$ter Ordnung bilden.

Es lässt sich also mit Hülfe der Formel (7) des \S~10 die Summe $s_m$ allgemein bestimmen, wenn man $s_0,s_1,\ldots,s_{n+1}$ als bekannt annimmt.

Um die erzeugende Function $F(x)$ von $s_m$ zu finden, wenn $f(x)$ die erzeugende Function von $u_m$ ist, setzt man
\[
 F(0)=f(0),\quad F(1)=f(0)+f(1),\quad F(2)=f(0)+f(1)+f(2),\ldots
\]
und hat dann in der Formel (7), \S~10,
\[
 F(x)=F(0)+D_0B_1^{(x)}+D_0'B_2^{(x)}+\cdots
\]
zu setzen
\[
\begin{aligned}
D_0&=F(1)-F(0)=f(1),& D_0'&=f(2)-f(1)=\Delta_1,\\
D_1&=F(2)-F(1)=f(2),& D_1'&=f(3)-f(2)=\Delta_2,\\
D_2&=F(3)-F(2)=f(3),& D_0''&=\Delta_2-\Delta_1=\Delta_1'.
\end{aligned}
\]
So erhält man
\begin{equation}
 F(x)=f(0)+f(1)B_1^{(x)}+\Delta_1B_2^{(x)}+\Delta_1'B_3^{(x)}+\cdots.
\tag{4}
\end{equation}
Nehmen wir z. B. $f(x)=x^2$, so giebt uns $F(m)$ die Summe der $m$ ersten Quadratzahlen. Es ist
\[
 \Delta_x=2x+1,
 \qquad \Delta_x'=2,
\]
also
\[
 F(x)=x+3{x(x-1)\over 1\cdot2}+2{x(x-1)(x-2)\over 1\cdot2\cdot3}
 ={x(x+1)(2x+1)\over6}.
\]
Für $f(x)=x^3$ ergiebt dieselbe Rechnung
\[
 F(x)=\left({x(x+1)\over2}\right)^2.
\]
Die Summe der $m$ ersten Cuben ist also gleich dem Quadrat der $m$ten Trigonalzahl.

\section*{\S~12. Der polynomische Lehrsatz.}

Im \S~8 ist für den binomischen Lehrsatz die Form abgeleitet:
\begin{equation}
 (x+y)^n=\Pi(n)\sum^{\alpha,\beta}{x^\alpha y^\beta\over \Pi(\alpha)\Pi(\beta)},
\tag{1}
\end{equation}
in der sich die Summe auf alle Combinationen zweier Zahlen $\alpha,\beta$ erstreckt, deren keine negativ ist und die der Bedingung
\begin{equation}
 \alpha+\beta=n
\tag{2}
\end{equation}
genügen.

Diese Form gestattet, zunächst durch Induction, eine Verallgemeinerung auf die $n$te Potenz eines Polynoms:
\begin{equation}
 (x+y+z+\cdots)^n=\Pi(n)\sum^{\alpha,\beta,\gamma,\ldots}
 {x^\alpha y^\beta z^\gamma\cdots
\over \Pi(\alpha)\Pi(\beta)\Pi(\gamma)\cdots},
\tag{3}
\end{equation}
mit der Bestimmung, dass $\alpha,\beta,\gamma,\ldots$ alle positiven oder verschwindenden ganzzahligen Werthe durchlaufen, die der Bedingung
\begin{equation}
 \alpha+\beta+\gamma+\cdots=n
\tag{4}
\end{equation}
genügen. Um aber die Richtigkeit dieser Formel allgemein zu beweisen, nehmen wir an, sie sei bewiesen, wenn das Polynom ein Glied weniger enthält, wie sie es in der That ist, wenn das Polynom nur zwei Glieder enthält.

Wir setzen dann
\begin{equation}
 u=y+z+\cdots
\tag{5}
\end{equation}
und wenden auf $(x+u)$ die Formel (1) an, aus der sich ergiebt:
\begin{equation}
 (x+y+z+\cdots)^n=\Pi(n)\sum^{\alpha,\nu}{x^\alpha u^\nu
\over \Pi(\alpha)\Pi(\nu)},
\tag{6}
\end{equation}
mit der Beschränkung
\begin{equation}
 \alpha+\nu=n.
\tag{7}
\end{equation}
Nun ist aber nach der Annahme schon bewiesen:
\begin{equation}
 u^\nu=\Pi(\nu)\sum^{\beta,\gamma,\ldots}
 {y^\beta z^\gamma\cdots
\over \Pi(\beta)\Pi(\gamma)\cdots},
\tag{8}
\end{equation}
\begin{equation}
 \beta+\gamma+\cdots=\nu,
\tag{9}
\end{equation}
und wenn dies in (6) eingesetzt wird, so ergiebt sich unmittelbar die Formel (3), und (7) geht in (4) über.

Die Coefficienten
\begin{equation}
 P_{\alpha,\beta,\gamma,\ldots}^{(n)}={\Pi(n)\over
 \Pi(\alpha)\Pi(\beta)\Pi(\gamma)\cdots},
\tag{10}
\end{equation}
die ihrer Bedeutung nach ganze Zahlen sind, heissen die Polynomialcoefficienten.

Beispielsweise erhält man für die dritte Potenz des Trinoms:
\begin{equation}
\begin{aligned}
(x+y+z)^3={}&x^3+y^3+z^3+3x^2y+3xy^2+3x^2z+3xz^2\\
&+3y^2z+3yz^2+6xyz.
\end{aligned}
\tag{11}
\end{equation}

\clearpage

\clearpage
\section*{Band I. Erster Abschnitt. Rationale Functionen.}

\section*{\S~13. Derivirte Functionen.}

Es sei
\begin{equation}
 f(x)=a_0x^n+a_1x^{n-1}+a_2x^{n-2}+\cdots+a_n
\tag{1}
\end{equation}
eine ganze rationale Function $n$ter Ordnung.

Wenn wir darin $x$ durch ein Binom $x+y$ ersetzen, so können wir auf jedes einzelne Glied den binomischen Lehrsatz anwenden, und können das Ergebniss nach fallenden oder nach steigenden Potenzen von $x$ oder von $y$ ordnen. Wir wollen die Ordnung nach steigenden Potenzen von $y$ ausführen. Die höchste Potenz von $y$, die vorkommt, ist die $n$te, und der Coefficient der nullten Potenz von $y$ ist die Function $f(x)$ selbst, wie man erkennt, wenn man $y=0$ setzt. Wir setzen also, indem wir die anderen Coefficienten mit
\[
 f'(x),\qquad {f''(x)\over \Pi(2)},\qquad {f'''(x)\over \Pi(3)},\ldots
\]
bezeichnen:
\begin{equation}
 f(x+y)=f(x)+yf'(x)+{y^2\over1\cdot2}f''(x)+\cdots
       =\sum_{0,n}^{\nu}{y^\nu\over\Pi(\nu)}\,f^{(\nu)}(x).
\tag{2}
\end{equation}

Die Functionen $f'(x),f''(x),f'''(x),\ldots$ heissen die erste, zweite, dritte, \ldots{} Derivirte oder Abgeleitete von $f(x)$. Es sind ganze Functionen von $x$ und $f^{(\nu)}(x)$ kann den Grad $n-\nu$ nicht übersteigen, da die Summe der Exponenten von $x$ und $y$ in keinem Gliede den Grad $n$ übersteigt.

Die erste Derivirte, die also der Coefficient der ersten Potenz von $y$ in der Entwickelung von $f(x+y)$ nach steigenden Potenzen von $y$ ist, erhält man durch Anwendung des binomischen Lehrsatzes auf (1):
\begin{equation}
 f'(x)=na_0x^{n-1}+(n-1)a_1x^{n-2}+(n-2)a_2x^{n-3}+\cdots .
\tag{3}
\end{equation}

Der Hauptsatz über die derivirten Functionen ergiebt sich aus (2), wenn wir $x$ in $x+z$ oder $y$ in $y+z$ verwandeln:
\begin{equation}
 f(x+y+z)=\sum_{0,n}^{\nu}{y^\nu\over\Pi(\nu)}f^{(\nu)}(x+z)
 =\sum_{0,n}^{\nu}{(y+z)^\nu\over\Pi(\nu)}f^{(\nu)}(x).
\tag{4}
\end{equation}
Bezeichnen wir mit $f^{(\nu,\mu)}(x)$ die $\mu$te Derivirte von $f^{(\nu)}(x)$, so ist nach (2):
\begin{equation}
 f^{(\nu)}(x+z)=\sum_{0,n-\nu}^{\mu}{z^\mu\over\Pi(\mu)}f^{(\nu,\mu)}(x)
\tag{5}
\end{equation}
und nach dem binomischen Satz:
\begin{equation}
 { (y+z)^\nu\over\Pi(\nu)}
 =\sum^{\beta,\gamma}{y^\beta z^\gamma\over\Pi(\beta)\Pi(\gamma)},
 \qquad \beta+\gamma=\nu.
\tag{6}
\end{equation}
Setzen wir dies in (4) ein, so folgt:
\begin{equation}
 \sum_{0,n}^{\nu}\sum_{0,n-\nu}^{\mu}
 {y^\nu z^\mu\over\Pi(\nu)\Pi(\mu)}f^{(\nu,\mu)}(x)
 =
 \sum^{\beta,\gamma}
 {y^\beta z^\gamma\over\Pi(\beta)\Pi(\gamma)}f^{(\beta+\gamma)}(x).
\tag{7}
\end{equation}
Die letzte Summe ist über alle nicht negativen Zahlen $\beta,\gamma$ zu erstrecken, deren Summe den Grad $n$ von $f(x)$ nicht übersteigt. Dieselben Zahlencombinationen durchlaufen aber auch die Exponenten $\nu,\mu$ auf der linken Seite, und die Vergleichung der Coefficienten gleicher Potenzen und Producte ergiebt (nach \S~1):
\begin{equation}
 f^{(\nu,\mu)}(x)=f^{(\nu+\mu)}(x),
\tag{8}
\end{equation}
also den Satz:

Die $\mu$te Derivirte von der $\nu$ten Derivirten ist die $(\nu+\mu)$te Derivirte der ursprünglichen Function.

Man erhält also die sämmtlichen höheren Derivirten, indem man nach der Regel (3) aus jeder vorangehenden die erste Derivirte bildet:
\begin{align}
 f(x)&=a_0x^n+a_1x^{n-1}+a_2x^{n-2}+a_3x^{n-3}+\cdots,\notag\\
 f'(x)&=na_0x^{n-1}+(n-1)a_1x^{n-2}+(n-2)a_2x^{n-3}+\cdots, \tag{9}\\
 f''(x)&=n(n-1)a_0x^{n-2}+(n-1)(n-2)a_1x^{n-3}+\cdots. \notag
\end{align}

Eine etwas einfachere Form nehmen diese Derivirten an, wenn man sich einer anderen Bezeichnungsweise bedient, die häufig im Gebrauch und für gewisse Zwecke fast unentbehrlich ist, die wir im Anschluss hieran besprechen wollen.

Es liegt wegen der Unbestimmtheit der Coefficienten $a_0,a_1,\ldots,a_n$ offenbar keine Beschränkung darin, wenn wir eine ganze rationale Function $n$ten Grades so darstellen:
\begin{equation}
 f(x)=a_0x^n+B_1^{(n)}a_1x^{n-1}+B_2^{(n)}a_2x^{n-2}+\cdots+a_n,
\tag{10}
\end{equation}
oder ausführlich:
\begin{equation}
 f(x)=a_0x^n+na_1x^{n-1}+{n(n-1)\over1\cdot2}a_2x^{n-2}+\cdots.
\tag{11}
\end{equation}
Wenn eine Function $f(x)$ so dargestellt ist, werden wir sagen, sie sei ``mit den Binomialcoefficienten geschrieben''.

Grössere Uebereinstimmung zeigen hierdurch bereits die Formeln (9), die dann so lauten:
\begin{align}
 f(x)&=a_0x^n+na_1x^{n-1}+{n(n-1)\over1\cdot2}a_2x^{n-2}+\cdots,\notag\\
 {1\over n}f'(x)&=a_0x^{n-1}+(n-1)a_1x^{n-2}+{(n-1)(n-2)\over1\cdot2}a_2x^{n-3}+\cdots, \tag{12}\\
 {1\over n(n-1)}f''(x)&=a_0x^{n-2}+(n-2)a_1x^{n-3}+{(n-2)(n-3)\over1\cdot2}a_2x^{n-4}+\cdots,\notag
\end{align}
worin die rechten Seiten alle auch mit den Binomialcoefficienten geschrieben erscheinen. Wir werden später den Nutzen dieser Bezeichnungsweise noch weiter kennen lernen, müssen aber schon hier hervorheben, dass die Wahl der einen oder anderen Darstellungsweise doch nicht ganz gleichgültig ist, die erste oft auch den Vorzug verdient. Besonders in den Fällen, wo die Coefficienten Zahlen sind und es auf das zahlentheoretische Verhalten dieser Coefficienten ankommt, darf man nicht ausser Acht lassen, dass durch die Binomialcoefficienten ein der Sache fremdes numerisches Element eingeführt wird. Dass Gauss in der Theorie der quadratischen Formen (in den Disq. ar.) die Schreibweise mit den Binomialcoefficienten anwendet, wenn er die quadratischen Formen durch $ax^2+2bxy+cy^2$ darstellt, und dass diese Bezeichnung allgemein Eingang gefunden hat, hat in der Zahlentheorie zu einer unnöthigen und sehr bedauerlichen Complication geführt.

\section*{\S~14. Derivirte eines Productes.}

Die derivirten Functionen, die wir hier betrachtet haben, sind keine anderen als die aus der Differentialrechnung bekannten Differentialquotienten; wir haben den Begriff aber hier, wo es sich um ganze rationale Functionen handelt, ohne Anwendung der Infinitesimalrechnung gewonnen aus den Entwickelungscoefficienten der Potenzen von $y$ in der Function $f(x+y)$. Bezeichnen wir die $\nu$te Ableitung von $f(x)$ mit $D_\nu f$, so ist nach (2), \S~13:
\begin{equation}
 f(x+y)=f(x)+yD_1f+{y^2\over1\cdot2}D_2f+{y^3\over1\cdot2\cdot3}D_3f+\cdots,
\tag{1}
\end{equation}
und daraus ergeben sich sofort die beiden Grundsätze, die sich in den Formeln
\begin{equation}
 D_\nu(Cf)=C D_\nu f;
\tag{2}
\end{equation}
\begin{equation}
 D_\nu(f+\varphi)=D_\nu f+D_\nu\varphi
\tag{3}
\end{equation}
ausdrücken, worin $C$ eine Constante, $\varphi$ eine zweite ganze rationale Function von $x$ ist.

Eine Verallgemeinerung der Formel (2) giebt die Darstellung der Derivirten des Productes $f\varphi$. Setzt man nämlich nach (1) abkürzend
\begin{align}
 f(x+y)&=u_0+yu_1+y^2u_2+\cdots+y^nu_n,\notag\\
 \varphi(x+y)&=v_0+yv_1+y^2v_2+\cdots+y^mv_m,
\tag{4}
\end{align}
also
\begin{equation}
 u_\nu={D_\nu f\over\Pi(\nu)},\qquad
 v_\nu={D_\nu\varphi\over\Pi(\nu)},
\tag{5}
\end{equation}
so ergiebt die Ausführung der Multiplication der beiden Formeln (4), wenn man in dem Product den Coefficienten von $y^\nu$ aufsucht:
\begin{equation}
 {D_\nu(f\varphi)\over\Pi(\nu)}
 =u_\nu v_0+u_{\nu-1}v_1+u_{\nu-2}v_2+\cdots+u_0v_\nu,
\tag{6}
\end{equation}
oder
\begin{equation}
 D_\nu(f\varphi)=\varphi D_\nu f+B_1^{(\nu)}D_{\nu-1}fD_1\varphi
 +B_2^{(\nu)}D_{\nu-2}fD_2\varphi+\cdots,
\tag{7}
\end{equation}
worin $B_1^{(\nu)},B_2^{(\nu)},\ldots$ die Binomialcoefficienten sind. In ähnlicher Weise kann man unter Anwendung der Polynomialcoefficienten die Derivirten eines Productes von mehr als zwei Factoren bilden.

Wir wollen die erste Derivirte, die wir jetzt mit $D$ statt mit $D_1$ bezeichnen, für ein Product von $n$ Factoren danach bilden. Für zwei Factoren erhalten wir nach (7):
\[
 D(f\varphi)=\varphi Df+fD\varphi=\varphi f'(x)+f\varphi'(x)
\]
und allgemein, wenn wir die Factoren mit $u_1,u_2,\ldots,u_n$, die Derivirten mit $u'_1,u'_2,\ldots,u'_n$ bezeichnen:
\begin{equation}
\begin{aligned}
D(u_1u_2\cdots u_n)
&=u'_1u_2\cdots u_n+u_1u'_2\cdots u_n+\cdots+u_1u_2\cdots u'_n,
\end{aligned}
\tag{8}
\end{equation}
oder kürzer:
\[
 {D(u_1u_2\cdots u_n)\over u_1u_2\cdots u_n}
 =\sum {D u_\nu\over u_\nu}.
\]

Wenn wir also ein Product aus linearen Factoren
\begin{equation}
 f(x)=(x-\alpha_1)(x-\alpha_2)\cdots(x-\alpha_n)
\tag{9}
\end{equation}
betrachten, so erhalten wir, da die ersten Derivirten von $x-\alpha_1$, $x-\alpha_2,\ldots,x-\alpha_n$ sämmtlich gleich $1$ sind,
\begin{equation}
\begin{aligned}
f'(x)&=(x-\alpha_2)(x-\alpha_3)\cdots(x-\alpha_n)\\
&\quad +(x-\alpha_1)(x-\alpha_3)\cdots(x-\alpha_n)+\cdots\\
&\quad +(x-\alpha_1)(x-\alpha_2)\cdots(x-\alpha_{n-1}),
\end{aligned}
\tag{10}
\end{equation}
wofür man auch setzen kann:
\begin{equation}
 f'(x)={f(x)\over x-\alpha_1}+{f(x)\over x-\alpha_2}+\cdots+{f(x)\over x-\alpha_n}.
\tag{11}
\end{equation}
Ein sehr wichtiges Resultat ergiebt sich hieraus, wenn $x$ gleich einem der Werthe $\alpha_1,\alpha_2,\ldots,\alpha_n$ gesetzt wird, nämlich
\begin{equation}
\begin{aligned}
 f'(\alpha_1)&=(\alpha_1-\alpha_2)(\alpha_1-\alpha_3)\cdots(\alpha_1-\alpha_n),\\
 f'(\alpha_2)&=(\alpha_2-\alpha_1)(\alpha_2-\alpha_3)\cdots(\alpha_2-\alpha_n),\\
 &\quad\cdots\\
 f'(\alpha_n)&=(\alpha_n-\alpha_1)(\alpha_n-\alpha_2)\cdots(\alpha_n-\alpha_{n-1}).
\end{aligned}
\tag{12}
\end{equation}

\section*{\S~15. Ganze Functionen mehrerer Veränderlichen: Formen.}

Wir haben bisher vorzugsweise die ganzen rationalen Functionen von einer Veränderlichen betrachtet; wir können uns aber nicht immer darauf beschränken und haben ja auch schon oben Functionen mehrerer Veränderlichen benutzt. Unter einer ganzen rationalen Function $n$ten Grades mehrerer Veränderlichen $F(x,y,z,\ldots)$ verstehen wir eine Summe von Gliedern:
\[
 \sum A_{\alpha,\beta,\gamma,\ldots}x^\alpha y^\beta z^\gamma\cdots,
\]
worin $\alpha,\beta,\gamma,\ldots$ ganzzahlige nicht negative Exponenten sind, deren Summe $\alpha+\beta+\gamma+\cdots$ den Werth $n$ nicht übersteigt, und wenigstens in einem Gliede auch wirklich erreicht. Der Grad wird also bestimmt durch den grössten Werth, den die Summe $\alpha+\beta+\gamma+\cdots$ annimmt.

Wenn die Summe der Exponenten $\alpha+\beta+\gamma+\cdots$ in allen Gliedern denselben Werth hat, so heisst die Function homogen.

Eine fundamentale Eigenschaft der homogenen Functionen $n$ten Grades ist die, dass, wenn alle Variablen mit demselben Factor vervielfältigt werden, der Erfolg derselbe ist, wie wenn die Function mit der $n$ten Potenz vervielfältigt wird; in Zeichen, wenn $t$ eine beliebige Veränderliche bedeutet:
\begin{equation}
 F(tx,ty,tz,\ldots)=t^nF(x,y,z,\ldots);
\tag{1}
\end{equation}
denn ersetzt man in dem Product $x^\alpha y^\beta z^\gamma\cdots$ die Variablen durch $tx,ty,tz,\ldots$, so erhält es den Factor
\[
 t^{\alpha+\beta+\gamma+\cdots};
\]
hat nun $\alpha+\beta+\gamma+\cdots$ in allen Gliedern denselben Werth $n$, so kann der Factor $t^n$ vor die Summe herausgenommen werden. Hat aber die Summe $\alpha+\beta+\gamma+\cdots$ in den einzelnen Gliedern verschiedene Werthe, so kann ein solcher gemeinschaftlicher Factor nicht herausgenommen werden, wenigstens nicht, ohne dass noch verschiedene Potenzen von $t$ in den einzelnen Gliedern bleiben.

Durch Vermehrung der Veränderlichen kann man jede nicht homogene Function in eine homogene von gleichem Grade verwandeln. Ist nämlich $m-1$ die Anzahl der Variablen in einer nicht homogenen Function $n$ten Grades, so setzen wir
\[
 x={x_1\over x_m},\qquad y={x_2\over x_m},\qquad z={x_3\over x_m},\ldots,
\]
und erhalten in
\[
 x_m^n F\!\left({x_1\over x_m},{x_2\over x_m},{x_3\over x_m},\ldots\right)
\]
eine ganze homogene Function $n$ten Grades der Variablen $x_1,x_2,\ldots,x_m$, die wir mit
\[
 \Phi(x_1,x_2,\ldots,x_m)
\]
bezeichnen.

Es empfiehlt sich bisweilen, die homogenen Functionen mehrerer Variablen mit den Polynomialcoefficienten zu schreiben. Wir setzen daher
\begin{equation}
 \Phi(x_1,x_2,\ldots,x_m)=
 \sum{\Pi(n)\over\Pi(\alpha_1)\Pi(\alpha_2)\cdots\Pi(\alpha_m)}
 A_{\alpha_1,\alpha_2,\ldots,\alpha_m}
 x_1^{\alpha_1}x_2^{\alpha_2}\cdots x_m^{\alpha_m},
\tag{2}
\end{equation}
wo sich die Summe auf alle nicht negativen, der Bedingung
\begin{equation}
 \alpha_1+\alpha_2+\cdots+\alpha_m=n
\tag{3}
\end{equation}
genügenden Zahlen erstreckt. Diese Bezeichnungsweise, ohne die Beschränkung (3), ist auch auf nicht homogene Functionen anwendbar.

Man kann aber die homogene Function auch so darstellen:
\begin{equation}
 \Phi(x_1,x_2,\ldots,x_m)=
 \sum A_{\nu_1,\nu_2,\ldots,\nu_n}x_{\nu_1}x_{\nu_2}\cdots x_{\nu_n},
\tag{4}
\end{equation}
worin jeder der Indices $\nu_1,\nu_2,\ldots,\nu_n$ von den übrigen unabhängig die Werthreihe $1,2,\ldots,m$ zu durchlaufen hat. Die Summe (4) besteht also aus $m^n$ Gliedern, die aber nicht alle von einander verschieden sind. Das Product $x_{\nu_1}x_{\nu_2}\cdots x_{\nu_n}$ bleibt nämlich ungeändert, wenn die Indices $\nu_1,\nu_2,\ldots,\nu_n$ beliebig unter einander permutirt werden. Die Anzahl der Permutationen von $n$ Elementen beträgt aber $\Pi(n)$. Sind unter diesen Elementen je $\alpha_1,\alpha_2,\ldots$ einander gleich, so reducirt sich die Zahl der Permutationen auf
\[
 {\Pi(n)\over \Pi(\alpha_1)\Pi(\alpha_2)\cdots},
\]
woraus sich ergiebt, dass in (4) irgend ein Product $x_1^{\alpha_1}x_2^{\alpha_2}\cdots$ genau
\[
 {\Pi(n)\over \Pi(\alpha_1)\Pi(\alpha_2)\cdots}
\]
mal vorkommt. Setzt man also noch fest, dass $A_{\nu_1,\nu_2,\ldots,\nu_n}$ sich nicht ändern soll, wenn die Indices beliebig permutirt werden, so erweisen sich die Bezeichnungsweisen (2) und (4) als identisch, wenn durch Zusammenfassen gleicher Factoren
\[
 x_{\nu_1}x_{\nu_2}\cdots x_{\nu_n}=x_1^{\alpha_1}x_2^{\alpha_2}\cdots x_m^{\alpha_m}
\]
und
\[
 A_{\nu_1,\nu_2,\ldots,\nu_n}=A_{\alpha_1,\alpha_2,\ldots,\alpha_m}
\]
gesetzt wird.

Bezeichnen wir die Anzahl der Glieder, die in der Function $\Phi$ [nach (2)] auftreten, mit $(m,n)$, so findet man, indem man zunächst die Glieder zählt, die den Factor $x_1$ haben und dann die übrigen, die eine homogene Function $n$ter Ordnung von den übrigen $m-1$ Variablen bilden, die Recursionsformel
\begin{equation}
 (m,n)=(m,n-1)+(m-1,n),
\tag{5}
\end{equation}
mit deren Hülfe man durch vollständige Induction den Ausdruck
\begin{equation}
 (m,n)={m(m+1)\cdots(m+n-1)\over 1\cdot2\cdots n}
 ={ \Pi(m+n-1)\over \Pi(n)\Pi(m-1)}
\tag{6}
\end{equation}
als richtig erweist.

Die ganzen homogenen Functionen werden auch Formen genannt. Man unterscheidet nach der Anzahl der Variablen unäre (einfache Potenzen), binäre, ternäre, quaternäre Formen. Die binären Formen sind es, die uns hier besonders interessiren, deren Theorie im Wesentlichen identisch ist mit der Theorie der ganzen rationalen Functionen einer Veränderlichen. Man gelangt von den binären Formen zu diesen Functionen zurück, wenn man eine der homogenen Variablen als constant ansieht, z. B. ihr den Werth $1$ giebt.

\section*{\S~16. Die Derivirten von Functionen mehrerer Variablen.}

Wir haben im \S~13 die derivirten Functionen einer ganzen rationalen Function einer Veränderlichen definirt. Der Begriff lässt sich unmittelbar übertragen auf Functionen mehrerer Variablen, indem man die Ableitungen in Bezug auf jede Variable für sich, als ob sie die einzige wäre, bildet. So erhält man, wenn man etwa wie in \S~15
\begin{equation}
 F(x,y,z,\ldots)=\sum A_{\alpha,\beta,\gamma,\ldots}x^\alpha y^\beta z^\gamma\cdots
\tag{1}
\end{equation}
setzt, die erste Derivirte nach $x$:
\begin{equation}
 F'(x)=\sum \alpha A_{\alpha,\beta,\gamma,\ldots}x^{\alpha-1}y^\beta z^\gamma\cdots,
\tag{2}
\end{equation}
oder nach $y$:
\begin{equation}
 F'(y)=\sum \beta A_{\alpha,\beta,\gamma,\ldots}x^\alpha y^{\beta-1}z^\gamma\cdots
\tag{3}
\end{equation}
u. s. f. Aus diesen Functionen kann man nach denselben Regeln wieder die Ableitungen nach den verschiedenen Variablen bilden und erhält so die höheren Ableitungen.

Um die Resultate übersichtlicher darzustellen, sei $\Phi(x_1,x_2,\ldots,x_m)$ eine ganze rationale Function $n$ter Ordnung der $m$ Veränderlichen $x_1,x_2,\ldots,x_m$. Wir ersetzen diese Veränderlichen durch Binome:
\[
 x_1+\xi_1,\quad x_2+\xi_2,\ldots,\quad x_m+\xi_m
\]
und entwickeln in jedem Gliede der Function
\[
 \Phi(x_1+\xi_1,x_2+\xi_2,\ldots,x_m+\xi_m)=\Phi(x+\xi)
\]
durch Ausführung der Multiplication
\[
 (x_1+\xi_1)^{\mu_1}(x_2+\xi_2)^{\mu_2}\cdots(x_m+\xi_m)^{\mu_m}
\]
nach Potenzen von $\xi_1,\xi_2,\ldots,\xi_m$. Fassen wir gleiche Potenzen und Producte der Variablen $\xi$ je in ein Glied zusammen, so ergiebt sich in der Bezeichnung (2) \S~15 für $\Phi(x+\xi)$ eine Darstellung, die in der Differentialrechnung die Taylor'sche Entwickelung heisst:
\begin{equation}
 \Phi(x+\xi)=\sum
 {\xi_1^{\alpha_1}\xi_2^{\alpha_2}\cdots\xi_m^{\alpha_m}
 \over \Pi(\alpha_1)\Pi(\alpha_2)\cdots\Pi(\alpha_m)}
 D_{\alpha_1,\alpha_2,\ldots,\alpha_m}\Phi .
\tag{4}
\end{equation}
Die Coefficienten, die wir mit
\[
 D_{\alpha_1,\alpha_2,\ldots,\alpha_m}\Phi
\]
bezeichnen, sind Functionen der Variablen $x$ und heissen, wenn
\[
 \alpha_1+\alpha_2+\cdots+\alpha_m=\nu
\]
ist, die Derivirten $\nu$ter Ordnung der Function $\Phi$.

Man stellt sie auch nach der in der Differentialrechnung gebräuchlichen Bezeichnungsweise so dar:
\begin{equation}
 D_{\alpha_1,\alpha_2,\ldots,\alpha_m}\Phi
 =
 {\partial^\nu\Phi\over
 \partial x_1^{\alpha_1}\partial x_2^{\alpha_2}\cdots\partial x_m^{\alpha_m}}.
\tag{5}
\end{equation}

Das Bildungsgesetz der Derivirten lässt sich in folgende Sätze zusammenfassen, wobei wir der Kürze wegen die Indices bei dem Zeichen $D$ weglassen.

I. Ist $C$ eine Constante, so ist
\[
 D(C\Phi)=CD\Phi.
\]
II. Sind $\Phi$ und $\Psi$ irgend zwei Functionen, so ist
\[
 D(\Phi+\Psi)=D\Phi+D\Psi.
\]
Beides folgt unmittelbar aus (4).

Wir können also leicht die derivirten Functionen allgemein bilden, wenn wir sie für den speciellen Fall kennen, in dem $\Phi$ ein Product von Potenzen ist, also wenn wir
\[
 D_{\alpha_1,\alpha_2,\ldots,\alpha_m}(x_1^{\mu_1}x_2^{\mu_2}\cdots x_m^{\mu_m})
\]
kennen, worin die $\mu$ beliebige, nicht negative Exponenten sind.

Nun ist aber
\[
 (x_1+\xi_1)^{\mu_1}
 =\sum^{\alpha_1}{\Pi(\mu_1)\over\Pi(\alpha_1)\Pi(\mu_1-\alpha_1)}
 \xi_1^{\alpha_1}x_1^{\mu_1-\alpha_1}
\]
und folglich:
\begin{equation}
\begin{aligned}
&(x_1+\xi_1)^{\mu_1}\cdots(x_m+\xi_m)^{\mu_m}\\
&=\sum^{\alpha_1\ldots\alpha_m}
 {\Pi(\mu_1)\cdots\Pi(\mu_m)x_1^{\mu_1-\alpha_1}\cdots x_m^{\mu_m-\alpha_m}
 \over
 \Pi(\mu_1-\alpha_1)\cdots\Pi(\mu_m-\alpha_m)\Pi(\alpha_1)\cdots\Pi(\alpha_m)}
 \xi_1^{\alpha_1}\cdots\xi_m^{\alpha_m}.
\end{aligned}
\tag{6}
\end{equation}
und es ergiebt also die Vergleichung mit (4) und (6)
\begin{equation}
\begin{aligned}
&D_{\alpha_1,\ldots,\alpha_m}(x_1^{\mu_1}\cdots x_m^{\mu_m})\\
&\quad =
{\Pi(\mu_1)\cdots\Pi(\mu_m)\over
 \Pi(\mu_1-\alpha_1)\cdots\Pi(\mu_m-\alpha_m)}
 x_1^{\mu_1-\alpha_1}\cdots x_m^{\mu_m-\alpha_m},
\end{aligned}
\tag{7}
\end{equation}
so lange $\alpha_1\leq\mu_1,\ldots,\alpha_m\leq\mu_m$. Dagegen ist
\begin{equation}
 D_{\alpha_1,\ldots,\alpha_m}(x_1^{\mu_1}\cdots x_m^{\mu_m})=0,
\tag{8}
\end{equation}
sobald einer der Indices $\alpha$ grösser ist als der entsprechende Exponent $\mu$.

Bezeichnen wir nun mit $\beta_1,\beta_2,\ldots,\beta_m$ ein zweites System von Indices, und bilden von der Function (7) die Ableitung $D_{\beta_1,\beta_2,\ldots,\beta_m}$, so ergiebt sich durch nochmalige Anwendung derselben Formeln (7) und (8):
\begin{equation}
 D_{\beta_1,\ldots,\beta_m}D_{\alpha_1,\ldots,\alpha_m}
 (x_1^{\mu_1}\cdots x_m^{\mu_m})
 =
 D_{\beta_1+\alpha_1,\ldots,\beta_m+\alpha_m}
 (x_1^{\mu_1}\cdots x_m^{\mu_m}),
\tag{9}
\end{equation}
und daraus folgt nach II. die allgemeine Gültigkeit des Satzes
\[
\text{III.}\qquad
D_{\beta_1,\ldots,\beta_m}D_{\alpha_1,\ldots,\alpha_m}\Phi
 =
D_{\beta_1+\alpha_1,\ldots,\beta_m+\alpha_m}\Phi,
\]
was eine Verallgemeinerung des Satzes (8), \S~13 ist, und man kann also die höheren Derivirten durch fortgesetzte Ableitung der niederen bilden.

Für die ersten Derivirten einer Function $\Phi$
\[
D_{1,0,\ldots,0}\Phi,\quad D_{0,1,\ldots,0}\Phi,\ldots,D_{0,0,\ldots,1}\Phi
\]
brauchen wir auch die kürzeren Zeichen
\[
\Phi'(x_1),\quad \Phi'(x_2),\ldots,\Phi'(x_m),
\]
ebenso für die zweiten
\[
D_{2,0,\ldots,0},\quad D_{1,1,\ldots,0},\ldots
\]
die Zeichen
\[
\Phi''(x_1,x_1),\quad \Phi''(x_1,x_2)=\Phi''(x_2,x_1),\ldots.
\]
Auch diese Bezeichnung lässt sich verallgemeinern und würde zu einem der Formel (4), \S~15 entsprechenden Ausdruck führen.

Wir erwähnen des häufigen Gebrauches wegen die Formeln für die quadratischen Formen besonders. Setzen wir
\begin{equation}
 \Phi(x)=\sum a_{i,k}x_i x_k,
\tag{10}
\end{equation}
worin $i,k$ von einander unabhängig die Reihe der Zahlen $1,2,\ldots,m$ durchlaufen und $a_{i,k}=a_{k,i}$ ist, so ist:
\begin{equation}
\begin{aligned}
 {1\over2}\Phi'(x_1)&=a_{1,1}x_1+a_{1,2}x_2+\cdots+a_{1,m}x_m,\\
 {1\over2}\Phi'(x_2)&=a_{2,1}x_1+a_{2,2}x_2+\cdots+a_{2,m}x_m,\\
 &\quad\cdots\\
 {1\over2}\Phi'(x_m)&=a_{m,1}x_1+a_{m,2}x_2+\cdots+a_{m,m}x_m,
\end{aligned}
\tag{11}
\end{equation}
und wir setzen noch
\begin{equation}
 \Phi(x,\xi)=\Phi(\xi,x)=\sum \xi_i\Phi'(x_i)=\sum x_i\Phi'(\xi_i).
\tag{12}
\end{equation}
Dann ist
\begin{equation}
 \Phi(x+\xi)=\Phi(x)+\Phi(x,\xi)+\Phi(\xi).
\tag{13}
\end{equation}
Die Function $\Phi(x,\xi)$ wird die Polare von $\Phi$ genannt. Sie ist linear und homogen sowohl in Beziehung auf die $x$, wie in Beziehung auf die $\xi$.

Sie kann ausgedrückt werden durch
\begin{equation}
 \Phi(x,\xi)=2\sum a_{ik}\xi_i x_k
\tag{14}
\end{equation}
und genügt der Bedingung
\begin{equation}
 \Phi(x,x)=2\Phi(x).
\tag{15}
\end{equation}

\section*{\S~17. Das Euler'sche Theorem über homogene Functionen.}

Aus den vorstehenden Entwickelungen lässt sich mit Leichtigkeit ein Fundamentalsatz über homogene Functionen herleiten, der von Euler entdeckt und nach ihm benannt ist.

Wir erhalten ihn am einfachsten aus der Formel (4) des vorigen Paragraphen, wenn wir mit $t$ eine beliebige Veränderliche bezeichnen,
\begin{equation}
 \xi_1=tx_1,\quad \xi_2=tx_2,\ldots,\quad \xi_m=tx_m
\tag{1}
\end{equation}
setzen und dann die Fundamentalformel \S~15 (1) für die homogenen Functionen anwenden. Wir erhalten so zunächst:
\begin{equation}
 (1+t)^n\Phi(x)
 =
 \sum { (tx_1)^{\alpha_1}(tx_2)^{\alpha_2}\cdots(tx_m)^{\alpha_m}
 \over \Pi(\alpha_1)\Pi(\alpha_2)\cdots\Pi(\alpha_m)}
 D_{\alpha_1,\alpha_2,\ldots,\alpha_m}\Phi.
\tag{2}
\end{equation}
Wendet man auf der linken Seite von (2) den binomischen Satz an, und setzt dann die Coefficienten gleich hoher Potenzen von $t$ beiderseits einander gleich, so ergiebt sich für jedes $\nu=1,2,\ldots,n$:
\begin{equation}
 {\Pi(n)\over\Pi(n-\nu)}\Phi(x_1,x_2,\ldots,x_m)
 =
 \sum^{\alpha}
 {\Pi(\nu)\over\Pi(\alpha_1)\Pi(\alpha_2)\cdots\Pi(\alpha_m)}
 x_1^{\alpha_1}x_2^{\alpha_2}\cdots x_m^{\alpha_m}
 D_{\alpha_1,\alpha_2,\ldots,\alpha_m}\Phi,
\tag{3}
\end{equation}
worin sich die Summe auf alle der Bedingung
\begin{equation}
 \alpha_1+\alpha_2+\cdots+\alpha_m=\nu
\tag{4}
\end{equation}
genügenden Werthsysteme der $\alpha$ erstreckt.

In dieser Form ist das zu erweisende Theorem in seiner Allgemeinheit enthalten. Für den besonderen Fall $\nu=1$ erhalten wir die Formel
\begin{equation}
 n\Phi(x_1,x_2,\ldots,x_m)
 =x_1\Phi'(x_1)+x_2\Phi'(x_2)+\cdots+x_m\Phi'(x_m),
\tag{5}
\end{equation}
wovon die Formel (15) des vorigen Paragraphen ein specieller Fall ist, und für $\nu=2$:
\begin{equation}
 n(n-1)\Phi(x_1,x_2,\ldots,x_m)
 =
 \sum_{i,k}^{m}x_i x_k\Phi''(x_i,x_k),
\tag{6}
\end{equation}
worin die Summe von $i=1$ bis $i=m$ und von $k=1$ bis $k=m$ zu erstrecken ist, so dass jedes Glied mit ungleichen $i,k$ zweimal in der Summe auftritt.

Setzen wir, wenn die $\alpha$ der Bedingung (4) unterworfen sind,
\begin{equation}
 \Phi_\nu(\xi,x)=
 \sum {\xi_1^{\alpha_1}\xi_2^{\alpha_2}\cdots\xi_m^{\alpha_m}
 \over \Pi(\alpha_1)\Pi(\alpha_2)\cdots\Pi(\alpha_m)}
 D_{\alpha_1,\alpha_2,\ldots,\alpha_m}\Phi,
\tag{7}
\end{equation}
so ist nach (4) des \S~16:
\begin{equation}
 \Phi(x+\xi)=\Phi(x)+\Phi_1(x,\xi)+\Phi_2(x,\xi)+\cdots+\Phi_n(x,\xi),
\tag{8}
\end{equation}
und da die linke Seite ungeändert bleibt, wenn $x$ mit $\xi$ vertauscht wird, so ergiebt sich die Relation:
\begin{equation}
 \Phi_{n-\nu}(x,\xi)=\Phi_\nu(\xi,x),
\tag{9}
\end{equation}
also insbesondere
\begin{equation}
 \Phi_n(x,\xi)=\Phi(\xi).
\tag{10}
\end{equation}
Die Function $\Phi_\nu(x,\xi)$ wird, als Function von $x$ betrachtet, die $\nu$te Polare der Function $\Phi$ für das Werthsystem $\xi$ genannt.

Wir wollen die Formel (3) noch für den Fall einer binären Form ($m=2$) specialisiren. Wir bezeichnen die Variablen mit $x,y$ und setzen zur Abkürzung:
\[
 \Phi(x,y)=u,\qquad D_{h,\nu-h}\Phi=u_h
\]
und erhalten aus (4):
\begin{equation}
 {\Pi(n)\over\Pi(n-\nu)}u
 =
 \sum_{0,\nu}^{h}{\Pi(\nu)\over\Pi(h)\Pi(\nu-h)}
 u_h x^h y^{\nu-h},
\tag{9}
\end{equation}
worin $\nu$ jeden beliebigen Werth, der nicht grösser als $n$ ist, annehmen kann.


\clearpage
\begin{center}
{\large Zweiter Abschnitt.}\\[0.5em]
{\Large Determinanten.}\\[0.4em]
\end{center}

\section*{\S\,18. Permutationen von \(n\) Elementen.}

Wir betrachten ein System von \(n\) unterschiedenen Elementen irgend welcher Art, z.~B. die \(n\) Ziffern
\[
1,2,3,\ldots,n,
\]
deren Complex in dieser bestimmten Anordnung wir mit \(\mathfrak A\) bezeichnen wollen. Die Elemente von \(\mathfrak A\) lassen sich auf verschiedene Arten anordnen, z.~B.
\[
2,1,3,\ldots,n.
\]
Der Uebergang von einer Anordnung zu einer anderen heisst eine Permutation.

Bezeichnen wir die Anzahl der verschiedenen Anordnungen, die nur von der Anzahl \(n\) der Elemente abhängen kann, mit \(\Pi(n)\), so ergiebt sich zunächst
\[
\Pi(1)=1,\qquad \Pi(2)=2,
\]
und um die Zahl allgemein zu bestimmen, denken wir uns zu \(n-1\) Elementen ein \(n\)tes hinzugefügt. In jeder Anordnung der \(n-1\) Elemente kann nun das \(n\)te Element an \(n\) verschiedene Stellen gesetzt werden, nämlich vor das erste, zwischen das erste und zweite, zwischen das zweite und dritte u.~s.~f., endlich nach dem \((n-1)\)ten, und alle die so entstandenen Anordnungen sind von einander verschieden. Daraus folgt:
\begin{equation}
\Pi(n)=n\Pi(n-1),
\tag{1}
\end{equation}
woraus sich durch vollständige Induction
\begin{equation}
\Pi(n)=1\cdot2\cdot3\cdots n
\tag{2}
\end{equation}
ergiebt, so dass das Zeichen \(\Pi(n)\) hier dieselbe Bedeutung hat, wie im ersten Abschnitt (\S~7).

Irgend eine Anordnung des Systems \(\mathfrak A\) bezeichnen wir mit \(\mathfrak A'\), oder ausführlicher, wenn \(\alpha_1,\alpha_2,\ldots,\alpha_n\) die Ziffern \(1,2,\ldots,n\) in irgend einer Reihenfolge bedeuten, mit
\begin{equation}
\mathfrak A'=\alpha_1,\alpha_2,\ldots,\alpha_n.
\tag{3}
\end{equation}

Man kann auf sehr verschiedene Arten aus einer Anordnung eine beliebige andere ableiten, d.~h. eine Permutation ausführen. Unter den verschiedenen Möglichkeiten sind für uns die durch sogenannte Transpositionen, d.~h. durch successive Vertauschung von nur zwei Elementen ausgeführten, von besonderem Interesse. Durch mehrere, nach einander ausgeführte Transpositionen lässt sich aus jeder Anordnung, z.~B. aus \(\mathfrak A\), jede andere \(\mathfrak A'\) herleiten. Man kann zu diesem Zweck etwa so verfahren, dass man in \(\mathfrak A\) zunächst das Element \(1\) mit dem, was in \(\mathfrak A'\) an erster Stelle steht, also mit \(\alpha_1\), vertauscht (falls nicht \(\alpha_1=1\) ist), dann, wenn \(\alpha_2\) nicht schon \(=2\) ist, \(2\) mit \(\alpha_2\) u.~s.~f.

Um z.~B. von \((1,2,3,4)\) zu \((4,3,2,1)\) zu gelangen, bildet man die Anordnungen
\[
(1,2,3,4),\qquad (4,2,3,1),\qquad (4,3,2,1).
\]
Bezeichnen wir eine Transposition kurz durch die beiden vertauschten Ziffern, also die Vertauschung von \(1\) mit \(2\) durch \((1,2)\), so haben wir hier nach einander die Transpositionen
\[
(1,4),\quad (2,3)
\]
ausgeführt.

Es ist zu bemerken, dass der Uebergang von einer Anordnung zu einer bestimmten anderen auf unendlich viele verschiedene Arten durch auf einander folgende Transpositionen erreicht werden kann. So geht die Anordnung \((1,2,3,4)\) auch durch die Transpositionen
\[
(1,2),\quad (1,3),\quad (2,4),\quad (1,2)
\]
in \((4,3,2,1)\) über.

\section*{\S\,19. Permutationen erster und zweiter Art.}

Die \(\Pi(n)\) Anordnungen von \(n\) Elementen lassen sich nach folgendem Gesichtspunkte in zwei Arten zerlegen.

Aus den \(n\) Elementen unseres Systems lassen sich
\[
{n(n-1)\over2}
\]
und nicht mehr Paare bilden. Wir wollen nun den \(n\) Elementen \(1,2,\ldots,n\) in bestimmter Weise \(n\) reelle Zahlwerthe \(a_1,a_2,\ldots,a_n\) zuordnen, und aus diesen Zahlwerthen die \(\frac{n(n-1)}2\) Differenzen
\[
a_1-a_2,\quad a_1-a_3,\ldots
\]
bilden, wobei der niedrigere Index dem Minuenden angehören soll. Das Differenzenproduct
\begin{equation}
P=(a_1-a_2)(a_1-a_3)\cdots(a_1-a_n)
(a_2-a_3)\cdots(a_2-a_n)\cdots(a_{n-1}-a_n)
\tag{1}
\end{equation}
wird, wenn die \(a_1,a_2,\ldots,a_n\) von einander verschiedene Zahlwerthe sind, einen von Null verschiedenen Werth haben, z.~B. einen positiven, wenn \(a_1>a_2>a_3>\cdots>a_n\) angenommen war.

Wenn wir nun die Indices \(1,2,3,\ldots,n\) irgendwie unter einander vertauschen, also etwa von \(\mathfrak A\) zu \(\mathfrak A'\) übergehen, so geht \(P\) in
\begin{equation}
\begin{aligned}
P'={}&(a_{\alpha_1}-a_{\alpha_2})(a_{\alpha_1}-a_{\alpha_3})\cdots
(a_{\alpha_1}-a_{\alpha_n})\\
&\cdot(a_{\alpha_2}-a_{\alpha_3})\cdots(a_{\alpha_2}-a_{\alpha_n})
\cdots(a_{\alpha_{n-1}}-a_{\alpha_n})
\end{aligned}
\tag{2}
\end{equation}
über, und dies Product besteht, abgesehen vom Vorzeichen, aus denselben Factoren wie \(P\), d.~h. es ist \(P'\) entweder gleich \(P\) oder entgegengesetzt zu \(P\).

I. Wir rechnen nun die Anordnung \(\mathfrak A'\) und also auch die Permutation, die \(\mathfrak A\) in \(\mathfrak A'\) verwandelt, zur ersten oder zur zweiten Art, je nachdem \(P\) mit \(P'\) gleich oder entgegengesetzt ist, so dass \(\mathfrak A\) selbst zur ersten Art gehört.

II. Durch eine einfache Transposition \((h,k)\), worin \(h,k\) irgend zwei der Ziffern \(1,2,\ldots,n\) bezeichnen, ändert sowohl \(P\) als \(P'\) sein Vorzeichen.

Denn die Factoren, die \(h\) und \(k\) gar nicht enthalten, werden durch diese Transposition nicht berührt; dann haben wir in \(P\) und \(P'\) den Factor \(\pm(a_h-a_k)\) und die Factorenpaare
\[
\pm(a_h-a_\nu)(a_k-a_\nu),
\]
wo \(\nu\) die Reihe der Zahlen \(1,2,\ldots,n\), mit Ausnahme von \(h,k\), durchläuft. Der erstere Factor ändert aber sein Zeichen, während das Factorenpaar ungeändert bleibt bei der Transposition \((h,k)\). Daraus folgt:

III. Die Permutationen der ersten Art sind aus einer geraden Anzahl von Transpositionen zusammengesetzt, und die der zweiten Art aus einer ungeraden Anzahl.

Zu der ersten Art ist dann auch die sogenannte identische Permutation zu rechnen, die \(\mathfrak A\) ungeändert lässt.

Daraus ergiebt sich noch die Folgerung: Auf wie verschiedenen Wegen man auch \(\mathfrak A'\) aus \(\mathfrak A\) durch Transpositionen ableiten mag, die Anzahl dieser Transpositionen ist bei allen diesen Arten übereinstimmend gerade oder ungerade (je nachdem \(\mathfrak A'\) zur ersten oder zur zweiten Art gehört).

Wenn wir in den sämmtlichen Anordnungen
\[
\mathfrak A,\mathfrak A',\mathfrak A'',\ldots
\]
der \(n\) Elemente eine Transposition, etwa \((1,2)\), vornehmen, so geht jede dieser Anordnungen in eine bestimmte andere über, etwa \(\mathfrak A\) in \(\mathfrak B\), \(\mathfrak A'\) in \(\mathfrak B'\), \(\mathfrak A''\) in \(\mathfrak B''\), und wenn wir dieselbe Transposition noch einmal wiederholen, so geht \(\mathfrak B\) wieder in \(\mathfrak A\), \(\mathfrak B'\) wieder in \(\mathfrak A'\) u.~s.~f. über. Daraus folgt, dass die Anordnungen \(\mathfrak B,\mathfrak B',\mathfrak B'',\ldots\) alle von einander verschieden sind und folglich in ihrer Gesammtheit mit der Gesammtheit der \(\mathfrak A\) übereinstimmen. Da nun, wie wir oben gesehen haben, die sämmtlichen Differenzenproducte
\[
P,P',P'',\ldots
\]
die aus \(P\) mit den verschiedenen Anordnungen \(\mathfrak A,\mathfrak A',\mathfrak A'',\ldots\) gebildet sind, durch eine Transposition das Zeichen ändern, so folgt, dass jedem \(\mathfrak A\) der ersten Art ein \(\mathfrak B\) der zweiten Art entspricht und jedem \(\mathfrak A\) der zweiten Art ein \(\mathfrak B\) der ersten Art.

IV. Hiernach ist die Anzahl der Anordnungen der ersten Art ebenso gross, wie die Anzahl der Anordnungen der zweiten Art, nämlich
\[
{1\over2}\Pi(n).
\]

Für \(n=3\) haben wir die folgenden sechs Anordnungen, von denen die erste Horizontalreihe die erste Art bildet:
\[
\begin{array}{ccc}
(1,2,3),&(2,3,1),&(3,1,2),\\
(3,2,1),&(2,1,3),&(1,3,2).
\end{array}
\]
Diese Sätze sind hier aus der Betrachtung des Productes \(P\), also einer Zahlgrösse, gewonnen. Wie man ohne Benutzung einer solchen Function zu denselben Ergebnissen gelangen kann, werden wir im XIV. Abschnitt sehen.

\section*{\S\,20. Determinanten.}

Wir betrachten jetzt ein System von \(n^2\) beliebigen Grössen, mit denen die rationalen Rechenoperationen ausgeführt werden können. Zu einer einfachen Bezeichnung dieser Grössen wählen wir einen Buchstaben mit einem doppelten Index \(a_i^{(k)}\), worin \(i\) sowohl als \(k\) die Reihe der Ziffern \(1,2,3,\ldots,n\) durchlaufen soll. Zur besseren Uebersicht ordnen wir diese Grössen in ein Quadrat, so dass alle \(a\) mit demselben oberen Index in einer Horizontalreihe, alle \(a\) mit demselben unteren Index in einer Verticalreihe stehen, und bezeichnen dies Quadrat mit \(\Delta\), also:
\begin{equation}
\Delta=
\begin{vmatrix}
a_1^{(1)}&a_2^{(1)}&a_3^{(1)}&\cdots&a_n^{(1)}\\
a_1^{(2)}&a_2^{(2)}&a_3^{(2)}&\cdots&a_n^{(2)}\\
a_1^{(3)}&a_2^{(3)}&a_3^{(3)}&\cdots&a_n^{(3)}\\
\vdots&\vdots&\vdots&&\vdots\\
a_1^{(n)}&a_2^{(n)}&a_3^{(n)}&\cdots&a_n^{(n)}
\end{vmatrix}.
\tag{1}
\end{equation}
Der Kürze halber nennt man die Horizontalreihen Zeilen, die Verticalreihen Colonnen.

Wir wollen aber unter dem zwischen verticalen Strichen eingeschlossenen Quadrat nicht nur den Complex der Grössen \(a\) verstehen, sondern eine bestimmte arithmetische Verbindung dieser Grössen, die sich ausrechnen lässt, sobald die \(a\) numerisch gegeben sind, und die wir jetzt beschreiben wollen.

Man bilde das Product der in der von links oben nach rechts unten gehenden Diagonale stehenden Glieder:
\begin{equation}
M=a_1^{(1)}a_2^{(2)}a_3^{(3)}\cdots a_n^{(n)};
\tag{2}
\end{equation}
leite daraus \(\Pi(n)\) Producte \(M,M',M'',\ldots\) her, indem man die unteren Indices permutirt, und gebe jedem so entstandenen Product das positive oder negative Zeichen, je nachdem die angewandte Permutation zur ersten oder zur zweiten Art gehört, also nach der Bezeichnung des vorigen Paragraphen
\begin{equation}
M'=\pm a_{\alpha_1}^{(1)}a_{\alpha_2}^{(2)}\cdots a_{\alpha_n}^{(n)}.
\tag{3}
\end{equation}

Die Summe aus diesen Producten soll \(\Delta\) sein. \(\Delta\) wird die Determinante der \(n^2\) Elemente \(a_i^{(k)}\) genannt, und zwar, wenn die Unterscheidung nothwendig ist, eine \(n\)-reihige Determinante (auch Determinante \(n\)ten Grades oder \(n\)ter Ordnung). Das Glied \(M\) dieser Summe, d.~h. also das Product aller in der Diagonale des Quadrats stehenden Elemente, wird das Hauptglied genannt.

Nehmen wir z.~B. \(n=2\), so erhalten wir
\begin{equation}
\Delta=
\begin{vmatrix}
a_1^{(1)}&a_2^{(1)}\\
a_1^{(2)}&a_2^{(2)}
\end{vmatrix}
=a_1^{(1)}a_2^{(2)}-a_2^{(1)}a_1^{(2)},
\tag{4}
\end{equation}
und für \(n=3\):
\begin{equation}
\begin{aligned}
\Delta={}&a_1^{(1)}a_2^{(2)}a_3^{(3)}
+a_2^{(1)}a_3^{(2)}a_1^{(3)}
+a_3^{(1)}a_1^{(2)}a_2^{(3)}\\
&-a_1^{(1)}a_3^{(2)}a_2^{(3)}
-a_2^{(1)}a_1^{(2)}a_3^{(3)}
-a_3^{(1)}a_2^{(2)}a_1^{(3)}.
\end{aligned}
\tag{5}
\end{equation}
Oder in anderer Bezeichnung:
\begin{equation}
\begin{vmatrix}
a&b&c\\
a'&b'&c'\\
a''&b''&c''
\end{vmatrix}
=ab'c''+bc'a''+ca'b''-ac'b''-ba'c''-cb'a''.
\tag{6}
\end{equation}

\section*{\S\,21. Hauptsätze über Determinanten.}

Aus dem Begriff der Determinante ergeben sich leicht die ersten Sätze, die für die Anwendung geeignet sind.

Wenn wir in dem Product
\begin{equation}
M'=\pm a_{\alpha_1}^{(1)}a_{\alpha_2}^{(2)}\cdots a_{\alpha_n}^{(n)}
\tag{1}
\end{equation}
die Factoren umstellen, so ändert sich sein Werth nicht. Wir können also die Factoren auch so anordnen, dass die unteren Indices in ihrer natürlichen Reihenfolge \(1,2,\ldots,n\) erscheinen. Dabei werden dann die oberen Indices in einer gewissen Weise permutirt erscheinen, also \(M'\) die Form erhalten:
\begin{equation}
\pm a_1^{(\beta_1)}a_2^{(\beta_2)}\cdots a_n^{(\beta_n)},
\tag{2}
\end{equation}
worin
\[
(\beta_1,\beta_2,\ldots,\beta_n)=\mathfrak B
\]
ebenso wie \((\alpha_1,\alpha_2,\ldots,\alpha_n)\) eine Anordnung der Ziffern \(1,2,\ldots,n\) bedeutet. Man kann die Anordnung \(\mathfrak B\) dadurch erhalten, dass man in den Factoren von \(M'\) die Transpositionen, die zu \(\mathfrak A\) geführt haben, von der letzten anfangend, rückgängig macht, um in der Reihe der unteren Indices wieder die ursprüngliche Anordnung zu erhalten. Die dabei sich ergebende Reihenfolge der oberen Indices ist dann die Anordnung \(\mathfrak B\). Es folgt daraus, dass \(\mathfrak B\) zur ersten oder zur zweiten Art gehört, je nachdem \(\mathfrak A\) zur ersten oder zur zweiten Art gehört, da beide durch die gleiche Anzahl von Transpositionen entstehen. Die Gesammtheit der \(\mathfrak B\) stellt ebenso wie die Gesammtheit der \(\mathfrak A\) alle Permutationen der \(n\) Elemente dar, da zwei verschiedene \(\mathfrak A\) niemals zu demselben \(\mathfrak B\) führen können. Damit ist bewiesen:

I. Die Determinante \(\Delta\) kann auch dadurch gebildet werden, dass man in dem Hauptglied \(a_1^{(1)}a_2^{(2)}\cdots a_n^{(n)}\) die oberen Indices auf alle möglichen Arten permutirt, jedem der so gebildeten Producte das positive oder negative Zeichen giebt, je nachdem die angewandte Permutation zur ersten oder zweiten Art gehört, und dann die Summe aller dieser Producte nimmt.

In der Darstellung von \(\Delta\) werden durch die oberen Indices die Zeilen, durch die unteren Indices die Colonnen gekennzeichnet, und demnach können wir diesem Satze auch den folgenden Ausdruck geben:

II. Eine Determinante ändert sich nicht, wenn die Zeilen zu Colonnen und die Colonnen zu Zeilen gemacht werden.

Wenn wir in den sämmtlichen Anordnungen \(\mathfrak A,\mathfrak A',\mathfrak A'',\ldots\) der \(n\) Elemente irgend zwei Elemente mit einander vertauschen, so bleibt die Gesammtheit dieser Anordnungen ungeändert, aber es geht jede Anordnung erster Art in eine Anordnung zweiter Art über und umgekehrt. Wenn wir also in den Gliedern \(M,M',M'',\ldots\), aus denen \(\Delta\) zusammengesetzt ist, irgend zwei untere Indices vertauschen, so geht jedes Glied mit positivem Zeichen in ein anderes über, das in \(\Delta\) mit dem negativen Zeichen behaftet war und umgekehrt, also es ändert \(\Delta\) sein Vorzeichen. Daraus folgt mit Hülfe von II. der Satz:

III. Wenn man in \(\Delta\) zwei untere oder zwei obere Indices mit einander vertauscht, so ändert die Determinante nur ihr Vorzeichen.

Etwas anders ausgedrückt: wenn man zwei Zeilen oder zwei Colonnen mit einander vertauscht, so ändert die Determinante nur ihr Vorzeichen; und daraus allgemeiner:

IV. Wenn in einer Determinante die Zeilen oder die Colonnen permutirt werden, so ändert sich der absolute Werth nicht, und das Vorzeichen ändert sich nicht oder geht in das entgegengesetzte über, je nachdem die angewandte Permutation zur ersten oder zweiten Art gehört.

Aus III. erhält man den folgenden Fundamentalsatz:

V. Wenn in zwei Zeilen oder in zwei Colonnen die an gleicher Stelle stehenden Glieder einander gleich sind (kürzer ausgedrückt: wenn zwei Reihen einander gleich sind), so hat die Determinante den Werth Null.

Denn die Vertauschung der zwei Reihen ändert nach III. das Zeichen, kann aber andererseits, da beide Reihen identisch sind, nichts ändern, so dass für \(\Delta\) nur der Werth Null übrig bleibt.

Man drückt den Satz V. nur anders aus, wenn man sagt:

VI. Man erhält eine verschwindende Determinante, wenn man die Elemente einer Reihe durch die entsprechenden Elemente einer anderen Reihe, oder, kurz gesagt, wenn man einen unteren oder oberen Index durch einen anderen ersetzt.

\section*{\S\,22. Unterdeterminanten.}

In jedem Gliede der entwickelten Determinante
\[
\sum \pm a_{\alpha_1}^{(1)}a_{\alpha_2}^{(2)}\cdots a_{\alpha_n}^{(n)},
\]
deren Werth wir jetzt mit \(A\) bezeichnen wollen, kommt jede der Zahlen \(1,2,\ldots,n\) ein und nur einmal als unterer Index vor. Es wird also ein gewisser Complex von Gliedern den Factor \(a_1^{(1)}\) enthalten, ein anderer Complex den Factor \(a_1^{(2)}\) u.~s.~f., endlich ein Complex den Factor \(a_1^{(n)}\); jedes Glied der Determinante kommt in einem und nur in einem dieser Complexe vor.

Bezeichnen wir also den ersten dieser Complexe mit \(a_1^{(1)}A_1^{(1)}\), den zweiten mit \(a_1^{(2)}A_1^{(2)}\), den letzten mit \(a_1^{(n)}A_1^{(n)}\), so können wir die Determinante folgendermaassen darstellen:
\begin{equation}
A=a_1^{(1)}A_1^{(1)}+a_1^{(2)}A_1^{(2)}+\cdots+a_1^{(n)}A_1^{(n)}.
\tag{1}
\end{equation}
An Stelle des unteren Index \(1\) hätten wir ebenso gut jeden anderen, \(r\), herausgreifen und daher
\begin{equation}
A=a_r^{(1)}A_r^{(1)}+a_r^{(2)}A_r^{(2)}+\cdots+a_r^{(n)}A_r^{(n)}
\tag{2}
\end{equation}
setzen können. Darin bedeutet das Product \(a_r^{(\nu)}A_r^{(\nu)}\) den Complex aller Glieder der Determinante, die den Factor \(a_r^{(\nu)}\) enthalten.

Da dieselben Regeln wie für die unteren so auch für die oberen Indices gelten, so kann man die Determinante auch noch in der folgenden Weise schreiben:
\begin{equation}
A=a_1^{(\mu)}A_1^{(\mu)}+a_2^{(\mu)}A_2^{(\mu)}+\cdots+a_n^{(\mu)}A_n^{(\mu)},
\tag{3}
\end{equation}
worin \(\mu\) gleichfalls jeden der Indices \(1,2,\ldots,n\) bedeuten kann.

Die hierdurch vollständig definirten Grössen \(A_r^{(\nu)}\) heissen die Unterdeterminanten der Determinante \(A\). Um ihre Bildungsweise genau kennen zu lernen, betrachten wir zunächst den Complex \(a_1^{(1)}A_1^{(1)}\). Man erhält ihn, wenn man in dem Product
\[
a_1^{(1)}a_2^{(2)}\cdots a_n^{(n)}
\]
den unteren Index \(1\) ungeändert lässt und nur die übrigen Indices \(2,3,\ldots,n\) auf alle Arten permutirt und die Summe der entstandenen Glieder mit Rücksicht auf die Zeichenregel bildet, d.~h. es ist \(A_1^{(1)}\) die \((n-1)\)-reihige Determinante:
\begin{equation}
A_1^{(1)}=
\begin{vmatrix}
a_2^{(2)}&a_3^{(2)}&\cdots&a_n^{(2)}\\
a_2^{(3)}&a_3^{(3)}&\cdots&a_n^{(3)}\\
\vdots&\vdots&&\vdots\\
a_2^{(n)}&a_3^{(n)}&\cdots&a_n^{(n)}
\end{vmatrix},
\tag{4}
\end{equation}
oder die Determinante, die man aus \(A\) erhält, wenn man in dem \(A\) darstellenden Quadrat die erste Zeile und die erste Colonne weglässt.

Daraus ergiebt sich leicht die Bedeutung von \(A_\mu^{(\nu)}\): man kann, indem man \(\nu-1\) Zeilenvertauschungen vornimmt, die \(\nu\)te Zeile zur ersten machen, und wenn man noch \(\mu-1\) Vertauschungen der Colonnen hinzunimmt, die \(\mu\)te Colonne zur ersten; im Uebrigen bleiben die Reihen in ihrer Aufeinanderfolge ungeändert. Die Determinante selbst hat den Factor \((-1)^{\mu+\nu}\) angenommen und ist dem absoluten Werthe nach ungeändert geblieben. In der so umgeänderten Reihenfolge ist aber das Element \(a_\mu^{(\nu)}\) an die Stelle des Elementes \(a_1^{(1)}\) getreten, und daraus schliesst man auf folgendes Bildungsgesetz:

Man erhält die Unterdeterminante \(A_\mu^{(\nu)}\) dadurch, dass man in dem die Determinante darstellenden Quadrat die beiden Reihen weglässt, die sich in \(a_\mu^{(\nu)}\) kreuzen, und den Factor \((-1)^{\mu+\nu}\) hinzufügt.

So erhält man z.~B. für die dreireihige Determinante die folgende Darstellung:
\begin{equation}
\begin{vmatrix}
a&b&c\\
a'&b'&c'\\
a''&b''&c''
\end{vmatrix}
=a(b'c''-c'b'')+b(c'a''-a'c'')+c(a'b''-b'a'').
\tag{5}
\end{equation}

Da der untere Index \(\nu\) in \(A_\mu^{(\nu)}\) gar nicht vorkommt, so ändert sich \(A_\mu^{(\nu)}\) nicht, wenn der untere Index \(\nu\) durch einen anderen ersetzt wird. Dann aber verschwindet nach \S~21, VI. die Determinante. Wir erhalten demnach aus (2) die folgende wichtige Relation, in der \(\mu,\nu\) irgend zwei von einander verschiedene Ziffern \(1,2,\ldots,n\) sein können:
\begin{equation}
0=a_1^{(\mu)}A_1^{(\nu)}+a_2^{(\mu)}A_2^{(\nu)}+\cdots+a_n^{(\mu)}A_n^{(\nu)},
\tag{6}
\end{equation}
und ebenso bekommt man aus (3):
\begin{equation}
0=a_\mu^{(1)}A_\nu^{(1)}+a_\mu^{(2)}A_\nu^{(2)}+\cdots+a_\mu^{(n)}A_\nu^{(n)}.
\tag{7}
\end{equation}

Beispielsweise ergiebt sich aus (5), wenn \(a,b,c\) durch \(a',b',c'\) ersetzt werden:
\begin{equation}
a'(b'c''-c'b'')+b'(c'a''-a'c'')+c'(a'b''-b'a'')=0,
\tag{8}
\end{equation}
eine Formel, von deren Richtigkeit man sich durch die einfachste Rechnung überzeugt.

Wenn wir die Relation (6) mit einem beliebigen Factor \(\lambda\) multipliciren und zu (2) addiren, so erhalten wir die Formel:
\begin{equation}
\begin{aligned}
A={}&(a_r^{(1)}+\lambda a_s^{(1)})A_r^{(1)}
+(a_r^{(2)}+\lambda a_s^{(2)})A_r^{(2)}\\
&+\cdots+(a_r^{(n)}+\lambda a_s^{(n)})A_r^{(n)},
\end{aligned}
\tag{9}
\end{equation}
die uns den folgenden Satz ausdrückt:

VII. Die Determinante ändert ihren Werth nicht, wenn man zu den Elementen einer Zeile, die mit einem beliebigen gemeinschaftlichen Factor multiplicirten entsprechenden Elemente einer anderen Zeile addirt.

Derselbe Satz gilt auch von den Colonnen. Er wird zur Vereinfachung und numerischen Berechnung von Determinanten oft mit Nutzen verwendet. Wir fügen noch folgende Sätze bei, die sich aus den Darstellungen (2), (3) sofort ablesen lassen.

VIII. Wenn alle Elemente einer Zeile oder einer Colonne einen gemeinschaftlichen Factor haben, so kann dieser weggelassen und als Factor vor die Determinante gesetzt werden.

Denn es ist nach (2):
\[
p a_r^{(1)}A_r^{(1)}+p a_r^{(2)}A_r^{(2)}+\cdots+p a_r^{(n)}A_r^{(n)}=pA.
\]

IX. Wenn in einer Zeile oder in einer Colonne alle Elemente bis auf eines verschwinden, so reducirt sich die Determinante auf das Product dieses einen Elementes mit der entsprechenden Unterdeterminante.

Denn wenn \(a_r^{(1)},a_r^{(2)},\ldots,a_r^{(n)}\) mit Ausnahme von \(a_r^{(\nu)}\) verschwinden, so ist nach (2):
\[
A=a_r^{(\nu)}A_r^{(\nu)}.
\]
Der Werth von \(A\) ist dann von den \(a_1^{(\nu)},a_2^{(\nu)},\ldots,a_n^{(\nu)}\) (mit Ausnahme von \(a_r^{(\nu)}\)) ganz unabhängig.

Um von diesen Sätzen eine Anwendung zu machen, wollen wir den Werth der Determinante
\[
\Delta=
\begin{vmatrix}
1&a&a^2\\
1&b&b^2\\
1&c&c^2
\end{vmatrix}
\]
bestimmen, worin \(a,b,c\) beliebige Grössen seien. Multipliciren wir die zweite Colonne mit \(a\) und subtrahiren sie von der dritten, darauf die erste mit \(a\) und subtrahiren sie von der zweiten, so folgt nach VII:
\[
\Delta=
\begin{vmatrix}
1&0&0\\
1&b-a&b(b-a)\\
1&c-a&c(c-a)
\end{vmatrix},
\]
und nach IX. und VIII.:
\begin{equation}
\Delta=(b-a)(c-a)
\begin{vmatrix}
1&b\\
1&c
\end{vmatrix}
=(b-a)(c-a)(c-b).
\tag{10}
\end{equation}

Auf die gleiche Weise kann man auch die \(n\)-reihige Determinante
\[
\begin{vmatrix}
1&a_1&a_1^2&\cdots&a_1^{n-1}\\
1&a_2&a_2^2&\cdots&a_2^{n-1}\\
\vdots&\vdots&\vdots&&\vdots\\
1&a_n&a_n^2&\cdots&a_n^{n-1}
\end{vmatrix}
\]
behandeln und findet ihren Werth gleich
\begin{equation}
(a_2-a_1)(a_3-a_1)\cdots(a_n-a_1)
(a_3-a_2)\cdots(a_n-a_2)\cdots(a_n-a_{n-1}).
\tag{11}
\end{equation}

Ordnet man die Colonnen in umgekehrter Reihenfolge, so sind dazu, je nachdem \(n\) gerade oder ungerade ist,
\[
{n\over2}(n-1)\quad \hbox{oder}\quad {n-1\over2}n
\]
Vertauschungen erforderlich, so dass sich die so geordnete Determinante von \(\Delta\) durch den Factor
\[
(-1)^{n(n-1)/2}
\]
unterscheidet. Es kommt auf dasselbe hinaus, wenn man den \(\frac{n(n-1)}2\) Factoren des Productes (11) das entgegengesetzte Vorzeichen giebt. Es besteht also zugleich mit (11) die Gleichung:
\begin{equation}
\begin{vmatrix}
a_1^{\,n-1}&a_1^{\,n-2}&\cdots&a_1&1\\
a_2^{\,n-1}&a_2^{\,n-2}&\cdots&a_2&1\\
\vdots&\vdots&&\vdots&\vdots\\
a_n^{\,n-1}&a_n^{\,n-2}&\cdots&a_n&1
\end{vmatrix}
=
\prod_{1\le i<k\le n}(a_i-a_k).
\tag{12}
\end{equation}

Wir wollen hier noch eine Bezeichnungsweise der Unterdeterminanten erwähnen, die der Differentialrechnung entnommen ist und oft mit Nutzen verwendet wird, besonders wenn es sich um die Bildung von Derivirten handelt. Wenn die Grössen \(a_\mu^{(\nu)}\) als unabhängige Variable betrachtet werden, so ist die Determinante \(A\) in Bezug auf jede von ihnen nur vom ersten Grade. Die nach \(a_\mu^{(\nu)}\) genommene Ableitung oder der Differentialquotient ist also gleich dem Coefficienten von \(a_\mu^{(\nu)}\), also:
\[
{\partial A\over \partial a_\mu^{(\nu)}}=A_\mu^{(\nu)}.
\]
Wenn demnach z.~B. die \(a_\mu^{(\nu)}\) Functionen einer Variablen \(t\) sind, so ist auch \(A\) eine Function von \(t\), und man erhält nach den ersten Regeln der Differentialrechnung die Ableitung von \(A\) in Bezug auf \(t\) in der Form
\[
A'(t)=\sum A_\mu^{(\nu)}{\partial a_\mu^{(\nu)}\over\partial t},
\]
wenn \({\partial a_\mu^{(\nu)}}/{\partial t}\) die Ableitung von \(a_\mu^{(\nu)}\) nach \(t\) bedeutet.


\section*{\S\,23. Die Unterdeterminanten im weiteren Sinne}

Wir können nun die Betrachtungen des vorigen Paragraphen in folgender Weise verallgemeinern.  Wie wir vorhin von der Aufgabe ausgegangen sind, alle Glieder in der entwickelten Determinante \(A\) aufzusuchen, die den Faktor \(a_1^{(1)}\) enthalten, so wollen wir jetzt alle die Glieder aufsuchen, die den Faktor
\[
        a_1^{(1)}a_2^{(2)}\cdots a_\nu^{(\nu)}
\]
enthalten, worin \(\nu\) eine beliebige Zahl unter \(n\) sein kann.  Diese Glieder erhalten wir aus dem Hauptgliede
\[
        a_1^{(1)}a_2^{(2)}\cdots a_\nu^{(\nu)}a_{\nu+1}^{(\nu+1)}\cdots a_n^{(n)},
\]
wenn wir bei der Permutation der unteren Indices \(1,2,\ldots,\nu\) unverändert lassen und nur \(\nu+1,\ldots,n\) auf alle Arten permutieren, unter Berücksichtigung der Vorzeichenregel.  Demnach ist der Inbegriff der gesuchten Glieder
\begin{equation}
        a_1^{(1)}a_2^{(2)}\cdots a_\nu^{(\nu)}
        \begin{vmatrix}
        a_{\nu+1}^{(\nu+1)}&\cdots&a_n^{(\nu+1)}\\
        \vdots&&\vdots\\
        a_{\nu+1}^{(n)}&\cdots&a_n^{(n)}
        \end{vmatrix}.
\tag{1}
\end{equation}

I. Die hier als Faktor auftretende Determinante von \(n-\nu\) Reihen, die wir mit
\[
        A_{1,2,\ldots,\nu}^{1,2,\ldots,\nu}
\]
bezeichnen, entsteht aus \(A\) durch Weglassen der \(\nu\) ersten Zeilen und Kolonnen.

Dieses Resultat wollen wir nun verallgemeinern.  Wir wählen irgend \(\nu\) Elemente
\begin{equation}
        a_{\alpha_1}^{(\beta_1)},
        a_{\alpha_2}^{(\beta_2)},\ldots,
        a_{\alpha_\nu}^{(\beta_\nu)},
\tag{3}
\end{equation}
jedoch so, daß nicht zwei Elemente in derselben Zeile oder in derselben Kolonne vorkommen.  Den Inbegriff der Glieder der Determinante, die das Produkt dieser Elemente als Faktor enthalten, bezeichnen wir mit
\begin{equation}
        A_{\alpha_1,\alpha_2,\ldots,\alpha_\nu}^{\beta_1,\beta_2,\ldots,\beta_\nu}.
\tag{2}
\end{equation}

Man kann durch Umstellen von Zeilen und Kolonnen, wodurch höchstens das Zeichen der Determinante geändert wird, immer erreichen, daß die Elemente (3) an die Stelle von \(a_1^{(1)},a_2^{(2)},\ldots,a_\nu^{(\nu)}\) gelangen.  Dann läßt sich Regel I anwenden und es ergibt sich:

II. Man erhält, vom Vorzeichen abgesehen, \(A_{\alpha_1,\ldots,\alpha_\nu}^{\beta_1,\ldots,\beta_\nu}\) als \((n-\nu)\)-reihige Determinante, wenn man in \(A\) alle Zeilen und Kolonnen wegläßt, die sich in einem der Elemente (3) schneiden, und die übrig bleibenden Zeilen und Kolonnen in ihrer Reihenfolge stehen läßt.

Für die Zeichenbestimmung ordne man die unteren und die oberen Indices \(1,2,
\ldots,n\) in der Weise
\begin{equation}
        \alpha_1,\alpha_2,\ldots,\alpha_\nu,\alpha_{\nu+1},\ldots,\alpha_n,
\tag{4}
\end{equation}
\begin{equation}
        \beta_1,\beta_2,\ldots,\beta_\nu,\beta_{\nu+1},\ldots,\beta_n,
\tag{5}
\end{equation}
indem man die fehlenden Indices der Größe nach aufeinander folgen läßt.

III. Die in II beschriebene \((n-\nu)\)-reihige Determinante erhält das positive oder negative Zeichen, je nachdem die beiden Anordnungen (4), (5) der Ziffern \(1,2,\ldots,n\) beide zu derselben oder zu verschiedenen Arten gehören.

Die so definierten Größen heißen die \(\nu\)-ten Unterdeterminanten oder Unterdeterminanten \(\nu\)-ter Ordnung.  Sie sind dargestellt durch \((n-\nu)\)-reihige Determinanten.  Es folgt:

IV. Die Unterdeterminante \(A_{\alpha_1,
\ldots,
\alpha_\nu}^{\beta_1,
\ldots,
\beta_\nu}\) ändert nur ihr Vorzeichen, wenn zwei ihrer unteren oder zwei ihrer oberen Indices vertauscht werden; dem absoluten Werte nach bleibt sie ungeändert.

Bezeichnen wir mit einer überstrichenen Folge irgend eine Anordnung der Indices, so enthält die Determinante auch die Glieder
\[
        \pm a_{\bar\alpha_1}^{(\beta_1)}a_{\bar\alpha_2}^{(\beta_2)}\cdots
        a_{\bar\alpha_\nu}^{(\beta_\nu)}
        A_{\bar\alpha_1,\ldots,\bar\alpha_\nu}^{\beta_1,\ldots,\beta_\nu}.
\]
Die hierbei auftretende \(\nu\)-reihige Determinante wird die complementäre Unterdeterminante genannt und mit
\[
        B_{\alpha_1,\ldots,\alpha_\nu}^{\beta_1,\ldots,\beta_\nu}
\]
bezeichnet.  Der betreffende Complex der Glieder ist also
\begin{equation}
        A_{\alpha_1,\ldots,\alpha_\nu}^{\beta_1,\ldots,\beta_\nu}
        B_{\alpha_1,\ldots,\alpha_\nu}^{\beta_1,\ldots,\beta_\nu}.
\tag{7}
\end{equation}
Daraus folgt der Satz von Laplace:
\begin{equation}
        A=\sum
        A_{\alpha_1,\ldots,\alpha_\nu}^{\beta_1,\ldots,\beta_\nu}
        B_{\alpha_1,\ldots,\alpha_\nu}^{\beta_1,\ldots,\beta_\nu}.
\tag{8}
\end{equation}

Eine andere Darstellung von \(A\) durch erste und zweite Unterdeterminanten ist
\begin{equation}
        A=a_\nu^{(\mu)}A_\nu^{(\mu)}+
        \sum a_\nu^{(i)}a_k^{(\mu)}A_{\nu k}^{\mu i},
\tag{9}
\end{equation}
oder, nach der Vorzeichenregel,
\begin{equation}
        A=a_\nu^{(\mu)}A_\nu^{(\mu)}-
        \sum a_\nu^{(i)}a_k^{(\mu)}A_{\nu k}^{\mu i}.
\tag{10}
\end{equation}
Für die geränderte Determinante
\begin{equation}
        U=\begin{vmatrix}
        a_1^{(1)}&a_2^{(1)}&\cdots&a_n^{(1)}&u_1\\
        a_1^{(2)}&a_2^{(2)}&\cdots&a_n^{(2)}&u_2\\
        \vdots&\vdots&&\vdots&\vdots\\
        a_1^{(n)}&a_2^{(n)}&\cdots&a_n^{(n)}&u_n\\
        v_1&v_2&\cdots&v_n&q
        \end{vmatrix}
\tag{11}
\end{equation}
ergibt sich
\begin{equation}
        U=qA-\sum u_i v_k A_k^{(i)}.
\tag{12}
\end{equation}
Auch bei höheren Unterdeterminanten ist die Bezeichnung durch Differentialquotienten zweckmäßig, z.B.
\begin{equation}
        A_{\nu k}^{\mu i}=
        {\partial^2 A\over \partial a_\nu^{(\mu)}\,\partial a_k^{(i)}} .
\tag{13}
\end{equation}

\section*{\S\,24. Lineare homogene Gleichungen}

Die hauptsächlichste Anwendung der Determinanten, der die ganze Theorie ihren Ursprung verdankt, ist die Auflösung linearer Gleichungen.  Wir betrachten ein System von \(m\) Gleichungen ersten Grades, in denen \(n\) Unbekannte \(x_1,x_2,
\ldots,x_n\) homogen vorkommen:
\begin{equation}
\begin{aligned}
 a_1^{(1)}x_1+a_2^{(1)}x_2+\cdots+a_n^{(1)}x_n&=0,\\
 a_1^{(2)}x_1+a_2^{(2)}x_2+\cdots+a_n^{(2)}x_n&=0,\\
 &\ \vdots\\
 a_1^{(m)}x_1+a_2^{(m)}x_2+\cdots+a_n^{(m)}x_n&=0.
\end{aligned}
\tag{1}
\end{equation}
Eine Lösung ist stets
\begin{equation}
        x_1=0,\ x_2=0,\ldots,x_n=0.
\tag{2}
\end{equation}

Das System der Coefficienten schreiben wir als Matrix
\begin{equation}
\begin{array}{cccc}
 a_1^{(1)}&a_2^{(1)}&\cdots&a_n^{(1)}\\
 a_1^{(2)}&a_2^{(2)}&\cdots&a_n^{(2)}\\
 \vdots&\vdots&&\vdots\\
 a_1^{(m)}&a_2^{(m)}&\cdots&a_n^{(m)}.
\end{array}
\tag{3}
\end{equation}
Unter den aus der Matrix hervorgehenden Determinanten sei wenigstens eine \(\nu\)-reihige von Null verschieden, während alle \((\nu+1)\)-reihigen verschwinden.  Ohne Beschränkung der Allgemeinheit sei
\begin{equation}
        A=\begin{vmatrix}
        a_1^{(1)}&a_2^{(1)}&\cdots&a_\nu^{(1)}\\
        \vdots&\vdots&&\vdots\\
        a_1^{(\nu)}&a_2^{(\nu)}&\cdots&a_\nu^{(\nu)}
        \end{vmatrix}\ne0.
\tag{4}
\end{equation}

I. Ist zunächst \(\nu=n\), so haben die Gleichungen keine andere Lösung als (2).  Insbesondere gilt: Wenn ein System von \(n\) linearen homogenen Gleichungen mit \(n\) Unbekannten eine von Null verschiedene Determinante hat, so sind sämtliche Unbekannte Null; hat es eine Lösung, bei der nicht alle Unbekannten verschwinden, so verschwindet die Determinante des Systems.

Ist \(\nu<n\), so werden die ersten \(\nu\) Gleichungen in die Form
\begin{equation}
\begin{aligned}
 a_1^{(1)}x_1+\cdots+a_\nu^{(1)}x_\nu&=-a_{\nu+1}^{(1)}x_{\nu+1}-\cdots-a_n^{(1)}x_n,\\
 &\vdots\\
 a_1^{(\nu)}x_1+\cdots+a_\nu^{(\nu)}x_\nu&=-a_{\nu+1}^{(\nu)}x_{\nu+1}-\cdots-a_n^{(\nu)}x_n
\end{aligned}
\tag{6}
\end{equation}
geschrieben.  Dann folgt
\begin{equation}
        A x_\mu=-\sum_{h=\nu+1}^{n} C_{\mu h}x_h,
\tag{7}
\end{equation}
wo
\begin{equation}
        C_{\mu h}=\sum_{i=1}^{\nu}a_h^{(i)}A_\mu^{(i)}.
\tag{8}
\end{equation}
Hierdurch sind \(x_1,\ldots,x_\nu\) linear durch \(x_{\nu+1},\ldots,x_n\) ausgedrückt.

III. Durch die Ausdrücke (7) sind die Gleichungen (1) befriedigt, welche Werte auch \(x_{\nu+1},\ldots,x_n\) haben mögen.  Also bleiben \(n-\nu\) Unbekannte willkürlich.

Für \(m=n-1\), \(\nu=n-1\), sind die Verhältnisse der Unbekannten völlig bestimmt:
\begin{equation}
        x_1:x_2:\cdots:x_n=A_1:A_2:\cdots:A_n.
\tag{16}
\end{equation}
Für \(n=3\) erhält man aus
\begin{equation}
        ax+by+cz=0,
        \qquad a'x+b'y+c'z=0
\tag{17}
\end{equation}
die Lösung
\begin{equation}
        x:y:z=bc'-cb':ca'-ac':ab'-ba'.
\tag{18}
\end{equation}

\section*{\S\,25. Elimination aus linearen Gleichungen}

Es kommt bisweilen vor, daß es sich bei einem gegebenen System linearer Gleichungen nicht um die wirkliche Ermittelung der Unbekannten handelt, sondern um die Beurteilung der Möglichkeit ihrer Lösung.  Die Aufstellung der Bedingungsgleichungen heißt Elimination.

Für ein System von \(m\) linearen homogenen Gleichungen mit \(n\) Unbekannten ist, wenn \(n\le m\), die notwendige und hinreichende Bedingung für eine nicht verschwindende Lösung die, daß alle \(n\)-reihigen Determinanten der Matrix verschwinden.  Die Anzahl dieser Determinanten ist
\[
        {m(m-1)\cdots(m-n+1)\over 1\cdot2\cdots n}.
\]
Um unabhängige Bedingungen zu erhalten, fragen wir danach, wann \(\nu\) Unbekannte durch die übrigen \(n-\nu\) bestimmt werden.  Ist
\begin{equation}
        A=\begin{vmatrix}
        a_1^{(1)}&a_2^{(1)}&\cdots&a_\nu^{(1)}\\
        \vdots&\vdots&&\vdots\\
        a_1^{(\nu)}&a_2^{(\nu)}&\cdots&a_\nu^{(\nu)}
        \end{vmatrix}\ne0,
\tag{1}
\end{equation}
so lauten die Bedingungen
\begin{equation}
\begin{vmatrix}
        a_1^{(1)}&\cdots&a_\nu^{(1)}&a_h^{(1)}\\
        \vdots&&\vdots&\vdots\\
        a_1^{(\nu)}&\cdots&a_\nu^{(\nu)}&a_h^{(\nu)}\\
        a_1^{(k)}&\cdots&a_\nu^{(k)}&a_h^{(k)}
\end{vmatrix}=0,
\tag{2}
\end{equation}
oder
\begin{equation}
        A a_h^{(k)}-\sum_{i=1}^{\nu}\sum_{\mu=1}^{\nu}
        a_h^{(i)}a_\mu^{(k)}A_\mu^{(i)}=0,
\tag{3}
\end{equation}
für \(h=\nu+1,\ldots,n\), \(k=\nu+1,\ldots,m\).  Ihre Anzahl ist
\begin{equation}
        (n-\nu)(m-\nu),
\tag{5}
\end{equation}
und sie sind voneinander unabhängig.

\section*{\S\,26. Unhomogene lineare Gleichungen}

Wir betrachten ein System
\begin{equation}
\begin{aligned}
 a_1^{(1)}x_1+a_2^{(1)}x_2+\cdots+a_n^{(1)}x_n&=y_1,\\
 &\vdots\\
 a_1^{(n)}x_1+a_2^{(n)}x_2+\cdots+a_n^{(n)}x_n&=y_n.
\end{aligned}
\tag{1}
\end{equation}
Sei
\begin{equation}
        A=\sum\pm a_1^{(1)}a_2^{(2)}\cdots a_n^{(n)}
\tag{2}
\end{equation}
die Determinante des Systems.  Dann ist
\begin{equation}
        A x_\mu=\sum_{k=1}^{n}y_k A_\mu^{(k)}.
\tag{4}
\end{equation}
Ist \(A\ne0\), so sind die Unbekannten eindeutig bestimmt.  Setzt man \(A_\mu^{(k)}=A_{k\mu}\), so erhält man
\begin{equation}
\begin{aligned}
 A x_1&=A_{11}y_1+A_{12}y_2+\cdots+A_{1n}y_n,\\
 &\vdots\\
 A x_n&=A_{n1}y_1+A_{n2}y_2+\cdots+A_{nn}y_n.
\end{aligned}
\tag{5}
\end{equation}

Wenn \(A=0\), ergeben sich Bedingungen für die \(y\).  Sei
\begin{equation}
        B=\begin{vmatrix}
        a_1^{(1)}&a_2^{(1)}&\cdots&a_\nu^{(1)}\\
        \vdots&\vdots&&\vdots\\
        a_1^{(\nu)}&a_2^{(\nu)}&\cdots&a_\nu^{(\nu)}
        \end{vmatrix}\ne0,
\tag{6}
\end{equation}
und seien alle \((\nu+1)\)-reihigen Unterdeterminanten Null.  Dann ist
\begin{equation}
        Bx_\mu=\sum_{i=1}^{\nu}B_\mu^{(i)}y_i
        -\sum_{h=\nu+1}^{n}x_h\sum_{i=1}^{\nu}B_\mu^{(i)}a_h^{(i)}.
\tag{7}
\end{equation}
Die Lösbarkeitsbedingungen sind
\begin{equation}
        B y_\lambda=\sum_{i=1}^{\nu}\sum_{\mu=1}^{\nu}
        B_\mu^{(i)}a_\mu^{(\lambda)}y_i,
        \qquad \lambda=\nu+1,\ldots,n.
\tag{10}
\end{equation}
Ist eine dieser Bedingungen nicht erfüllt, so hat das System keine Lösung.  Wird (1) als lineare Substitution aufgefaßt, so ist die Determinante \(A\) die Substitutionsdeterminante; nur für \(A\ne0\) sind die \(y\) unabhängige Variable und (5) gibt die inverse Substitution.


\section*{§ 27. Multiplication von Determinanten.}

Der Satz, den wir jetzt noch beweisen wollen, lehrt, wie man das Product zweier Determinanten von gleich viel Reihen durch eine einzige Determinante von ebenso viel Reihen darstellen kann. Man wird am einfachsten darauf geführt, wenn man die Auflösung von zwei Systemen linearer Gleichungen betrachtet.

Es seien jetzt die Coefficienten $a_i^{(k)}, b_h^{(k)}$, wenn $i,h$ die Reihe der Zahlen $1,2,\ldots,n$ durchlaufen, beliebige veränderliche Grössen, ebenso die Grössen $x_i,y_i,z_i$, die nur an die Relationen gebunden sind:
\begin{align}
 \sum_i a_i^{(k)}x_i&=y_k,\tag{1}\\
 \sum_k b_h^{(k)}y_k&=z_h.\tag{2}
\end{align}
Wenn nun die Aufgabe gestellt wird, die Variablen $x$ durch die Variablen $z$ zu bestimmen, so kann diese Aufgabe auf doppelte Art gelöst werden. Man kann nach § 26, (4) die Gleichungen (1) in Bezug auf $x$, die Gleichungen (2) in Bezug auf $y$ auflösen, und die letzteren Ausdrücke in die ersteren einsetzen. Bezeichnen wir mit $A,B$ die Determinanten der beiden Gleichungssysteme, also
\[
A=\begin{vmatrix}
 a_1^{(1)}&a_2^{(1)}&\cdots&a_n^{(1)}\\
 a_1^{(2)}&a_2^{(2)}&\cdots&a_n^{(2)}\\
 \cdots&\cdots&&\cdots\\
 a_1^{(n)}&a_2^{(n)}&\cdots&a_n^{(n)}
\end{vmatrix},\qquad
B=\begin{vmatrix}
 b_1^{(1)}&b_2^{(1)}&\cdots&b_n^{(1)}\\
 b_1^{(2)}&b_2^{(2)}&\cdots&b_n^{(2)}\\
 \cdots&\cdots&&\cdots\\
 b_1^{(n)}&b_2^{(n)}&\cdots&b_n^{(n)}
\end{vmatrix},
\]
und mit $A_i^{(k)},B_i^{(k)}$ die ersten Unterdeterminanten, so erhält man
\begin{align}
 A x_k&=\sum_i A_k^{(i)}y_i,\tag{3}\\
 B y_i&=\sum_h B_i^{(h)}z_h,\tag{4}
\end{align}
und durch Substitution von (4) in (3):
\begin{equation}
 ABx_k=\sum_h z_h\sum_i A_k^{(i)}B_i^{(h)}.\tag{5}
\end{equation}
Man kann aber auch so verfahren, dass man aus (1) die Ausdrücke für $y$ in (2) einsetzt, wodurch man, wenn
\begin{equation}
 c_i^{(h)}=\sum_s a_i^{(s)}b_h^{(s)}\tag{6}
\end{equation}
gesetzt wird, erhält:
\begin{equation}
 \sum_i c_i^{(h)}x_i=z_h.\tag{7}
\end{equation}
Wenn man nun
\[
C=\begin{vmatrix}
 c_1^{(1)}&c_2^{(1)}&\cdots&c_n^{(1)}\\
 c_1^{(2)}&c_2^{(2)}&\cdots&c_n^{(2)}\\
 \cdots&\cdots&&\cdots\\
 c_1^{(n)}&c_2^{(n)}&\cdots&c_n^{(n)}
\end{vmatrix}\tag{8}
\]
setzt, und mit $C_i^{(k)}$ die Unterdeterminanten von $C$ bezeichnet, so ergiebt die Auflösung von (7):
\begin{equation}
 Cx_k=\sum_h z_h C_k^{(h)},\tag{9}
\end{equation}
und die Vergleichung von (5) mit (9) ergiebt:
\begin{equation}
 AB=C,\tag{10}
\end{equation}
und für die Unterdeterminanten:
\begin{equation}
 C_k^{(h)}=\sum_i A_k^{(i)}B_i^{(h)}.\tag{11}
\end{equation}
In dieser Schlussweise ist aber noch eine Lücke, bedingt durch den Zweifel, ob sich (5) von (9) nicht durch einen beiden Seiten gemeinschaftlichen Factor unterscheiden könnte. Um diesem Zweifel zu begegnen, wollen wir die Gleichung (10) direct beweisen, woraus dann die Richtigkeit von (11) folgt.

Wir wollen das Hauptglied der Determinante $C$ nach (6) bilden, indem wir mit $s_1,s_2,\ldots,s_n$ von einander unabhängige Summationsbuchstaben bezeichnen, die von $1$ bis $n$ laufen:
\begin{align}
 c_1^{(1)}c_2^{(2)}\cdots c_n^{(n)}
 &=\sum_{s_1}a_1^{(s_1)}b_1^{(s_1)}
   \sum_{s_2}a_2^{(s_2)}b_2^{(s_2)}\cdots
   \sum_{s_n}a_n^{(s_n)}b_n^{(s_n)}\notag\\
 &=\sum_{s_1,s_2,\ldots,s_n}a_1^{(s_1)}a_2^{(s_2)}\cdots a_n^{(s_n)}
 b_1^{(s_1)}b_2^{(s_2)}\cdots b_n^{(s_n)}.\tag{12}
\end{align}
Die Permutation der unteren Indices der $c$ entspricht der Permutation der unteren Indices der $a$, und wenn man also mit Rücksicht auf die Vorzeichenregel die Determinante $C$ bildet, so erhält man
\begin{align}
 \sum \pm c_1^{(1)}c_2^{(2)}\cdots c_n^{(n)}
 =\sum_{s_1,s_2,\ldots,s_n} b_1^{(s_1)}b_2^{(s_2)}\cdots b_n^{(s_n)}
 \sum\pm a_1^{(s_1)}a_2^{(s_2)}\cdots a_n^{(s_n)}.\tag{13}
\end{align}
Nun ist
\[
\sum\pm a_1^{(s_1)}a_2^{(s_2)}\cdots a_n^{(s_n)}
\]
nach § 21, VI. immer dann gleich Null, wenn unter den $s_1,s_2,\ldots,s_n$ zweimal dieselbe Ziffer vorkommt; es behalten also nur die Glieder in (13) einen von Null verschiedenen Werth, in denen $s_1,s_2,\ldots,s_n$ eine Anordnung der Indices $1,2,\ldots,n$ ist, und zwar ist dieser Werth $+A$ oder $-A$, je nachdem diese Anordnung zur ersten oder zur zweiten Art gehört (§ 21, IV).

Demnach wird die rechte Seite von (13):
\[
 A\sum_{s_1,s_2,\ldots,s_n}\pm b_1^{(s_1)}b_2^{(s_2)}\cdots b_n^{(s_n)},
\]
die Summe erstreckt über alle Permutationen $s_1,s_2,\ldots,s_n$. Diese Summe ist aber gerade die Determinante $B$ und daher die Formel
\[
 C=AB
\]
bewiesen.

Das in der Formel (6) enthaltene Bildungsgesetz der Elemente $c_i^{(h)}$ können wir in Worten so ausdrücken: Um die Elemente der Determinante $C$, die das Product der beiden Determinanten $A,B$ ist, zu erhalten, multiplicirt man die Elemente je einer Colonne von $A$ mit den entsprechenden Elementen einer Colonne von $B$ und addirt die Producte.

Nach den Sätzen über die Determinanten kann man die Form von $C$ in mannigfacher Weise abändern. Wir wollen darüber Folgendes bemerken. Wenn man zwei Colonnen in $A$ oder in $B$ vertauscht, so ändern sich die Elemente $c_i^{(h)}$ nicht, sondern vertauschen sich nur unter einander. Wenn man aber zwei Zeilen in $A$ oder in $B$ vertauscht, so ändern sich die $c_i^{(h)}$, indem die Factoren in den einzelnen Producten der Summe anders zusammengefasst werden; wenn man aber entsprechende Vertauschungen in den Zeilen von $A$ und von $B$ gleichzeitig vornimmt, so bleiben die $c_i^{(h)}$ ungeändert, weil dadurch nur die einzelnen Glieder der Summe vertauscht werden.

Indem man in $A$ oder in $B$ oder in beiden zugleich die Zeilen zu Colonnen macht, erhält man noch drei verschiedene Arten für die Bildung des Products zweier Determinanten in Determinantenform. Letzteres kann man auch so ausdrücken: Um das Product zweier Determinanten zu bilden, kann man die Elemente der einzelnen Zeilen oder Colonnen des einen Factors mit den entsprechenden Elementen der Zeilen oder der Colonnen des anderen Factors multipliciren und die Producte addiren, und diese Productsummen als Elemente einer neuen Determinante auffassen.

Auf ein Product zweier Determinanten mit verschiedener Elementenzahl lässt sich die Multiplicationsregel dadurch anwenden, dass man die Determinante mit geringerer Reihenzahl durch den Satz § 22, IX. in eine andere mit mehr Reihen verwandelt.

\section*{§ 28. Determinanten der Unterdeterminanten.}

Wir machen hier gleich eine Anwendung von dem Multiplicationsgesetz der Determinanten. Es sei wie bisher
\[
A=\begin{vmatrix}
 a_1^{(1)}&a_2^{(1)}&\cdots&a_n^{(1)}\\
 a_1^{(2)}&a_2^{(2)}&\cdots&a_n^{(2)}\\
 \cdots&\cdots&&\cdots\\
 a_1^{(n)}&a_2^{(n)}&\cdots&a_n^{(n)}
\end{vmatrix},\tag{1}
\]
und
\[
\begin{matrix}
 A_1^{(1)}&A_2^{(1)}&\cdots&A_n^{(1)}\\
 A_1^{(2)}&A_2^{(2)}&\cdots&A_n^{(2)}\\
 \cdots&\cdots&&\cdots\\
 A_1^{(n)}&A_2^{(n)}&\cdots&A_n^{(n)}
\end{matrix}\tag{2}
\]
das System der Unterdeterminanten. Bilden wir aus (2) die Determinante, die wir mit $\Delta$ bezeichnen wollen, so können wir auf das Product $A\Delta$ die Multiplicationsregel anwenden. Dies giebt aber nach § 22, (3) und (7):
\[
 A\Delta=\begin{vmatrix}
 A&0&\cdots&0\\
 0&A&\cdots&0\\
 \cdots&\cdots&&\cdots\\
 0&0&\cdots&A
\end{vmatrix}=A^n,
\]
und daraus durch Division mit $A$
\begin{equation}
 \Delta=A^{n-1}.\tag{3}
\end{equation}
Es ist also $\Delta$ die $(n-1)$te Potenz von $A$. Bei dieser Ableitung ist allerdings zunächst vorausgesetzt, dass $A$ von Null verschieden sei. Da aber (3) in Bezug auf die Elemente $a_i^{(h)}$ eine Identität ist, d. h. auch dann gilt, wenn diese Grössen unabhängige Variable sind, so folgt, dass auch noch in diesem Ausnahmefall die Formel (3) gilt, d. h. dass, wenn $A$ verschwindet, auch $\Delta$ verschwindet.

Dies Ergebniss ist ein specieller Fall eines allgemeineren Satzes, nach dem jede beliebige Determinante der Matrix (2) gebildet werden kann. Betrachten wir die $\nu$-reihige Unterdeterminante
\[
\Delta_\nu=\begin{vmatrix}
 A_1^{(1)}&A_2^{(1)}&\cdots&A_\nu^{(1)}\\
 A_1^{(2)}&A_2^{(2)}&\cdots&A_\nu^{(2)}\\
 \cdots&\cdots&&\cdots\\
 A_1^{(\nu)}&A_2^{(\nu)}&\cdots&A_\nu^{(\nu)}
\end{vmatrix},\tag{4}
\]
aus der man durch Permutation der oberen und unteren Indices alle anderen $\nu$-reihigen Unterdeterminanten ableiten kann, so kann man die Multiplicationsregel anwenden, indem man $\Delta_\nu$ nach der Schlussbemerkung des letzten Paragraphen in eine $n$-reihige Determinante verwandelt.

Setzt man
\[
\Delta_\nu=\begin{vmatrix}
 A_1^{(1)}&\cdots&A_\nu^{(1)}&A_{\nu+1}^{(1)}&\cdots&A_n^{(1)}\\
 \cdots&&\cdots&\cdots&&\cdots\\
 A_1^{(\nu)}&\cdots&A_\nu^{(\nu)}&A_{\nu+1}^{(\nu)}&\cdots&A_n^{(\nu)}\\
 0&\cdots&0&1&\cdots&0\\
 \cdots&&\cdots&\cdots&&\cdots\\
 0&\cdots&0&0&\cdots&1
\end{vmatrix},\tag{5}
\]
wobei $n-\nu$ Zeilen und Colonnen beigefügt sind, von denen die ersteren ausser in den Diagonalgliedern lauter Nullen haben. Bildet man jetzt das Product $A\Delta_\nu$, so folgt
\[
 A\Delta_\nu=A^\nu
 \begin{vmatrix}
 a_{\nu+1}^{(\nu+1)}&\cdots&a_n^{(\nu+1)}\\
 \cdots&&\cdots\\
 a_{\nu+1}^{(n)}&\cdots&a_n^{(n)}
 \end{vmatrix},\tag{6}
\]
und dies ist nach dem Satz IX, § 22. Dividirt man hier durch $A$ und wendet die Bezeichnung des § 23 an, so folgt
\begin{equation}
 \Delta_\nu=A^{\nu-1}A_{1,2,\ldots,\nu}^{1,2,\ldots,\nu}.\tag{7}
\end{equation}
Für $\nu=2$ ergiebt sich das specielle Resultat
\begin{equation}
 A_1^{(1)}A_2^{(2)}-A_1^{(2)}A_2^{(1)}=A A_{1,2}^{1,2}.\tag{8}
\end{equation}
Die Formel (8) werden wir später öfter benutzen. In der Bezeichnung durch die Differentialquotienten lässt sie sich so darstellen
\begin{equation}
 A\frac{\partial^2A}{\partial a_1^{(1)}\partial a_2^{(2)}}
 =\frac{\partial A}{\partial a_1^{(1)}}\frac{\partial A}{\partial a_2^{(2)}}
 -\frac{\partial A}{\partial a_1^{(2)}}\frac{\partial A}{\partial a_2^{(1)}}.\tag{9}
\end{equation}
Besonders wichtig ist sie in dem Fall, wo $A$ eine symmetrische Determinante ist, wo also $a_i^{(k)}=a_k^{(i)}$ ist; dann ist auch
\[
 \frac{\partial A}{\partial a_1^{(2)}}=\frac{\partial A}{\partial a_2^{(1)}},
\]
worin bei der Differentiation nicht Rücksicht genommen ist auf die Abhängigkeit $a_i^{(k)}=a_k^{(i)}$, dann wird die Formel (9)
\begin{equation}
 A\frac{\partial^2A}{\partial a_1^{(1)}\partial a_2^{(2)}}
 =\frac{\partial A}{\partial a_1^{(1)}}\frac{\partial A}{\partial a_2^{(2)}}
 -\left(\frac{\partial A}{\partial a_1^{(2)}}\right)^2.\tag{10}
\end{equation}

\section*{§ 29. Interpolation.}

Wir haben schon in § 9 unter einer besonderen Voraussetzung die Aufgabe gelöst, eine ganze rationale Function zu bestimmen, die für eine genügende Anzahl vorgeschriebener Werthe des Arguments gegebene Werthe annimmt. Wir wollen nun durch Anwendung der Determinanten diese Aufgabe allgemein lösen. Wir wollen aber folgende Bemerkungen vorausschicken.

Im § 22, (12) ist der Werth der Determinante
\[
\begin{vmatrix}
 \alpha_1^{n-1}&\alpha_1^{n-2}&\cdots&\alpha_1&1\\
 \alpha_2^{n-1}&\alpha_2^{n-2}&\cdots&\alpha_2&1\\
 \cdots&\cdots&&\cdots&\cdots\\
 \alpha_n^{n-1}&\alpha_n^{n-2}&\cdots&\alpha_n&1
\end{vmatrix}
\]
gleich dem Differenzenproduct
\begin{equation}
(\alpha_1-\alpha_2)(\alpha_1-\alpha_3)\cdots(\alpha_1-\alpha_n)
(\alpha_2-\alpha_3)\cdots(\alpha_2-\alpha_n)\cdots(\alpha_{n-1}-\alpha_n)\tag{1}
\end{equation}
gefunden, und wir wollen jetzt zur Abkürzung dieses Differenzenproduct mit
\begin{equation}
 [\alpha_1,\alpha_2,\alpha_3,\ldots,\alpha_n]\tag{2}
\end{equation}
bezeichnen, so dass die Grösse $[\alpha_1,\alpha_2,\alpha_3,\ldots,\alpha_n]$ ihr Vorzeichen ändert, wenn zwei der Elemente $\alpha_1,\alpha_2,\ldots,\alpha_n$ vertauscht werden.

Wir definiren ferner eine ganze rationale Function $n$ten Grades $f(x)$ der Veränderlichen $x$ durch die Gleichung
\begin{equation}
 f(x)=(x-\alpha_1)(x-\alpha_2)\cdots(x-\alpha_n),\tag{3}
\end{equation}
so dass nach § 13
\begin{align}
 f'(\alpha_1)&=(\alpha_1-\alpha_2)(\alpha_1-\alpha_3)\cdots(\alpha_1-\alpha_n),\notag\\
 f'(\alpha_2)&=(\alpha_2-\alpha_1)(\alpha_2-\alpha_3)\cdots(\alpha_2-\alpha_n),\tag{4}\\
 &\vdots\notag\\
 f'(\alpha_n)&=(\alpha_n-\alpha_1)(\alpha_n-\alpha_2)\cdots(\alpha_n-\alpha_{n-1}).\notag
\end{align}
In dem Product aller dieser Ausdrücke kommt jeder Factor des Productes (1) zweimal mit entgegengesetztem Zeichen vor, so dass wir, da die Anzahl der Factoren von (1) gleich $\frac{n(n-1)}{2}$ ist, die Gleichung erhalten
\begin{equation}
 f'(\alpha_1)f'(\alpha_2)\cdots f'(\alpha_n)=(-1)^{\frac{n(n-1)}2}[\alpha_1,\alpha_2,\ldots,\alpha_n]^2.\tag{5}
\end{equation}
Wir können unser Product (1) aber auch mit Hülfe der Relationen (4) folgendermaassen darstellen:
\begin{align}
 [\alpha_1,\alpha_2,\ldots,\alpha_n]
 &= f'(\alpha_1)[\alpha_2,\alpha_3,\ldots,\alpha_n]\notag\\
 &=-f'(\alpha_2)[\alpha_1,\alpha_3,\ldots,\alpha_n]\tag{6}\\
 &=\cdots=(-1)^{n-1}f'(\alpha_n)[\alpha_1,\alpha_2,\ldots,\alpha_{n-1}],\notag
\end{align}
wo auf der rechten Seite jeder der Klammerausdrücke ein Element weniger enthält als der Klammerausdruck auf der linken Seite.

Es soll also jetzt die Aufgabe gestellt sein, eine ganze rationale Function $(n-1)$ten Grades $\varphi(x)$ zu bestimmen, die für die Argumentwerthe $x=\alpha_1,\alpha_2,\ldots,\alpha_n$ der Reihe nach die vorgeschriebenen Werthe $\varphi(\alpha_1),\varphi(\alpha_2),\ldots,\varphi(\alpha_n)$ annimmt.

Setzt man
\begin{equation}
 \varphi(x)=a_1x^{n-1}+a_2x^{n-2}+\cdots+a_{n-1}x+a_n,\tag{7}
\end{equation}
so sind die $n$ unbekannten Coefficienten $a_1,a_2,\ldots,a_n$ aus den linearen Gleichungen
\begin{align}
 \varphi(\alpha_1)&=a_1\alpha_1^{n-1}+a_2\alpha_1^{n-2}+\cdots+a_{n-1}\alpha_1+a_n,\notag\\
 \varphi(\alpha_2)&=a_1\alpha_2^{n-1}+a_2\alpha_2^{n-2}+\cdots+a_{n-1}\alpha_2+a_n,\tag{8}\\
 &\vdots\notag\\
 \varphi(\alpha_n)&=a_1\alpha_n^{n-1}+a_2\alpha_n^{n-2}+\cdots+a_{n-1}\alpha_n+a_n\notag
\end{align}
zu bestimmen. Die Determinante dieses Systems ist aber genau die Grösse (2), und die Aufgabe ist also nach § 26 immer lösbar, wenn diese Grösse von Null verschieden ist, d. h. wenn von den $n$ Grössen $\alpha_1,\alpha_2,\ldots,\alpha_n$ keine zwei einander gleich sind.

Statt aber die Gleichungen (8) nach § 26 aufzulösen und die gefundenen Ausdrücke in (7) einzusetzen, ist es vorzuziehen, nach § 24, II. die homogenen Grössen $1,a_1,a_2,\ldots,a_n$ aus den $n+1$ Gleichungen (7) und (8) zu eliminiren, wodurch man die Determinantengleichung erhält
\[
\begin{vmatrix}
 \varphi(x)&x^{n-1}&x^{n-2}&\cdots&x&1\\
 \varphi(\alpha_1)&\alpha_1^{n-1}&\alpha_1^{n-2}&\cdots&\alpha_1&1\\
 \cdots&\cdots&\cdots&&\cdots&\cdots\\
 \varphi(\alpha_n)&\alpha_n^{n-1}&\alpha_n^{n-2}&\cdots&\alpha_n&1
\end{vmatrix}=0.\tag{9}
\]
Diese Determinante entwickeln wir nun nach den Elementen der ersten Colonne. Für die dabei auftretenden Unterdeterminanten können wir dann die Bezeichnung (2) anwenden und erhalten so
\begin{align}
&\varphi(x)[\alpha_1,\alpha_2,\ldots,\alpha_n]
-\varphi(\alpha_1)[x,\alpha_2,\alpha_3,\ldots,\alpha_n]
+\varphi(\alpha_2)[x,\alpha_1,\alpha_3,\ldots,\alpha_n]
\notag\\
&\qquad+\cdots+(-1)^n\varphi(\alpha_n)[x,\alpha_1,\alpha_2,\ldots,\alpha_{n-1}]=0.\tag{10}
\end{align}
Nun ist nach (1)
\[
 [x,\alpha_2,\alpha_3,
 \ldots,\alpha_n]=(x-\alpha_2)(x-\alpha_3)\cdots(x-\alpha_n)[\alpha_2,\alpha_3,
 \ldots,\alpha_n],
\]
und wenn man also (10) durch $[\alpha_1,\alpha_2,\ldots,\alpha_n]$ dividirt, und die verschiedenen Ausdrücke (6) dabei berücksichtigt, so erhält man schliesslich
\begin{equation}
 \varphi(x)=f(x)\left\{\frac{\varphi(\alpha_1)}{(x-\alpha_1)f'(\alpha_1)}+\frac{\varphi(\alpha_2)}{(x-\alpha_2)f'(\alpha_2)}+\cdots+\frac{\varphi(\alpha_n)}{(x-\alpha_n)f'(\alpha_n)}\right\}.\tag{11}
\end{equation}
wodurch die Function $\varphi(x)$ völlig bestimmt ist. Dieser Ausdruck für $\varphi(x)$ heisst die Interpolationsformel von Lagrange. Dass sie die gestellten Forderungen erfüllt, lässt sich nachträglich sehr leicht verificiren, wenn man $x=\alpha_1,\alpha_2,
\ldots,\alpha_n$ setzt. Die Nenner $x-\alpha_1,\ldots,x-\alpha_n$ sind nur scheinbar darin enthalten.




\section*{Dritter Abschnitt. Die Wurzeln algebraischer Gleichungen.}

\section*{§ 30. Begriff der Wurzeln. Mehrfache Wurzeln.}

Nachdem in den beiden ersten Abschnitten die algebraischen Grössen mehr von der formalen Seite betrachtet waren, wobei es sich um identische Umformungen von Buchstabenausdrücken handelte, in denen die Buchstaben durchweg als Symbole für variable Grössen aufgefasst werden konnten, treten nun die Zahlengrössen mehr in den Vordergrund.

Wir verstehen hier unter Zahlen, gemäss dem in der Einleitung Festgesetzten, reelle oder imaginäre Grössen von der Form $a+bi$, und stellen die reellen Grössen zur Veranschaulichung durch die Punkte einer geraden Linie, die imaginären durch die Punkte einer Ebene dar. Unter dem absoluten Werth einer imaginären Grösse $a+bi$ verstehen wir $\sqrt{a^2+b^2}$ und bezeichnen ihn nach Weierstrass mit
\[
 |a+bi|.
\]
Es sei nun
\begin{equation}
 f(x)=a_0x^n+a_1x^{n-1}+a_2x^{n-2}+\cdots+a_n
\tag{1}
\end{equation}
eine ganze Function von $x$, worin die Coefficienten irgend welche reelle oder imaginäre Zahlen sind, und der erste, $a_0$, von Null verschieden vorausgesetzt wird. Wenn $\alpha$ eine Zahl ist, die für $x$ gesetzt die Function $f(x)$ zu Null macht, die also der Bedingung $f(\alpha)=0$ genügt, so heisst $\alpha$ eine Wurzel der Gleichung
\[
 f(x)=0.
\]
Wir sagen auch kurz, $\alpha$ ist eine Wurzel von $f(x)$. Nach § 4 lässt sich dann $f(x)$ durch $x-\alpha$ ohne Rest theilen, so dass man
\begin{equation}
 f(x)=(x-\alpha)f_1(x)
\tag{2}
\end{equation}
setzen kann, worin $f_1(x)$ nur vom $(n-1)$ten Grad ist, und denselben ersten Coefficienten $a_0$ hat, wie $f(x)$, also
\[
 f_1(x)=a_0x^{n-1}+\cdots .
\]
Die Coefficienten von $f_1(x)$ sind in § 4, (6) angegeben.

Jede Wurzel von $f_1(x)$ ist also zugleich Wurzel von $f(x)$; und umgekehrt ist jede Wurzel von $f(x)$ entweder $=\alpha$ oder eine Wurzel von $f_1(x)$. Wenn eine Wurzel $\alpha$ von $f(x)$ bekannt ist, so ist die Aufgabe, die übrigen zu finden, auf die Lösung einer Gleichung $(n-1)$ten Grades zurückgeführt.

Das Ziel der Betrachtungen dieses Abschnittes besteht in dem Nachweis, dass jede Function $f(x)$ vom $n$ten Grad wenigstens eine Wurzel hat. Dieser Satz heisst der Fundamentalsatz der Algebra. Zunächst ziehen wir aus (2) die wichtige Folgerung:

\medskip
\noindent I. Eine Gleichung $n$ten Grades kann nicht mehr als $n$ Wurzeln haben.
\medskip

Denn hätte $f(x)$ mehr als $n$ Wurzeln, so hätte $f_1(x)$, was nur vom $(n-1)$ten Grade ist, mehr als $n-1$ Wurzeln. Eine Gleichung ersten Grades hat aber nicht mehr als eine Wurzel, woraus die Richtigkeit unseres Satzes durch vollständige Induction folgt.

Man giebt ihm bisweilen auch den Ausdruck:

\medskip
\noindent II. Wenn eine Function $n$ten Grades $f(x)$ mehr als $n$ Wurzeln hat, so müssen alle ihre Coefficienten Null sein.
\medskip

Ist $\beta$ eine Wurzel von $f_1(x)$, so lässt sich ebenso setzen
\[
 f_1(x)=(x-\beta)f_2(x),
\]
worin $f_2(x)=a_0x^{n-2}+\cdots$ nur vom $(n-2)$ten Grade ist. Nimmt man also an, dass jede der Functionen, die man durch diese Division erhält, $f_1(x), f_2(x),\ldots$, wenigstens eine Wurzel habe, so erhält man schliesslich
\begin{equation}
 f(x)=a_0(x-\alpha)(x-\beta)\cdots(x-\nu),
\tag{3}
\end{equation}
und wir können also den Satz, den wir vorhin als das Ziel unserer Betrachtungen bezeichnet haben, auch so aussprechen: Es soll bewiesen werden:

\medskip
\noindent Eine Function $n$ten Grades lässt sich in $n$ lineare Factoren zerlegen.
\medskip

Die Grössen $\alpha,\beta,\ldots,\nu$ sind dann alle Wurzeln von $f(x)$ und ihre Zahl ist also $n$. Es ist aber nicht ausgeschlossen, dass unter den $\alpha,\beta,\ldots,\nu$ dieselbe Zahl mehrmals vorkommt. Dann würde $f(x)$ weniger als $n$ Wurzeln haben, während doch der Satz bestehen bleibt, dass $f(x)$ in $n$ lineare Factoren zerlegbar ist. Um die Übereinstimmung herzustellen, ist man übereingekommen, wenn $x-\alpha$ mehrmals in $f(x)$ aufgeht, $\alpha$ unter den Wurzeln mehrfach zu zählen und also von einfachen, zweifachen, dreifachen etc. Wurzeln zu sprechen.

Nach dem Begriff der derivirten Functionen [§ 13, (2)] ist, wenn $\alpha$ eine beliebige Grösse ist,
\begin{equation}
 f(x)=f(\alpha)+(x-\alpha)f'(\alpha)+\frac{(x-\alpha)^2}{1\cdot2}f''(\alpha)
 +\frac{(x-\alpha)^3}{1\cdot2\cdot3}f'''(\alpha)+\cdots .
\tag{4}
\end{equation}
Wird also nun wieder angenommen, dass $f(\alpha)$ verschwindet, so ergiebt sich
\[
 f_1(x)=f'(\alpha)+\frac{x-\alpha}{1\cdot2}f''(\alpha)+\frac{(x-\alpha)^2}{1\cdot2\cdot3}f'''(\alpha)+\cdots,
\]
also
\[
 f_1(\alpha)=f'(\alpha).
\]
Es ist also $\alpha$ eine Doppelwurzel von $f(x)$, wenn mit $f(\alpha)$ gleichzeitig $f'(\alpha)$ verschwindet. Auch die Formel (4) zeigt, dass nur unter dieser Voraussetzung $f(x)$ durch $(x-\alpha)^2$ theilbar ist. Diese Schlussweise lässt sich weiter ausdehnen und führt zu dem Satze:

\medskip
\noindent III. Die nothwendige und hinreichende Bedingung dafür, dass $\alpha$ eine $m$-fache Wurzel von $f(x)$ ist, ist die, dass $f(\alpha),f'(\alpha),f''(\alpha),\ldots,f^{(m-1)}(\alpha)$ verschwinden, $f^{(m)}(\alpha)$ nicht verschwindet.

\section*{§ 31. Stetigkeit ganzer Functionen.}

Wenn, wie bisher,
\[
 f(x)=a_0x^n+a_1x^{n-1}+\cdots+a_n
\]
eine ganze Function von $x$ ist, so beweisen wir zunächst folgenden Satz:

\medskip
\noindent IV. $f(x)$ wird zugleich mit $x$ unendlich gross.
\medskip

Das will sagen, wenn $C$ eine beliebig gegebene positive Zahl ist, so kann man die positive Zahl $R$ so wählen, dass
\[
 |f(x)|>C
\]
wird, sobald der absolute Werth von $x$ grösser als $R$ wird; oder geometrisch ausgedrückt: man kann in der $x$-Ebene einen Kreis um den Nullpunkt so beschreiben, dass ausserhalb dieses Kreises der absolute Werth von $f(x)$ nicht mehr unter $C$ herabsinkt, wie gross auch $C$ angenommen ist. Der Satz ist einleuchtend für den Fall einer einfachen Potenz $x^n$; denn ist $r$ der absolute Werth von $x$, so ist $r^n$ der absolute Werth von $x^n$, und $r^n$ ist, sobald $r$ grösser als $1$ geworden ist, grösser als $r$ und wächst also mit $r$ ins Unendliche.

Um den Satz allgemein zu beweisen, bezeichnen wir die absoluten Werthe von $x,a_0,a_1,a_2,\ldots,a_n$ mit $r,c_0,c_1,c_2,\ldots,c_n$. Dann ist
\begin{equation}
 |f(x)|=r^n\left|a_0+\frac{a_1}{x}+\frac{a_2}{x^2}+\cdots+\frac{a_n}{x^n}\right|,
\tag{2}
\end{equation}
und, weil nach den in der Einleitung bewiesenen Sätzen der absolute Werth einer Summe nicht kleiner als die Differenz und nicht grösser als die Summe der absoluten Werthe der Summanden sein kann,
\[
\left|a_0+\frac{a_1}{x}+\frac{a_2}{x^2}+\cdots+\frac{a_n}{x^n}\right|
\ge c_0-\left|\frac{a_1}{x}+\frac{a_2}{x^2}+\cdots+\frac{a_n}{x^n}\right|
\ge c_0-\frac{c_1}{r}-\frac{c_2}{r^2}-\cdots-\frac{c_n}{r^n}.
\]
Da die $c_0,c_1,\ldots,c_n$ fest gegebene Constanten sind, so kann man $R$ so gross wählen, dass, sobald $r>R$ ist,
\[
 \frac{c_1}{r}+\frac{c_2}{r^2}+\cdots+\frac{c_n}{r^n}
\]
beliebig klein, z. B. kleiner als $\tfrac12c_0$ wird; dann ist
\[
 \left|a_0+\frac{a_1}{x}+\frac{a_2}{x^2}+\cdots+\frac{a_n}{x^n}\right|>\tfrac12c_0,
\]
und also nach (2)
\[
 |f(x)|>\tfrac12c_0R^n,
\]
also, wenn man
\begin{equation}
 R^n>\frac{2C}{c_0}
\tag{3}
\end{equation}
annimmt,
\[
 |f(x)|>C.
\]
Hieran schliesst sich der Satz:

\medskip
\noindent V. $f(x)$ ist eine stetige Function von $x$.
\medskip

Damit soll folgendes gesagt sein: Ist der absolute Werth von $x$ kleiner als eine beliebige endliche Grösse $R$, so kann man, wie klein auch die positive Grösse $\omega$ angenommen wird, eine positive Grösse $\varepsilon$ derart finden, dass
\begin{equation}
 |f(x+h)-f(x)|<\omega
\tag{4}
\end{equation}
wird, so lange der absolute Werth $\rho$ von $h$ kleiner als $\varepsilon$ bleibt, und zwar so, dass $\varepsilon$ nur von der Wahl von $\omega$, nicht aber von $x$ abhängig ist.

Nach § 13, (2) ist
\begin{equation}
 f(x+h)-f(x)=h f'(x)+\frac{h^2}{1\cdot2}f''(x)+\cdots+h^n a_0.
\tag{5}
\end{equation}
Nehmen wir nun den absoluten Werth von $h$ kleiner als $1$ an, so ist der absolute Werth der in der Klammer stehenden Grösse
\begin{equation}
 f'(x)+\frac{h}{1\cdot2}f''(x)+\cdots+h^{n-1}a_0
\tag{6}
\end{equation}
kleiner als eine bestimmte von Null verschiedene Zahl $K$, die man erhält, wenn man in (6) $h$ durch $1$, die Coefficienten $a_0,a_1,\ldots,a_{n-1}$ durch ihre absoluten Werthe und $x$ durch $R$ ersetzt, weil dadurch jedes einzelne Glied der Summe (6) durch seinen absoluten Werth oder durch eine grössere Zahl ersetzt wird. Wählt man $\varepsilon$ so klein, dass
\[
 \varepsilon K<\omega,
\]
so ist wegen (5) die Forderung (4) erfüllt, so lange $\rho<\varepsilon$ bleibt.

Wir können dem Satze von der Stetigkeit der Function $f(x)$ noch einen etwas anderen Ausdruck geben, der uns für die Folge wichtig ist. Nach dem Satze, dass der absolute Werth einer Summe von zwei Gliedern der Grösse nach zwischen der Summe und der Differenz der absoluten Werthe der Glieder liegt (vgl. Einleitung, S. 19), erhalten wir aus (5), wenn wir
\[
 \varphi(h)=h f'(x)+\frac{h^2}{1\cdot2}f''(x)+\cdots+h^n a_0
\]
setzen,
\[
 |f(x)|-|\varphi(h)|<|f(x+h)|<|f(x)|+|\varphi(h)|,
\]
oder
\[
 -|\varphi(h)|<|f(x+h)|-|f(x)|<|\varphi(h)|.
\]
Nun ist, wenn $\rho<\varepsilon$ ist, $|\varphi(h)|<\omega$, und wir können also auch sagen:

\medskip
\noindent VI. Der absolute Werth der Differenz
\[
 |f(x+h)|-|f(x)|
\]
ist kleiner als eine beliebig gegebene Grösse $\omega$, sobald der absolute Werth von $h$ kleiner als ein genügend kleines $\varepsilon$ ist.
\medskip

Daraus können wir noch schliessen, wenn wir $x=y+iz$ setzen, und $h$ entweder reell oder rein imaginär annehmen, dass $|f(y+zi)|$ bei feststehendem $z$ eine stetige Function von $y$ und bei feststehendem $y$ eine stetige Function von $z$ ist.

Sind die Coefficienten $a_0,a_1,\ldots,a_n$ und die Variable $x$ reell, so kann man $x$ dem absoluten Werthe nach so gross wählen, dass das Vorzeichen von $f(x)$ mit dem Vorzeichen des ersten Gliedes $a_0x^n$ übereinstimmt, also bei positivem $a_0$ und geradem $n$ positiv, bei ungeradem $n$ für negative $x$ negativ und für positive $x$ positiv wird.

Denn man kann in
\[
 f(x)=x^n\left(a_0+\frac{a_1}{x}+\frac{a_2}{x^2}+\cdots+\frac{a_n}{x^n}\right)
\]
$x$ dem absoluten Werthe nach so gross wählen, dass der absolute Werth von
\[
 \frac{a_1}{x}+\frac{a_2}{x^2}+\cdots+\frac{a_n}{x^n}
\]
unter den von $a_0$ heruntersinkt und also der Klammerausdruck
\[
 a_0+\frac{a_1}{x}+\frac{a_2}{x^2}+\cdots+\frac{a_n}{x^n}
\]
das Vorzeichen von $a_0$ hat, was es dann für jedes absolut grössere $x$ beibehält.

\section*{§ 32. Vorzeichenwechsel von $f(x)$. Wurzeln von Gleichungen ungeraden Grades und von reinen Gleichungen.}

Auf Grund der Sätze des vorigen Paragraphen können wir das folgende wichtige Theorem beweisen.

\medskip
\noindent VIII. Sind die Coefficienten von $f(x)$ reell und existiren zwei reelle Werthe $a,b$ von $x$, so dass $f(a)$ und $f(b)$ entgegengesetzte Vorzeichen haben, so giebt es zwischen $a$ und $b$ wenigstens eine Wurzel $\xi$ der Gleichung $f(x)=0$.
\medskip

Wir wollen annehmen, es sei $a<b$ und $f(a)$ negativ, $f(b)$ positiv. Wir theilen die Zahlen zwischen $a$ und $b$ so in zwei Theile $\mathfrak A$ und $\mathfrak B$, dass eine Zahl $\alpha$ zu $\mathfrak A$ gehört, wenn zwischen $\alpha$ und $a$, die Grenzen eingerechnet, kein $x$ liegt, wofür $f(x)$ positiv wird, und eine Zahl $\beta$ zu $\mathfrak B$, wenn zwischen $\beta$ und $a$ wenigstens ein Werth von $x$ liegt, für den $f(x)$ positiv wird. Ein drittes ist offenbar nicht möglich und jede Zahl zwischen $a$ und $b$ ist also in $\mathfrak A$ oder in $\mathfrak B$ untergebracht. Wir rechnen noch die Zahlen, die kleiner als $a$ sind, zu $\mathfrak A$, und die grösser als $b$ sind, zu $\mathfrak B$.

Gehört $\alpha$ zu $\mathfrak A$, so gehört auch jedes $x$, was kleiner ist als $\alpha$, zu $\mathfrak A$, und gehört $\beta$ zu $\mathfrak B$, so gehört jedes grössere $x$ zu $\mathfrak B$. Es ist also jedes $\beta$ grösser als jedes $\alpha$ und die beiden Zahlengebiete $\mathfrak A$, $\mathfrak B$ bilden einen Schnitt (Einleitung, S. 5) der durch eine Zahl $\xi$ erzeugt wird, so dass jede Zahl $\alpha$, die kleiner als $\xi$ ist, zu $\mathfrak A$ gehört, jede Zahl $\beta$, die grösser als $\xi$ ist, zu $\mathfrak B$.

Ist nun für irgend eine Zahl $x_0$ der Werth von $f(x_0)$ negativ oder positiv, so lässt sich nach V. des vorigen Paragraphen eine positive Grösse $\varepsilon$ so bestimmen, dass für jedes $x$ zwischen $x_0-\varepsilon$ und $x_0+\varepsilon$ die Werthe von $f(x)$ immer negativ oder immer positiv sind. Daraus ergiebt sich, dass $f(\xi)$ weder negativ noch positiv sein kann und also Null sein muss. Denn wäre $f(\xi)$ negativ, so wäre auch $f(x)$ zwischen $\xi$ und $\xi+\varepsilon$ negativ; diese Werthe würden also zu $\mathfrak A$ gehören, während sie doch grösser als $\xi$ sind; und wäre $f(\xi)$ positiv, so wäre $f(x)$ positiv, so lange $x$ zwischen $\xi-\varepsilon$ und $\xi$ liegt; diese Werthe würden also, obwohl sie kleiner als $\xi$ sind, zu $\mathfrak B$ gehören, und beides ist nach der Definition von $\xi$ unmöglich. Ganz ähnlich kann geschlossen werden, wenn $f(a)$ positiv und $f(b)$ negativ ist, und unser Satz ist also erwiesen.

Hieraus ziehen wir zwei wichtige Folgerungen. Wenn $f(\xi)$ verschwindet, so lässt sich ein Intervall
\[
 \xi-\varepsilon<x<\xi+\varepsilon
\]
so bestimmen, dass $f(x)$ ausser für $x=\xi$ in diesem Intervall nicht verschwindet. Es sind dann vier Fälle in Bezug auf die Vorzeichen von $f(x)$ in den beiden Intervallen $\xi-\varepsilon<x<\xi$ und $\xi<x<\xi+\varepsilon$ möglich, die im folgenden Schema zusammengestellt sind:
\[
\begin{array}{c|cc}
&\xi-\varepsilon<x<\xi&\xi<x<\xi+\varepsilon\\
1.&+&-\\
2.&-&+\\
3.&+&+\\
4.&-&-
\end{array}
\]
Ist $f^{(\nu)}(x)$ die erste unter den Derivirten von $f(x)$, die für $x=\xi$ von Null verschieden ist, so zeigt die Formel [§ 13, (2)]
\[
 f(\xi+h)=\frac{h^\nu}{\Pi(\nu)}f^{(\nu)}(\xi)+\cdots,
\]
dass diese vier Fälle durch folgende Kennzeichen unterschieden sind:
\[
\begin{array}{lll}
1.&\nu\text{ ungerade},& f^{(\nu)}(\xi)<0,\\
2.&\nu\text{ ungerade},& f^{(\nu)}(\xi)>0,\\
3.&\nu\text{ gerade},& f^{(\nu)}(\xi)>0,\\
4.&\nu\text{ gerade},& f^{(\nu)}(\xi)<0.
\end{array}
\]
Wir sagen, dass im ersten Falle $f(x)$ abnehmend, im zweiten wachsend durch Null geht. Im Falle 3 ist $f(\xi)=0$ ein Minimum, im Falle 4 ein Maximum.

Wenn insbesondere $f'(\xi)$ von Null verschieden, $\xi$ also eine einfache Wurzel ist, so ist das positive oder negative Vorzeichen von $f'(\xi)$ das Kennzeichen für das Wachsen oder Abnehmen der Function $f(x)$ beim Durchgang von $x$ durch $\xi$.

Wir beweisen mit diesen Hülfsmitteln noch zwei weitere wichtige Sätze.

\medskip
\noindent IX. Die Gleichung
\[
 x^n-a=0
\]
hat, wenn $a$ eine positive Zahl ist, eine und nur eine positive Wurzel.
\medskip

Denn setzen wir
\begin{equation}
 f(x)=x^n-a,
\tag{1}
\end{equation}
so ist $f(x)$ nach § 31, VII. für ein hinlänglich grosses positives $x$ positiv, und für $x=0$ negativ; also giebt es einen positiven Werth von $x$, den man mit $\sqrt[n]{a}$ bezeichnet, für den $f(x)$ verschwindet. Dass es aber nur einen solchen Werth geben kann, folgt daraus, dass $x^n$, so lange $x$ positiv bleibt, mit $x$ fortwährend wächst.

Ist $n$ ungerade, so existirt auch für negative $a$ eine und nur eine negative Wurzel von (1), was ebenso bewiesen wird.

\medskip
\noindent X. Eine Gleichung ungeraden Grades mit reellen Coefficienten hat immer wenigstens eine reelle Wurzel.
\medskip

Denn nach Satz VII. kann $f(x)$, wenn der Grad ungerade ist, sowohl positiv als negativ werden und $f(x)=0$ hat also nach VIII. eine Wurzel.

Wir beweisen endlich noch:

\medskip
\noindent XI. Eine reine Gleichung hat immer eine Wurzel.
\medskip

Unter einer reinen Gleichung verstehen wir eine Gleichung von der Form
\begin{equation}
 x^n-a=0,
\tag{2}
\end{equation}
wo $a$ eine beliebige complexe Grösse ist. Dass diese Gleichung für ein reelles positives $a$ immer eine Wurzel hat, ist oben schon bewiesen, und für $a=0$ wird sie offenbar durch $x=0$ befriedigt.

Der allgemeine Beweis kann in zwei Theile zerlegt werden; es genügt nämlich, wenn wir zunächst $n=2$, dann $n$ ungerade voraussetzen. Denn existirt die Grösse $\sqrt{a}$ und $\sqrt[n]{a}$ oder $\sqrt[n]{\alpha}$ für jedes $\alpha$, so ist also auch die Existenz von $\sqrt[n]{a}$ für jedes beliebige $n$ sichergestellt. Es würde sogar genügen, die Existenz von $\sqrt[n]{a}$ für den Fall zu beweisen, dass $n$ eine Primzahl ist. Daraus ist aber zunächst kein besonderer Vortheil zu ziehen.

Wir bezeichnen die complexe Grösse $a$ mit $b+ci$ und setzen, da $x$ im Allgemeinen auch complex sein wird,
\[
 x=y+iz.
\]
Dann haben wir zunächst zu zeigen, dass
\begin{equation}
 (y+iz)^2=b+ic,
\tag{3}
\end{equation}
welche reellen Werthe auch $b,c$ haben mögen, durch reelle $y,z$ befriedigt werden kann. Die Gleichung (3) ist aber erfüllt, wenn
\begin{equation}
 y^2-z^2=b,\\qquad 2yz=c,
\tag{4}
\end{equation}
also
\[
 (y^2+z^2)^2=b^2+c^2
\]
und folglich
\begin{equation}
 y^2+z^2=\sqrt{b^2+c^2}.
\tag{5}
\end{equation}
Nach IX. existirt immer eine positive Quadratwurzel $\sqrt{b^2+c^2}$. Aus (4) und (5) folgt nun
\[
 y^2=\frac{b+\sqrt{b^2+c^2}}{2},\qquad
 z^2=\frac{-b+\sqrt{b^2+c^2}}{2},
\]
und die beiden Grössen $\pm b+\sqrt{b^2+c^2}$ sind offenbar positiv, da $b^2+c^2>b^2$, also auch $\sqrt{b^2+c^2}>\sqrt{b^2}$. Wir erhalten daraus
\[
 y=\pm\sqrt{\frac{b+\sqrt{b^2+c^2}}{2}},\qquad
 z=\pm\sqrt{\frac{-b+\sqrt{b^2+c^2}}{2}},
\]
worin aber das eine Vorzeichen durch das andere bestimmt ist durch die Bedingung
\[
 2yz=c;
\]
eines der beiden Vorzeichen aber bleibt willkürlich, so dass wir nicht nur eine, sondern zwei (entgegengesetzte) Lösungen von (3) erhalten.

Es sei ferner $n$ ungerade; wir nehmen die zu lösende Gleichung in der Form an:
\begin{equation}
 (y+iz)^n=b+ci.
\tag{6}
\end{equation}
Ist $c=0$, also $a$ reell, so können wir $z=0$ setzen und erhalten nach IX. einen reellen Werth von $x$. Nehmen wir aber $c$ von Null verschieden an, so folgt aus (6) (Einleitung, S. 19)
\begin{equation}
 (y-iz)^n=b-ci.
\tag{7}
\end{equation}
Multipliciren wir die beiden Gleichungen (6), (7) mit einander, so folgt
\[
 (y^2+z^2)^n=b^2+c^2,
\]
und daraus (nach IX.)
\begin{equation}
 y^2+z^2=\sqrt[n]{b^2+c^2},
\tag{8}
\end{equation}
wodurch der absolute Werth von $x$ bestimmt ist. Es ergiebt sich ferner aus (6) und (7):
\begin{equation}
 \varphi(y,z)=(y-iz)^n(b+ic)-(y+iz)^n(b-ci)=0.
\tag{9}
\end{equation}
Nun ist $\varphi(y,z)$ eine ganze homogene Function $n$ten Grades der beiden Veränderlichen $y,z$, und zwar mit reellen Coefficienten (da die Vertauschung von $i$ mit $-i$ in $\varphi(y,z)$ nichts ändert), und wenn wir nach absteigenden Potenzen von $y$ ordnen, so ist das erste Glied $2ciy^n$, also nach Voraussetzung von Null verschieden. Setzen wir nun
\[
 y=\lambda z,
\]
so erhalten wir aus (9)
\begin{equation}
 \varphi(\lambda,1)=0,
\tag{10}
\end{equation}
also eine Gleichung $n$ten Grades für $\lambda$, die nach X. gewiss eine Wurzel hat. Ist $\lambda$ bestimmt, so folgt aus (8)
\begin{equation}
 y=\lambda\sqrt{\frac{\sqrt[n]{b^2+c^2}}{1+\lambda^2}},\qquad
 z=\sqrt{\frac{\sqrt[n]{b^2+c^2}}{1+\lambda^2}}.
\tag{11}
\end{equation}
Hierdurch sind die Gleichungen (8), (9) thatsächlich befriedigt; aus (9) aber folgt
\[
 \frac{(y+iz)^n}{b+ic}=\frac{(y-iz)^n}{b-ci},
\]
und aus (8), dass der gemeinsame Werth dieser beiden Brüche $+1$ ist. Wir haben daher
\[
 (y+iz)^n=\pm(b+ci),\qquad (y-iz)^n=\pm(b-ci),
\]
und es kann also das Vorzeichen der Quadratwurzel in (11) so bestimmt werden, dass die Gleichungen (6) und (7) erfüllt sind.

\section*{§ 33. Lösung reiner Gleichungen durch trigonometrische Functionen.}

Weit vollständiger und einfacher kann die Auflösbarkeit einer reinen Gleichung dargethan werden, wenn man die trigonometrischen Functionen und ihre einfachsten Eigenschaften als bekannt voraussetzt. Freilich sind diese Functionen der eigentlichen Algebra fremd und darum ist es befriedigender, die principiellen Fragen, so wie es im Vorhergehenden geschehen ist und auch noch später geschehen soll, ohne ihre Hülfe zu beantworten. Für die Anwendung und eine bequemere Anschauung wollen wir aber dieses Hülfsmittel doch nicht entbehren.

Setzen wir, wenn $b,c$ reelle Zahlen sind,
\begin{equation}
 b=r\cos\varphi,\qquad c=r\sin\varphi,
\tag{1}
\end{equation}
so ist, wenn wir $r$ positiv annehmen, der Winkel $\varphi$ hierdurch bis auf ein Vielfaches von $2\pi$ bestimmt. Er wird völlig bestimmt sein, wenn wir ein Intervall von der Grösse $2\pi$ festsetzen, in dem er liegen soll, z. B.
\[
 -\pi<\varphi\le \pi.
\]
In geometrischer Auffassung sind $r,\varphi$ die Polarcoordinaten des Punktes, dessen rechtwinklige Coordinaten $b,c$ sind. $r$ ist der absolute Werth der complexen Grösse $a=b+ci$ und $\varphi$ wollen wir die Phase dieser complexen Grösse nennen.

Das Product zweier complexen Grössen
\[
 a=r(\cos\varphi+i\sin\varphi),\qquad
 a'=r'(\cos\varphi'+i\sin\varphi')
\]
ergiebt sich in der Form
\begin{equation}
 aa'=rr'[\cos(\varphi+\varphi')+i\sin(\varphi+\varphi')],
\tag{3}
\end{equation}
und daraus erhält man, wenn $n$ eine beliebige ganze positive (oder selbst negative) Zahl ist, den nach Moivre benannten Lehrsatz
\begin{equation}
 (\cos\varphi+i\sin\varphi)^n=\cos n\varphi+i\sin n\varphi,
\tag{4}
\end{equation}
der dazu geführt hat, die Grösse $\cos\varphi+i\sin\varphi$ als Exponentialfunction mit imaginärem Exponenten zu betrachten, und demnach
\[
 \cos\varphi+i\sin\varphi=e^{i\varphi}
\]
zu setzen. Dadurch erhalten die Formeln (3) und (4) die Gestalt
\begin{equation}
 e^{i\varphi}e^{i\varphi'}=e^{i(\varphi+\varphi')},\qquad (e^{i\varphi})^n=e^{in\varphi}.
\tag{5}
\end{equation}
Wenn wir demnach
\begin{equation}
 a=r(\cos\varphi+i\sin\varphi)
\tag{6}
\end{equation}
setzen und unter $\sqrt[n]{r}$ den nach § 32, IX. existirenden und völlig bestimmten positiven Werth verstehen, dessen $n$te Potenz $=r$ ist, so ist
\begin{equation}
 x=\sqrt[n]{r}\left(\cos\frac{\varphi}{n}+i\sin\frac{\varphi}{n}\right)
\tag{7}
\end{equation}
eine Grösse, deren $n$te Potenz $=a$ ist, also eine Wurzel der Gleichung
\begin{equation}
 x^n=a,
\tag{8}
\end{equation}
wenn unter $n$ eine beliebige positive ganze Zahl verstanden wird. Derselben Forderung genügt aber auch jeder der Werthe
\begin{equation}
 x_k=\sqrt[n]{r}\left(\cos\frac{\varphi+2k\pi}{n}+i\sin\frac{\varphi+2k\pi}{n}\right),
\tag{9}
\end{equation}
wenn $k$ eine beliebige ganze Zahl ist. In (9) sind aber nicht mehr als $n$ verschiedene Werthe enthalten, die man erhält, wenn man
\begin{equation}
 k=0,1,2,\ldots,n-1
\tag{10}
\end{equation}
setzt; denn vermehrt man $k$ um ein Vielfaches von $n$, so ändert (9) seinen Werth nicht, während die $n$ Werthe (10) lauter verschiedene Werthe von $x$ ergeben, weil sie verschiedene Phasen haben.

Die Gleichung (8) hat demnach nicht nur eine, sondern $n$ verschiedene Wurzeln. Die geometrischen Bilder der Werthe $x_k$ liegen alle auf einem Kreise mit dem Radius $\sqrt[n]{r}$, und zwar um den Winkel $\frac{2\pi}{n}$ von einander entfernt. Man erhält den ersten dieser Punkte dadurch, dass man den gegebenen Winkel $\varphi$ in $n$ gleiche Theile theilt. Der Radius $\sqrt[n]{r}$ ist kleiner oder grösser als $r$, je nachdem $r$ kleiner oder grösser als $1$ ist. Die beistehende Figur zeigt diese Verhältnisse für $n=3$ unter der Voraussetzung, dass $r>1$ ist.

Man kann die verschiedenen Werthe von $x_k$ dadurch erhalten, dass man den ersten von ihnen
\[
 x_0=\sqrt[n]{r}\left(\cos\frac{\varphi}{n}+i\sin\frac{\varphi}{n}\right)
\]
mit den verschiedenen Werthen von
\begin{equation}
 \xi_k=\cos\frac{2k\pi}{n}+i\sin\frac{2k\pi}{n}
\tag{11}
\end{equation}
multiplicirt. Die Grössen $\xi_k$,
\begin{equation}
 \xi_k^n=1,
\tag{12}
\end{equation}
sind Wurzeln der Gleichung $x^n=1$, und heissen die $n$ten Einheitswurzeln. Es sind ihrer $n$ und nur $n$ verschiedene Werthe, von denen bei ungeradem $n$ nur einer (für $k=0$), bei geradem $n$ zwei (für $k=0$ und $k=\frac n2$) reell sind. Die geometrischen Bilder der $n$ten Einheitswurzeln sind die Eckpunkte eines dem Kreise mit dem Radius $1$ einbeschriebenen regelmässigen $n$-Ecks. Ihre algebraische Bestimmung ist Gegenstand der Kreistheilungslehre. Bezeichnen wir den Werth von $\xi_k$ für $k=1$ mit $\varepsilon$, setzen also
\[
 \varepsilon=\cos\frac{2\pi}{n}+i\sin\frac{2\pi}{n},
\]
so ist nach dem Moivre'schen Satze
\[
 \xi_k=\varepsilon^k,
\]
so dass alle $n$ten Einheitswurzeln als Potenzen von einer unter ihnen dargestellt sind.

\end{document}
